\documentclass[12pt,a4paper]{amsart}
\usepackage{amsfonts,amssymb,amsxtra,amsthm,amsmath,amscd,mathrsfs,cite}
\numberwithin{equation}{section}
\usepackage{graphicx}
\usepackage[hidelinks]{hyperref}
\hypersetup{hypertexnames=false,bookmarksdepth=3}
\makeatletter
\@ifundefined{subjclassname@2020}{%
  \@namedef{subjclassname@2020}{%
    \textup{2020} Mathematics Subject Classification}%
}{}
\makeatother
\textwidth=160mm
\oddsidemargin=0mm
\evensidemargin=0mm
\emergencystretch=1em
\raggedbottom
\renewcommand{\thefootnote}{\fnsymbol{footnote}}
\addtocounter{footnote}{1}
\def\e{{\rm e}}
\def\d{\,{\rm d}}
\def\leq{\leqslant}
\def\geq{\geqslant}
\def\le{\leqslant}
\def\ge{\geqslant}
\theoremstyle{plain}
\newtheorem{theorem}{Theorem}[section]
\newtheorem{proposition}{Proposition}[section]
\newtheorem{lemma}[proposition]{Lemma}
\newtheorem*{grh}{\(\mathrm{GRH}(N)\)}
\newtheorem*{hypothesisR}{Hypothesis R}
\newtheorem*{hypothesisS}{Hypothesis S}
\newtheorem{corollary}[theorem]{Corollary}
\theoremstyle{definition}
\newtheorem*{definition}{Definition}
\theoremstyle{remark}
\newtheorem{remark}{Remark}
\begin{document}
\title[A conditional proof of the Goldbach conjecture]{A conditional proof of the Goldbach conjecture}
\author{Jiamin Li \& Jianya Liu}
\address{Jiamin Li\\
Mathematical Research Center\\
Shandong University\\
Jinan, Shandong 250100\\
P. R. China}
\email{lijiamin@mail.sdu.edu.cn}
\address{Jianya Liu\\
State Key Laboratory of Cryptography and Digital Economy Security \& Mathematical Research Center\\
Shandong University\\
Jinan, Shandong 250100\\
P. R. China}
\email{jyliu@sdu.edu.cn}
\date{}
\keywords{Goldbach conjecture, Beta function, Dirichlet $L$-functions,
	quaternionic spectral function, Riemann hypothesis}
\subjclass[2020]{Primary 11P32; Secondary 11M06, 47B15}
\maketitle

\begin{flushright}
\begin{minipage}{0.82\textwidth}
\small\itshape
``The famous problem of which I propose to speak tonight is
probably as difficult as any of the unsolved problems of
mathematics, and my lecture cannot be entirely easy.''
\par\smallskip
\raggedleft\normalfont
---G.~H.~Hardy, \emph{Goldbach's Theorem}~\cite[p.~1]{Hardy1922}
\end{minipage}
\end{flushright}

\begin{abstract} 
	We give what is, to our knowledge, the first conditional proof of the binary Goldbach conjecture for all sufficiently large even integers, assuming the generalized Riemann hypothesis for Dirichlet \(L\)-functions together with either of two additional spectral hypotheses. Our approach is fundamentally different from the classical sieve and circle methods. For each even integer \(N\), odd Dirichlet characters modulo \(N\) give an exact projection of a smoothed Goldbach correlation. Subtracting the exact diagonal and a model built from Euler’s beta function isolates an arithmetic residual, which under GRH defines a compact selfadjoint Hilbert–Schmidt operator on a Hardy space. An exact quadratic-form identity expresses the weighted Goldbach error in terms of this operator. Its spectrum defines a symmetric quaternionic function whose zero locations determine the operator norm. A prescribed zero-free region yields the weighted Hardy–Littlewood asymptotic \(\mathcal G(N)\sim c_\phi\mathfrak S(N)\). Under the stronger hypothesis that all spectral zeros have real part \(1/2\), we obtain
	\[
	\mathcal G(N)=c_\phi\mathfrak S(N)
	+O_\phi\!\left(
	\sqrt{\frac{\mathfrak S(N)}{N}}(\log N)^{3/2}
	\right).
	\]
	Here \(\mathcal G(N)\) is the normalized weighted Goldbach correlation, \(\mathfrak S(N)\) is the Goldbach singular series, and \(c_\phi=\int_0^1\phi(u)^2\,du\) for a fixed nonzero, nonnegative, symmetric smooth weight supported in \([1/3,2/3]\). Under either spectral hypothesis together with GRH, every sufficiently large even integer \(N\) has \(\gg_\phi\mathfrak S(N)N/(\log N)^2\) ordered representations as a sum of two primes in \([N/3,2N/3]\). The method also extends naturally to other problems in additive number theory.
\end{abstract}

\section{Introduction and statement of results}\label{story:section1}

\subsection{Historical background}\label{story:history}

Goldbach's conjecture, proposed in 1742~\cite{GoldbachEuler1742},
states that every even integer
$N\ge4$ is a sum of two primes. The classical approaches to
this problem are the circle method and sieve methods.
For the related ternary problem, in 1923 Hardy and
Littlewood~\cite{HL23} used the circle method to show, under a
quasi-Riemann Hypothesis, that every sufficiently large odd
integer is a sum of three primes; Vinogradov~\cite{Vin37}
removed this hypothesis in 1937. For the binary problem, R\'enyi's prime $+$ almost-prime
theorem~\cite{Ren48} was followed by Chen's $(1+2)$ theorem~\cite{Chen6673}
and theorem $(1+1.9)$ of the authors~\cite{LiLiu19}.
The latter gives $N=p_1+r p_2$ for sufficiently
large even $N$, with $p_1, p_2$ prime and $r=1$ or prime satisfying
$r\le p_2^{0.9}$. Despite these advances, neither the circle
method nor sieve methods have so far resolved the binary
Goldbach conjecture. An excellent
account of these classical methods and their principal results
is given in the monograph of Chengdong Pan and Chengbiao
Pan~\cite{PanPan}.

Beyond existence, Hardy and Littlewood~\cite{HL23} conjectured
the following weighted asymptotic formula:
\begin{equation}\label{eq:hardy-littlewood-goldbach}
 \frac1N\sum_{\substack{m,n\ge1\\m+n=N}}
       \Lambda(m)\Lambda(n)\sim\mathfrak S(N)
 \qquad(N\to\infty,\ N\ \text{even}),
\end{equation}
where the von Mangoldt function is defined by $\Lambda(n)=\log p$
if $n=p^k$ for a prime $p$ and an
integer $k\ge1$, and $\Lambda(n)=0$ otherwise. The singular
series is
\begin{equation}\label{eq:singular-series}
 \mathfrak S(N)=2C_2\prod_{\substack{p\mid N\\p>2}}
                      \frac{p-1}{p-2},\qquad
 C_2=\prod_{p>2}\left(1-\frac1{(p-1)^2}\right).
\end{equation}
The singular series records the local arithmetic densities of
prime pairs.

No proof of the binary conjecture
is known under the generalized Riemann Hypothesis (GRH) alone,
even when GRH is assumed for all Dirichlet $L$-functions.
Under this hypothesis, Hardy and Littlewood~\cite{HL24V}
obtained the classical exceptional-set bound
$E(x)\ll_\varepsilon x^{1/2+\varepsilon}$ for every
$\varepsilon>0$, where $E(x)$ counts the even integers
$4\le N\le x$ that are not sums of two primes. This estimate
does not exclude exceptions altogether.

\subsection{Our results}\label{story:results}

We give a conditional proof of the binary Goldbach conjecture
for all sufficiently large even integers by supplementing GRH
with either of two spectral hypotheses.
Our route is fundamentally different from the circle method and
sieve methods. We draw directly on the classical ideas of
Dirichlet, Euler, Hilbert and P\'olya, Riemann, and Hardy.
Starting from the Goldbach problem, we use odd characters to
isolate the correlation and the Beta model to define its residual.
Under GRH, the residual yields a selfadjoint operator. Zeros of
its spectral function control the norm; a Hardy quadratic-form
identity transfers this control to the Goldbach error.

More precisely, we establish conditional asymptotic formulas
for a weighted form of the Hardy--Littlewood prediction.
Fix a nonzero function
$\phi\in C_c^\infty(\mathbb R)$ supported in $[1/3,2/3]$,
with $0\le\phi\le1$ and $\phi(u)=\phi(1-u)$. For even
$N\ge4$, write
\begin{align}
 \mathcal G(N)&=\frac1N
   \sum_{\substack{n+m=N\\(nm,N)=1}}
       \Lambda(n)\Lambda(m)\phi(n/N)\phi(m/N),
       \label{eq:goldbach-correlation}\\
 c_\phi&=\int_0^1\phi^2(u)\,\mathrm{d}u>0,
       \label{eq:c0}
\end{align}
The sum runs over ordered pairs of positive integers and gives a
normalized weighted count for the Goldbach equation $n+m=N$.
Our aim is to prove, under our hypotheses, that $\mathcal G(N)>0$
for all sufficiently large even $N$, with a positive contribution
from prime pairs remaining after the removal of proper prime powers.
The expected main term is $c_\phi\mathfrak S(N)$, with the same singular series
\eqref{eq:singular-series} as in \eqref{eq:hardy-littlewood-goldbach}.

Gauss's congruences gave arithmetic a unifying language.
Here the decisive choice is to fix the modulus at the very integer
$N$ to be represented. Within this one modulus, Dirichlet's odd
characters reveal Goldbach, an arithmetic operator $T_N$ becomes
accessible, and its spectrum produces a quaternionic function in
Riemann's spirit. Each $N$ thus carries an arithmetic world of its
own. This is our departure from the usual sieve and circle-method
routes.

For each $N$, the construction uses the following part of GRH.

\begin{grh}
Every nontrivial zero of every primitive Dirichlet $L$-function
whose conductor divides $N$ has real part $1/2$.
\end{grh}

The following five steps connect the arithmetic and spectral sides of the proof.

\emph{Dirichlet.} The elementary sign $\chi(-1)=-1$ is decisive:
odd-character orthogonality separates equal from opposite residue
classes, and the support of the weight makes the latter exactly
$n+m=N$.
Logarithmic derivatives of $L$-functions turn our exact projection
into a two-variable arithmetic kernel.

\emph{Euler.} The Beta model encodes the off-diagonal geometry
and recovers the Goldbach main term; subtracting the exact diagonal
and this model isolates the residual. The theory of several complex
variables reveals striking phenomena; natural arithmetic examples
help make them clear.
As Lebl~\cite[\S1.1]{LeblSCV} writes,
``There is something very special about the holomorphic
category.'' Here the Beta model supplies an example dictated by
arithmetic. After continuation and conjugation in the second
variable, the model and its residual are Fueter regular. In coordinates
$\omega=s+zj$, the residual satisfies
\[
 (\partial_{\bar s}+j\partial_z)\mathscr E_N(s,\bar z)=0.
\]
The four real coordinates retain the two-variable geometry of
the prime-pair problem.

\emph{Hilbert--P\'olya.} The idea of Hilbert--P\'olya seeks a selfadjoint
operator to prove a Riemann Hypothesis. We reverse this approach:
assuming GRH, we use such an operator to derive arithmetic consequences.
For $N>e^{10}>22026$ satisfying $\mathrm{GRH}(N)$, we realize the
Hilbert--P\'olya idea in the Goldbach setting by explicitly constructing
the selfadjoint operator $T_N$ from the residual in \eqref{qp:operator}.
This operator is compact and Hilbert--Schmidt; the construction is
detailed in Subsection~\ref{story:hilbert-polya}. We then use either of
the additional spectral hypotheses below to bound its operator norm.
The resulting bound makes the arithmetic error
$|\mathcal G(N)-c_\phi\mathfrak S(N)|$ smaller than the positive main term
$c_\phi\mathfrak S(N)$ for all sufficiently large even $N$.

\emph{Riemann.} Riemann connected primes with a complex function.
For prime pairs, our passage to a function of four real variables
goes through the operator: the function is built from its spectrum.
Write $\mu_j(N)$ for the nonzero real eigenvalues of $T_N$,
with multiplicity, so that $\sum_j\mu_j^2(N)<\infty$. Set
\begin{equation}\label{v3i:spectral-scale}
 \delta_N=\sqrt{\frac{\mathfrak S(N)}N}(\log N)^{3/2}.
\end{equation}
Define the symmetric spectral function
\begin{equation}\label{qp:spectral-function}
 \mathfrak X_N(q)=
 \prod_j\left(1-\frac{4\mu_j^2(N)}{\delta_N^2}q(1-q)\right),
 \qquad q\in\mathbb H,
\end{equation}
where $\mathbb H$ denotes the quaternions and $\Re q$ denotes
the real part of $q$. The product converges locally uniformly; an empty
product is $1$. Its symmetry under $q\mapsto1-q$ is similar to that of
Riemann's completed function~\cite{Riemann1859}.
This is an intrinsic slice-regular function of the
spectral variable $q$; the Fueter regularity above concerns the
arithmetic coordinates. The locations of its zeros determine the
magnitudes of the nonzero eigenvalues and hence $\|T_N\|$.
Bounds on this norm control $|\mathcal G(N)-c_\phi\mathfrak S(N)|$
through the Hardy quadratic-form identity below.

\emph{Hardy.} Linnik's dispersion method~\cite{Linnik1963}
exemplifies the power of quadratic estimates in additive problems.
Here the key step is to express the exact signed error
as a quadratic form in a Hardy space. The weight determines $v_N$ with
$\|v_N\|^2\ll_\phi1$. Before applying Cauchy--Schwarz, we have the exact identity
\[
 c_\phi\mathfrak S(N)-\mathcal G(N)=\langle T_Nv_N,v_N\rangle.
\]
Applying Cauchy--Schwarz then gives an upper bound for the absolute error:
\[
 |\langle T_Nv_N,v_N\rangle|
 \le\|T_N\|\,\|v_N\|^2\ll_\phi\|T_N\|.
\]
Appendix~\ref{proof:app:numerical} tests the counting method numerically
at $N=30030$, recovering the exact number of Goldbach representations.
Section~\ref{story:section2} develops these five steps.

To turn the last inequality into arithmetic estimates, we now
formulate two conditions on the zeros of $\mathfrak X_N$ that
give bounds for $\|T_N\|$. The first, Hypothesis S, excludes zeros
sufficiently close to $1$ and yields an asymptotic formula.

\begin{hypothesisS}
Fix beforehand a positive function $h(N)\to\infty$ such that,
eventually,
\begin{equation}\label{v3i:S-width}
 \eta_N=\frac{h^2(N)(\log N)^3}{N\mathfrak S(N)}\le\frac14.
\end{equation}
For every sufficiently large even $N$ satisfying
$\mathrm{GRH}(N)$, the function $\mathfrak X_N(q)$ has no
zeros in $\Re q>1-\eta_N$.
\end{hypothesisS}

In fact, GRH alone already gives a zero-free region for $\mathfrak X_N$
of the same form. Its width, however, falls short of that required by
Hypothesis S, so the resulting estimate is not strong enough to make
the Goldbach error $o(\mathfrak S(N))$.

\begin{theorem}\label{thm:goldbach-S}
Assume $\mathrm{GRH}(N)$ for every sufficiently large even
$N$, and Hypothesis S with the prescribed function $h$. Then
\begin{equation}\label{v3i:weak-result}
 \mathcal G(N)=c_\phi\mathfrak S(N)+
 O_\phi\!\left(\frac{\mathfrak S(N)}{h(N)}\right).
\end{equation}
In particular, $\mathcal G(N)\sim c_\phi\mathfrak S(N)$.
The implied constant is independent of $N$ and $h$.
\end{theorem}

Although Hypothesis S already suffices for our present purpose
of obtaining Goldbach representations, we nevertheless venture
the following stronger, ultimate conjecture on the spectral zeros.
Together with GRH, it would yield what we expect to be the best
possible remainder bound in the weighted Goldbach asymptotic formula.

\begin{hypothesisR}
For every sufficiently large even $N$ satisfying
$\mathrm{GRH}(N)$, every zero of $\mathfrak X_N(q)$ has
real part $1/2$.
\end{hypothesisR}

Hypothesis R implies Hypothesis S for every admissible choice
of $h$ and gives the following sharper estimate.

\begin{theorem}\label{thm:goldbach-R}
Assume $\mathrm{GRH}(N)$ for every sufficiently large even
$N$, and Hypothesis R. Then
\begin{equation}\label{v3i:strong-result}
 \mathcal G(N)=c_\phi\mathfrak S(N)+
 O_\phi\!\left(
       \sqrt{\frac{\mathfrak S(N)}N}(\log N)^{3/2}\right).
\end{equation}
\end{theorem}

Either theorem gives
$\gg_\phi\mathfrak S(N)N/(\log N)^2$ ordered representations
$N=p_1+p_2$, with both primes in $[N/3,2N/3]$, for every
sufficiently large even $N$.

The method developed here also extends naturally to other additive
problems. These extensions will be discussed in future work.

\medskip\noindent\textit{Notation.}
We use standard notation from analytic number theory. In
particular, $p$ denotes a prime, $(a,b)$ the greatest common
divisor of $a$ and $b$, and $d\mid n$ divisibility.
The symbols $\Lambda$, $\varphi$, and $\mu$ denote the von
Mangoldt, Euler totient, and M\"obius functions, respectively.
The notations $O,o,\ll,\gg,\asymp$, and $\sim$ have their
standard asymptotic meanings; subscripts indicate the allowed
dependence of implied constants.
We write $\mathbf1_{\mathcal A}$ for the indicator of the
condition $\mathcal A$, and $\int_{(\sigma)}$ for integration
upwards along $\Re s=\sigma$.
All logarithms are natural, and all limits in $N$ are taken
through even integers.

\section{Idea and outline of the proof}\label{story:section2}

The proof has two tasks: represent
$c_\phi\mathfrak S(N)-\mathcal G(N)$ as a quadratic form of
$T_N$ at a vector of uniformly bounded norm, and obtain the
required bound for $\|T_N\|$ from a condition on spectral zeros.
The construction uses $\mathrm{GRH}(N)$; the additional norm
bound comes from Hypothesis S or R.

The five steps below complete the argument. We will prove the analytic
results used here in
Sections~\ref{proof:sec:goldbach-projection}--\ref{proof:sec:v3-mellin-operator}.
We introduce notation as it enters the construction. In the final Hardy
step, we will establish the exact quadratic-form identity and a bound for
$\|v_N\|$ that is uniform in $N$. These results convert the bound for $\|T_N\|$ into
the required estimate for the Goldbach error. Removing the contribution
of proper prime powers then yields representations of $N$ as a sum of
two primes for all sufficiently large even $N$.
We retain the fixed weights and arithmetic notation of Section~\ref{story:section1}.

\subsection{Dirichlet: an exact projection of the Goldbach correlation}
\label{story:dirichlet}

\subsubsection{The finite projection}

The relation $n+m=N$ implies $n\equiv-m\ \mathrm{mod}\ N$. Let $X_N^-$ be
all odd characters modulo $N$, including imprimitive characters,
extended by $0$ at nonunits. Their orthogonality gives
\begin{equation}\label{gb:orthogonality}
 \frac2{\varphi(N)}\sum_{\chi\in X_N^-}
          \chi(n)\overline{\chi(m)}
 =\mathbf1_{(nm,N)=1}
    \left(\mathbf1_{n\equiv m\ \mathrm{mod}\ N}
               -\mathbf1_{n\equiv-m\ \mathrm{mod}\ N}\right).
\end{equation}
On the support of the weight, $|n-m|\le N/3$ and
$2N/3\le n+m\le4N/3$. The two congruences thus mean $n=m$
and $n+m=N$, respectively. Multiplying by
$\Lambda(n)\Lambda(m)\phi(n/N)\phi(m/N)/N$ and summing gives
\begin{equation}\label{gb:projection-identity}
 \mathcal G(N)=\mathcal D(N)-\mathcal P(N),
\end{equation}
where
\begin{equation}\label{gb:finite-moments}
 \begin{aligned}
 \mathcal D(N)&=\frac1N\sum_{(n,N)=1}\Lambda^2(n)\phi^2(n/N),\\
 \mathcal P(N)&=\frac2{\varphi(N)}
       \sum_{\chi\in X_N^-}|P_\chi|^2,
 \end{aligned}
\end{equation}
and
\begin{equation}\label{gb:prime-polynomial}
 P_\chi=\sum_{n\ge1}\frac{\Lambda(n)\chi(n)}{\sqrt n}\Phi(n/N)
        =\frac1{\sqrt N}\sum_{n\ge1}\Lambda(n)\chi(n)\phi(n/N).
\end{equation}
Here the companion weight used throughout is defined by
\begin{equation}\label{eq:fixed-goldbach-weight}
 \Phi(u)=\sqrt u\,\phi(u).
\end{equation}
The factor $1/\sqrt N$ gives the required normalization after
squaring. We will prove the unconditional projection identity in
Lemma~\ref{proof:v9:finite-projection}. It reduces the problem
to estimating the character moment relative to the exact diagonal.

\subsubsection{The analytic kernel and Mellin inversion}

For $\Re s,\Re z>1/2$, define
\begin{equation}\label{eq:goldbach-moment}
 \mathscr P_N(s,z)=\frac2{\varphi(N)}\sum_{\chi\in X_N^-}
       \frac{L'}L(1/2+s,\chi)
       \frac{L'}L(1/2+z,\bar\chi).
\end{equation}
The variables $s,z$ are independent, and the bar conjugates only
the character. We write $w=s+z$. Use the Mellin convention
\begin{equation}\label{gb:mellin-convention}
 \widetilde{\Phi}(s)=\int_0^\infty \Phi(u)u^{s-1}\d u,
 \quad
 \Phi(u)=\frac1{2\pi i}\int_{(\sigma)}\widetilde{\Phi}(s)u^{-s}\d s
 \quad(\sigma\in\mathbb R).
\end{equation}
The series $-L'/L(1/2+s,\chi)=\sum_n\Lambda(n)\chi(n)n^{-1/2-s}$
and Mellin inversion represent $P_\chi$ by an integral. Applying
this to $P_\chi P_{\bar\chi}=|P_\chi|^2$ and averaging gives
\begin{equation}\label{gb:mellin-moment}
 \mathcal P(N)=\frac1{(2\pi i)^2}
 \int_{(\sigma)}\int_{(\sigma)}
 N^{s+z}\widetilde{\Phi}(s)\widetilde{\Phi}(z)\mathscr P_N(s,z)
          \,\mathrm ds\,\mathrm dz.
\end{equation}
The two minus signs cancel. This formula holds
unconditionally for $\sigma>1/2$ by absolute convergence.

Under $\mathrm{GRH}(N)$, all inducing primitive characters have
their nontrivial zeros on $\Re s=1/2$. Since odd characters are
nonprincipal, $L'/L(1/2+s,\chi)$ is holomorphic for $\Re s>0$.
We will prove the joint holomorphy of $\mathscr P_N$ in
Proposition~\ref{proof:v3:kernel-holomorphy} and extend
\eqref{gb:mellin-moment} to $0<\sigma\le1/9$ in
Lemma~\ref{proof:v9:mellin-signed-families}, where we will justify
the contour shift using the required growth estimates.
GRH enters through this contour shift.

\subsubsection{The families to be modeled}
\label{story:dirichlet-families}

For $\Re s,\Re z>1/2$, the unweighted analytic kernel has the
absolutely convergent expansion
\[
 \mathscr P_N(s,z)=
 \sum_{\substack{n,m\ge1\\(nm,N)=1}}
 \frac{\Lambda(n)\Lambda(m)}{n^{1/2+s}m^{1/2+z}}
 \left(\mathbf1_{n\equiv m\ \mathrm{mod}\ N}
              -\mathbf1_{n\equiv-m\ \mathrm{mod}\ N}\right).
\]
We will prove this identity in
Proposition~\ref{proof:v3:kernel-holomorphy} using the
logarithmic-derivative series and odd-character orthogonality.
The positive congruence gives the diagonal $n=m$ and two difference
families; the negative congruence gives the sum family. The three
off-diagonal families are therefore
\[
 \begin{array}{c|c}
 \text{relation},\ k\ge1 & \text{sign in }\mathscr P_N\\ \hline
 m=n+kN & +\\
 n=m+kN & +\\
 n+m=kN & -
 \end{array}
\]
After Mellin inversion restores the weight, its support gives
$|n-m|\le N/3$ and $n+m\le4N/3$. Thus both difference families
and the sum terms with $k\ge2$ vanish; among the off-diagonal
families, only $n+m=N$ remains. Before this inversion, however,
the analytic model must retain all three families. In the next step,
we will keep the diagonal exactly and model their signed contribution;
subtracting this model will isolate the Goldbach error.

\subsection{Euler: the Beta model and the analytic residual}
\label{story:euler}

\subsubsection{The signed continuous model}

By \eqref{gb:projection-identity}, a model for $\mathscr P_N$ must
have weighted integral $\mathcal D(N)-c_\phi\mathfrak S(N)$.
We retain the exact diagonal and model the three off-diagonal
families using their local densities and Euler's Beta integral.

Consider the three families in the table in
Subsection~\ref{story:dirichlet-families}. Fix $k\ge1$ and set $h=kN$.
In the first two rows, $h=|n-m|$ is the positive difference; in the
third, $h=n+m$ is the sum. All three have the same local-density
factor $\mathfrak S(h)$. For $0<\Re s,\Re z<1/2$ and $w=s+z$,
their continuous integrals, in the order of the table, are
\[
 \begin{aligned}
 \int_0^\infty x^{-1/2-s}(x+h)^{-1/2-z}\,\mathrm{d}x
    &=h^{-w}B(1/2-s,w),\\
 \int_0^\infty x^{-1/2-z}(x+h)^{-1/2-s}\,\mathrm{d}x
    &=h^{-w}B(1/2-z,w),\\
 \int_0^h x^{-1/2-s}(h-x)^{-1/2-z}\,\mathrm{d}x
    &=h^{-w}B(1/2-s,1/2-z).
 \end{aligned}
\]
We will prove these evaluations and derive the common density factor
in Lemma~\ref{proof:v9:signed-beta-model}. Combining the three terms
with the signs $+,+,-$ from the table gives the model
$\mathfrak S(h)h^{-w}K_-(s,z)$, where
\begin{equation}\label{eq:beta-kernel}
 K_-(s,z)=B(1/2-s,s+z)+B(1/2-z,s+z)
                         -B(1/2-s,1/2-z).
\end{equation}
Here $B(s,z)=\Gamma(s)\Gamma(z)/\Gamma(s+z)$ is Euler's Beta function.
The local-density calculation specifies this continuous model without
assuming an asymptotic formula for prime pairs.

To collect the scales $h=kN$, introduce
\begin{equation}\label{v3i:series}
 \mathscr S_N(w)=\sum_{k\ge1}\frac{\mathfrak S(kN)}{k^w}
 \quad(\Re w>1).
\end{equation}
The exact diagonal, the density Euler product, and the resulting
off-diagonal model are
\begin{equation}\label{eq:goldbach-arithmetic-kernels}
 \begin{aligned}
 {\mathscr D}_N(w)&=\sum_{\substack{n\ge1\\(n,N)=1}}
                           \frac{\Lambda^2(n)}{n^{1+w}},\\
 \mathscr S_N(w)&=\mathfrak S(N)\zeta(w)P_N(w),
 \quad P_N(w)=\prod_{\substack{p>2\\p\nmid N}}
                       \left(1+\frac{p^{-w}}{p-2}\right),\\
 {\mathscr K}_N(s,z)&=N^{-s-z}\mathscr S_N(s+z)K_-(s,z).
 \end{aligned}
\end{equation}
The factor $N^{-w}$ comes from $(kN)^{-w}$; $P_N$ is an Euler
product, distinct from $P_\chi$. We will prove the Euler product identity
$\mathscr S_N(w)=\mathfrak S(N)\zeta(w)P_N(w)$ for the density series
in Lemma~\ref{proof:lem:singular-N}. There we will also continue
$\mathscr S_N$ meromorphically to $\Re w>0$, allowing its use on the
Beta contours.

\subsubsection{Continuation and the exact main term}

Gamma factorization and the zeta functional equation give
\begin{equation}\label{v3:model-functional-form}
 {\mathscr K}_N(s,z)=\mathfrak S(N)\left(\frac\pi N\right)^{s+z}
 P_N(s+z)\zeta(1-s-z)\mathcal X_1(s)\mathcal X_1(z),
\end{equation}
where $\mathcal X_1(s)=\Gamma((3/2-s)/2)/\Gamma((3/2+s)/2)$.
We will prove in Proposition~\ref{proof:v3:model-geometry} that this
expression is holomorphic on $0<\Re s,\Re z<3/2$, unconditionally.

When the weight is restored, $N^{s+z}$ cancels the model's
$N^{-s-z}$. At integer scales, support and symmetry give
\[
 \int_0^\infty\phi(u)\phi(u+k)\d u=0\quad(k\ge1),
 \qquad
 \int_0^k\phi(u)\phi(k-u)\d u=
 \begin{cases}
 c_\phi,&k=1,\\
 0,&k\ge2.
 \end{cases}
\]
Hence only the negative sum term at $k=1$ survives, with density
$\mathfrak S(N)$. For the continued model this evaluation requires
a contour shift: the combined continuous mass of the two difference
terms and the negative sum term is zero, cancelling the density
series pole at $w=1$. We will justify this contour shift and prove
the following identity for $0<\sigma\le1/9$ in
Lemma~\ref{proof:gb:main-term-inversion}:
\[
 \frac1{(2\pi i)^2}\int_{(\sigma)}\int_{(\sigma)}
 N^{s+z}\widetilde{\Phi}(s)\widetilde{\Phi}(z)
 {\mathscr K}_N(s,z)\,\d s\,\d z
 =-c_\phi\mathfrak S(N).
\]
We will also evaluate the diagonal integral as $\mathcal D(N)$ in the same lemma.
Thus the model has exactly the required weighted integral.

\subsubsection{The arithmetic residual}

Define
\[
 \mathscr E_N(s,z)
 =\mathscr P_N(s,z)-{\mathscr D}_N(s+z)-{\mathscr K}_N(s,z).
\]
We will prove in Proposition~\ref{proof:v7:residual-holomorphy} that,
under $\mathrm{GRH}(N)$, the residual is holomorphic on
$0<\Re s,\Re z<3/2$. Combining the kernels'
three Mellin evaluations with \eqref{gb:projection-identity} gives
\[
 \mathcal G(N)=c_\phi\mathfrak S(N)
 -\frac1{(2\pi i)^2}\int_{(\sigma)}\int_{(\sigma)}
 N^{s+z}\widetilde{\Phi}(s)\widetilde{\Phi}(z)
 \mathscr E_N(s,z)\,\d s\,\d z.
\]
Indeed, the residual integral equals
$\mathcal P(N)-\mathcal D(N)+c_\phi\mathfrak S(N)
=c_\phi\mathfrak S(N)-\mathcal G(N)$, fixing the sign.
The task is now to bound this integral by $o(\mathfrak S(N))$.
We next express it as a quadratic form of a compact selfadjoint
operator, whose norm will control the error.

\subsection{Hilbert--P\'olya: reversing the direction of dependence}
\label{story:hilbert-polya}

\subsubsection{The spectral strategy and the role of GRH}
\label{story:polya-strategy}

The Hilbert--P\'olya programme seeks a selfadjoint realization
whose real spectral parameters would place zeta zeros on the
critical line. Here the direction is reversed: for even
$N>e^{10}$ we assume $\mathrm{GRH}(N)$, construct an operator
from the arithmetic residual, and ask how its spectrum controls
Goldbach. GRH supplies holomorphy; the further spectral hypotheses
will supply smallness.

Indeed, every odd character is nonprincipal, and GRH makes
$g_\chi(s)=L'/L(1/2+s,\chi)$ holomorphic for $\Re s>0$.
Together with the diagonal and the continued Beta model, this
makes $\mathscr E_N$ holomorphic on $0<\Re s,\Re z<3/2$.
Conjugating its second variable gives a Hermitian kernel, since
\begin{equation}\label{qp:Gram}
 \mathscr P_N(s,\bar z)=\frac2{\varphi(N)}\sum_{\chi\in X_N^-}
                 g_\chi(s)\overline{g_\chi(z)},
 \quad g_\chi(s)=\frac{L'}L(1/2+s,\chi),
\end{equation}
and the diagonal and Beta model have the same symmetry.
Thus
\[
 \mathscr E_N(s,\bar z)=\overline{\mathscr E_N(z,\bar s)}.
\]
Subtraction need not preserve positivity, so eigenvalues of either
sign must be allowed.

\subsubsection{The Hardy strip and the operator}
\label{story:polya-space}

Set
\begin{equation}\label{hd:strip}
 \alpha_N=\frac1{\log N},\quad \beta=\frac18,\quad \sigma_*=\frac19,
 \quad \Delta_N=\beta-\alpha_N,\quad
 S_N=\{s:\alpha_N<\Re s<\beta\},
\end{equation}
and write
\begin{equation}\label{hd:space-and-boundary}
 \mathcal H_N=H^2(S_N),\quad
 \Gamma_N=(\alpha_N+i\mathbb R)\,\sqcup\,(\beta+i\mathbb R).
\end{equation}
Here $\mathcal H_N$ consists of holomorphic functions whose paired
vertical $L^2$ means remain bounded as the two lines approach
the boundaries. Their traces form a closed subspace of
$L^2(\Gamma_N)$, with inner product
\begin{equation}\label{hd:inner-product}
 \langle h,g\rangle_{\mathcal H_N}
       =\int_{\Gamma_N}h(s)\overline{g(s)}\,|\mathrm{d}s|,
\end{equation}
linear in the first variable. Both boundaries carry positive arc
length, with no $1/(2\pi)$ factor. The choice
$\alpha_N<1/10$ leaves $\sigma_*$ strictly inside the strip and
keeps $N^{2\alpha_N}=e^2$, which will make the weight functional
uniformly bounded.

\label{story:operator}
Insert decay at large heights and define
\begin{equation}\label{qp:operator}
 (T_Nh)(s)=\int_{\Gamma_N}
       \frac{\mathscr E_N(s,\bar z)}{(1+s)(1+\bar z)}h(z)\,|\mathrm{d}z|.
\end{equation}
We will prove in Proposition~\ref{proof:bd:HS} that the boundary
kernel is square integrable and that the operator maps into the
Hardy trace space. Hermitian symmetry then makes $T_N$
selfadjoint; square integrability makes it Hilbert--Schmidt and
hence compact. Consequently its nonzero eigenvalues $\mu_j(N)$,
counted with multiplicity, satisfy
\[
 \mu_j(N)\in\mathbb R,\qquad
 \sum_j\mu_j^2(N)=\|T_N\|_{\mathrm{HS}}^2<\infty,
 \qquad \|T_N\|=\sup_j|\mu_j(N)|.
\]
The last supremum is understood as zero when $T_N=0$.
These facts hold for fixed $N$; the arithmetic question is the
size of this norm as $N$ varies.

\subsubsection{The bridge to the arithmetic error}

For $\alpha_N<\sigma\le1/9$, we define the linear functional
$\ell_N:\mathcal H_N\to\mathbb C$ using the Mellin transform
$\widetilde\Phi$ of the companion weight $\Phi$ from
\eqref{eq:fixed-goldbach-weight}:
\[
 \ell_N(h)=\frac1{2\pi i}\int_{(\sigma)}
       N^s\widetilde\Phi(s)(1+s)h(s)\,\mathrm{d}s.
\]
It has a Riesz vector $v_N$, characterized by
$\ell_N(h)=\langle h,v_N\rangle$, with
$\mathcal Q_N:=\|v_N\|^2\ll_\phi1$ uniformly in $N$.
We will verify this bound and derive the following identity for the
residual in Subsection~\ref{story:hardy-quadratic}:
\[
 \mathcal G(N)=c_\phi\mathfrak S(N)-\langle T_Nv_N,v_N\rangle.
\]
Now Cauchy--Schwarz gives
\begin{equation}\label{v3f:master-error-formula}
 |\mathcal G(N)-c_\phi\mathfrak S(N)|
 \le\|T_Nv_N\|\,\|v_N\|
 \le\mathcal Q_N\|T_N\|.
\end{equation}
Thus $\|T_N\|=o(\mathfrak S(N))$ suffices for the relative
asymptotic, while $\|T_N\|\le\delta_N$ gives the error
$O_\phi(\delta_N)$ in Theorem~\ref{thm:goldbach-R}.
Before applying the inequality, the quadratic form remains the
exact arithmetic correction. We next express the required norm
bounds as conditions on spectral zeros.

\subsection{Riemann: spectral zeros and operator norm bounds}
\label{story:riemann}

For even $N>e^{10}$ satisfying $\mathrm{GRH}(N)$, inequality
\eqref{v3f:master-error-formula}, together with $\mathcal Q_N\ll_\phi1$,
gives the master estimate
\begin{equation}\label{v3f:master-estimate}
 |\mathcal G(N)-c_\phi\mathfrak S(N)|\ll_\phi\|T_N\|.
\end{equation}
Thus bounding the operator norm directly controls the error in the
weighted Goldbach count. We will express the required norm bounds as
conditions on spectral zeros and derive the two asymptotic formulas
under Hypotheses S and R, respectively.

\subsubsection{The spectral product and its zeros}
\label{story:riemann-spectrum}

Recall $\delta_N^2=\mathfrak S(N)(\log N)^3/N$ and the product
\eqref{qp:spectral-function}:
\[
 \mathfrak X_N(q)=\prod_j
 \left(1-\frac{4\mu_j^2(N)}{\delta_N^2}q(1-q)\right).
\]
The squared eigenvalues are summable, so the product converges
locally uniformly. On the usual complex slice, it has the Fredholm
determinant representation
\[
 \mathfrak X_N(q)=\det\!\left(\operatorname{Id}
              -4\delta_N^{-2}q(1-q)T_N^2\right).
\]
Its real coefficients give a consistent intrinsic slice-regular
extension to $q\in\mathbb H$. This spectral variable is independent
of the arithmetic variables $(s,z)$; its regularity is distinct
from the Fueter regularity of $\mathscr E_N(s,\bar z)$ in the
coordinates $\omega=s+zj$.

\label{story:riemann-zeros}
The factors immediately give
\begin{equation}\label{sym:normalization}
 \mathfrak X_N(0)=\mathfrak X_N(1)=1,
\end{equation}
\begin{equation}\label{sym:functional-equation}
 \mathfrak X_N(1-q)=\mathfrak X_N(q),
\end{equation}
and
\begin{equation}\label{sym:conjugation}
 \mathfrak X_N(\bar q)=\overline{\mathfrak X_N(q)}.
\end{equation}
For real $q\notin(0,1)$ every factor is at least $1$, so there
are no real zeros there. To find all zeros, put
$d_j=\delta_N^2/\mu_j^2(N)$ and write $q=x+yI$, where
$x,y\in\mathbb R$ and $I$ is a unit imaginary quaternion.
The factor equation is
\[
 q(1-q)=x(1-x)+y^2+(1-2x)yI=d_j/4.
\]
Its imaginary part forces $x=1/2$ whenever $y\ne0$, and the
remaining quadratic gives
\[
 \begin{cases}
 \Re q=\frac12,\quad |\Im q|=\frac12\sqrt{d_j-1},&d_j>1,\\[1mm]
 q=\frac12\quad\text{(a double zero of the factor)},&d_j=1,\\[1mm]
 q=\frac12\pm\frac12\sqrt{1-d_j},&0<d_j<1.
 \end{cases}
\]
The nonreal zeros form two-spheres on $\Re q=1/2$; every zero
off that hyperplane is real and comes from $|\mu_j|>\delta_N$.
Compactness allows only finitely many such exceptional pairs.
Absolute convergence excludes additional zeros of the infinite
product. We will give the convergence argument and prove the full
symmetry statements in Section~\ref{proof:sec:spectral-analysis}.

\subsubsection{The zero--norm dictionary}

Fix $0<\eta<1/2$. Only the larger root of an exceptional pair
can enter $\Re q>1-\eta$. Keeping it outside means
\[
 \frac12+\frac12\sqrt{1-\delta_N^2/\mu_j^2}
 \le1-\eta,
\]
which is equivalent to
\[
 |\mu_j|\le\frac{\delta_N}{2\sqrt{\eta(1-\eta)}}.
\]
Nonexceptional eigenvalues already satisfy this inequality.
Taking the supremum therefore gives the exact dictionary
\[
 \mathfrak X_N(q)\ne0\ (\Re q>1-\eta)
 \quad\Longleftrightarrow\quad
 \|T_N\|\le\frac{\delta_N}{2\sqrt{\eta(1-\eta)}}.
\]
The half-space is open, so boundary zeros and equality in the
norm bound are allowed. In particular, a bound
$\|T_N\|\le M_N$ with $M_N>\delta_N$ guarantees width
\begin{equation}\label{outline:norm-width-conversion}
 \eta(M_N)=\frac{1-\sqrt{1-\delta_N^2/M_N^2}}2,
 \qquad \eta(M_N)(1-\eta(M_N))=\frac{\delta_N^2}{4M_N^2}.
\end{equation}
Thus the width has order $\delta_N^2/M_N^2$.

\subsubsection{What GRH alone supplies}
\label{story:riemann-grh}

In Theorem~\ref{proof:grh:main-theorem}, we will prove the uniform
analytic estimate
\begin{equation}\label{grh:operator-bound}
 \|T_N\|\le\|T_N\|_{\mathrm{HS}}\ll(\log N)^2
\end{equation}
for sufficiently large even $N$ satisfying $\mathrm{GRH}(N)$.
Its proof splits the operator according to
\[
 \mathscr E_N=\mathscr P_N-\mathscr D_N-\mathscr K_N.
\]
Writing $\mathcal P_N$, $\mathcal D_N$, and $\mathcal K_N$
for these three kernels with the weighting and integration of
\eqref{qp:operator}, one has
\[
 T_N=\mathcal P_N-\mathcal D_N-\mathcal K_N.
\]
The character average controls the first component, positivity
controls the diagonal, and the continued Beta formula controls
the model. Their separate estimates each have size
$O((\log N)^2)$ at $\alpha_N=1/\log N$; they do not measure
cancellation between the components.

Putting $M_N=C(\log N)^2$ in the dictionary gives a guaranteed
zero-free region
\begin{equation}\label{grh:zero-free-width}
 \Re q>1-\eta_N^{\mathrm{GRH}},\qquad
 \eta_N^{\mathrm{GRH}}\asymp\frac{\mathfrak S(N)}{N\log N}.
\end{equation}
The constant is chosen sufficiently small using the absolute
constant implied in \eqref{grh:operator-bound}. This available norm
bound is too large to make the Goldbach error smaller than its
main term; S and R impose the required improvement.

\subsubsection{Completing the implications of S and R}
\label{story:riemann-hypotheses}

For any zero-free width $0<\eta_N\le1/4$, the dictionary and
the master estimate \eqref{v3f:master-estimate} give
\begin{equation}\label{v3f:relative-width}
 \frac{|\mathcal G(N)-c_\phi\mathfrak S(N)|}{c_\phi\mathfrak S(N)}
 \ll_\phi\sqrt{\frac{(\log N)^3}{N\eta_N\mathfrak S(N)}}.
\end{equation}
It therefore suffices for the asymptotic that
\begin{equation}\label{v3f:minimal-width}
 \frac{N\eta_N\mathfrak S(N)}{(\log N)^3}\longrightarrow\infty.
\end{equation}
Writing the left side as $h(N)^2$ explains precisely the scale
chosen in Hypothesis S.

\emph{Hypothesis S.} Write $\mathrm S(N;h)$ for its zero-exclusion
assertion at the fixed modulus, with the prescribed admissible
$h$ and width \eqref{v3i:S-width}. Substitution gives
\begin{equation}\label{v3i:S-norm}
 \mathrm S(N;h)\quad\Longleftrightarrow\quad
 \|T_N\|\le\frac{\mathfrak S(N)}{2h(N)\sqrt{1-\eta_N}}.
\end{equation}
Since $1-\eta_N\ge3/4$, the master estimate
\eqref{v3f:master-estimate} proves
$\mathcal G(N)=c_\phi\mathfrak S(N)+O_\phi(\mathfrak S(N)/h(N))$,
which is Theorem~\ref{thm:goldbach-S}. The implied constant is
independent of $N$ and $h$; the required lower bound on $N$
may depend on $\phi$ and the preassigned function $h$. The gain over the available
GRH width is
\begin{equation}\label{outline:S-gap}
 \frac{\eta_N}{\eta_N^{\mathrm{GRH}}}
 \asymp h^2(N)\frac{(\log N)^4}{\mathfrak S^2(N)}\longrightarrow\infty,
\end{equation}
using $\mathfrak S(N)\ll\log\log N$.
The excluded zeros are real pairs near the endpoints; the
Siegel-zero analogy concerns their location, without asserting
uniqueness, simplicity, or repulsion.

\emph{Hypothesis R.} Write $\mathrm R(N)$ for its assertion
at the fixed modulus. The zero classification shows directly that
\begin{equation}\label{v3i:R-equivalence}
 \mathrm R(N)\quad\Longleftrightarrow\quad\|T_N\|\le\delta_N.
\end{equation}
The master estimate \eqref{v3f:master-estimate} therefore gives
$\mathcal G(N)=c_\phi\mathfrak S(N)+O_\phi(\delta_N)$,
proving Theorem~\ref{thm:goldbach-R}. All exceptional real pairs
are excluded, so R also implies S for every admissible choice of
$h$.

In the following Hardy subsection, we will establish the
exact quadratic-form identity and the uniform bound for $\|v_N\|$
needed to convert the spectral estimates into these arithmetic conclusions.

\subsection{Hardy: the arithmetic error as an exact quadratic form}
\label{story:hardy}

The final step evaluates the residual at the vector selected by the
Goldbach weight. We obtain an exact identity first; Cauchy--Schwarz
then turns it into the estimate needed for the two theorems.
Throughout, assume $\mathrm{GRH}(N)$ and $N>e^{10}$, with
$\Phi(u)=\sqrt u\,\phi(u)$ as in \eqref{eq:fixed-goldbach-weight}.

\subsubsection{Mellin inversion and the arithmetic residual}
\label{story:mellin}

Mellin inversion of the prime polynomial gives
\begin{equation}\label{gb:mellin-P}
 P_\chi=-\frac1{2\pi i}\int_{(\sigma)}
    N^s\widetilde{\Phi}(s)\frac{L'}L(1/2+s,\chi)\,\mathrm ds.
\end{equation}
It holds first on $\sigma=1$ by absolute convergence, and then
on $0<\sigma\le1/9$: GRH excludes poles between the contours,
while the smooth weight makes the horizontal integrals vanish.
Multiplying the formulas for $\chi$ and $\bar\chi$ and averaging
over $X_N^-$ proves \eqref{gb:mellin-moment} on these lines.
The two subtracted kernels have the evaluations
\begin{align}
 \frac1{(2\pi i)^2}\int_{(\sigma)}\int_{(\sigma)}
  N^{s+z}\widetilde{\Phi}(s)\widetilde{\Phi}(z){\mathscr D}_N(s+z)\,\d s\,\d z
 &=\mathcal D(N),\label{gb:diagonal-inversion}\\
 \frac1{(2\pi i)^2}\int_{(\sigma)}\int_{(\sigma)}
  \widetilde{\Phi}(s)\widetilde{\Phi}(z)
  \mathscr S_N(s+z)K_-(s,z)\,\d s\,\d z
 &=-c_\phi\mathfrak S(N).\label{gb:beta-inversion}
\end{align}
We will prove both evaluations in
Lemma~\ref{proof:gb:main-term-inversion}.

The residual integral is therefore
\begin{equation}\label{v3w:residual}
 \begin{aligned}
 r_\sigma(N)&=\frac1{(2\pi i)^2}
 \int_{(\sigma)}\int_{(\sigma)}
 N^{s+z}\widetilde\Phi(s)\widetilde\Phi(z)\mathscr E_N(s,z)
 \,\mathrm ds\,\mathrm dz\\
 &=\mathcal P(N)-\mathcal D(N)+c_\phi\mathfrak S(N).
 \end{aligned}
\end{equation}
The exact projection \eqref{gb:projection-identity} now yields
\begin{equation}\label{v3w:exact}
 \mathcal G(N)=c_\phi\mathfrak S(N)-r_\sigma(N).
\end{equation}
In particular, this signed discrepancy is independent of $\sigma$;
write it as $r(N)$.

\subsubsection{The residual as a quadratic form}
\label{story:hardy-quadratic}

For an interior line $\alpha_N<\sigma<\beta$, define
\begin{equation}\label{bm:functional}
 \ell_N(h)=\frac1{2\pi i}\int_{(\sigma)}
          N^s\widetilde\Phi(s)(1+s)h(s)\,\mathrm{d}s,
 \qquad h\in\mathcal H_N.
\end{equation}
The factor $1+s$ compensates for the denominator inserted in
the operator kernel. Put $V(u)=\Phi(u)-u\Phi'(u)$, so that
$\widetilde V(s)=(1+s)\widetilde\Phi(s)$.
We will prove the following formula for the exact squared norm in
Proposition~\ref{proof:bm:norm}, using the Fourier--Laplace
representation of the strip Hardy space:
\begin{equation}\label{bm:norm-integral}
 \mathcal Q_N:=\|\ell_N\|^2
 =\frac1{2\pi}\int_0^\infty
       \frac{(Nu)^{2\alpha_N}|V(u)|^2}
            {1+(Nu)^{-2\Delta_N}}\,\frac{\mathrm{d}u}{u}.
\end{equation}
This representation also proves independence of the interior contour.
Since $\alpha_N=1/\log N$ and $V$ is supported in $[1/3,2/3]$,
\begin{equation}\label{bm:uniform-norm}
 \mathcal Q_N\le\frac{e^2}{2\pi}
       \int_0^\infty|\Phi(u)-u\Phi'(u)|^2\,\frac{\mathrm{d}u}{u}
       \ll_\phi1.
\end{equation}
The choice of the moving strip thus prevents a loss growing with $N$.

Let $v_N$ be the Riesz vector, with
$\ell_N(h)=\langle h,v_N\rangle$ and
$\|v_N\|^2=\mathcal Q_N$; the inner product is linear in its first
entry. Let $k_z$ be the evaluation vector, defined by
$h(z)=\langle h,k_z\rangle$. The operator kernel is
\[
 \langle T_Nk_z,k_s\rangle
   =\frac{\mathscr E_N(s,\bar z)}{(1+s)(1+\bar z)}.
\]
Apply $\ell_N$ to each kernel variable, conjugating the second.
The factors $1+s$ and $1+\bar z$ cancel; the reality of $\Phi$
then converts the second contour to the holomorphic Mellin contour.
This recovers precisely \eqref{v3w:residual}, and hence
\begin{equation}\label{bm:quadratic-form}
 r(N)=\langle T_Nv_N,v_N\rangle.
\end{equation}
Only now do we estimate. Cauchy--Schwarz gives
\begin{equation}\label{bm:direct-error}
 |r(N)|\le\|T_Nv_N\|\,\|v_N\|
       \le\mathcal Q_N\|T_N\|\ll_\phi\|T_N\|.
\end{equation}
Together with \eqref{v3w:exact}, this proves the master estimate
\eqref{v3f:master-estimate} and hence
Theorems~\ref{thm:goldbach-S} and \ref{thm:goldbach-R} under S and R.
Retaining the signed quadratic form gives the exact weighted
Goldbach identity.

\subsubsection{Goldbach representations}
\label{story:primes-and-zeros}

Goldbach representations follow readily from either
Theorem~\ref{thm:goldbach-S} or Theorem~\ref{thm:goldbach-R}:
under its hypotheses, every sufficiently large even integer $N$
can be written as $N=p_1+p_2$, with $p_1,p_2$ prime,
and the number of representations is $\gg_\phi N(\log N)^{-2}$.

\section{Dirichlet: an exact projection of the Goldbach correlation}
\label{proof:sec:goldbach-projection}\label{proof:v7:dirichlet}
\label{proof:sec:arithmetic-motivation}\label{proof:v7:main-thread}

We express the weighted Goldbach correlation as the difference
between a diagonal moment and a moment over odd characters modulo $N$.
The identity is finite and unconditional. We then pass to logarithmic
derivatives of Dirichlet $L$-functions and use GRH to move the Mellin
contours into the region needed for the continuous model.

Fix the weight $\phi$ used in \eqref{eq:goldbach-correlation}
and its companion $\Phi$ defined by \eqref{eq:fixed-goldbach-weight}.
The symbol $\varphi(N)$ denotes Euler's
totient, whereas $\phi$ always denotes the weight. With $\Lambda$
the von Mangoldt function, the correlation to be studied is
\[
 \mathcal G(N)=\frac1N
 \sum_{\substack{n+m=N\\(nm,N)=1}}
       \Lambda(n)\Lambda(m)\phi(n/N)\phi(m/N),
\]
as in~\eqref{eq:goldbach-correlation}.
The sum runs over ordered pairs of positive integers. Its expected
main term is $c_\phi\mathfrak S(N)$, where
\[
 c_\phi=\int_0^1\phi^2(u)\,\d u,
\]
as in~\eqref{eq:c0},
and
\[
 \mathfrak S(N)=2C_2\prod_{\substack{p\mid N\\p>2}}\frac{p-1}{p-2},
 \qquad C_2=\prod_{p>2}\left(1-\frac1{(p-1)^2}\right),
\]
as in~\eqref{eq:singular-series}.
We write
\[
 X_N^-=\{\chi\ \mathrm{mod}\ N:\chi(-1)=-1\}
\]
for all odd characters modulo $N$, including imprimitive characters,
and extend each character by zero when its argument is not coprime
to $N$. These conventions account for the coprimality condition in
the projection below.

\subsection{The finite projection and its size}
\label{proof:v7:weight-class}\label{proof:v7:finite-character-moment}

We will study the two moments
\[
 \begin{aligned}
 \mathcal D(N)&=\frac1N\sum_{(n,N)=1}\Lambda^2(n)\phi^2(n/N),\\
 \mathcal P(N)&=\frac2{\varphi(N)}\sum_{\chi\in X_N^-}|P_\chi|^2,
 \end{aligned}
\]
as in~\eqref{gb:finite-moments},
where $P_\chi$ is the character polynomial
\[
 P_\chi=\sum_{n\ge1}\frac{\Lambda(n)\chi(n)}{\sqrt n}\Phi(n/N)
       =\frac1{\sqrt N}\sum_{n\ge1}\Lambda(n)\chi(n)\phi(n/N),
\]
as in~\eqref{gb:prime-polynomial}.
Thus $\mathcal D(N)$ retains equal indices, while $\mathcal P(N)$
averages the squared character sums with the same normalization.

\begin{lemma}[Finite projection]\label{proof:v9:finite-projection}
For every even $N\ge4$, the moments just defined satisfy the exact identity
\[
 \mathcal G(N)=\mathcal D(N)-\mathcal P(N)
\]
as in~\eqref{gb:projection-identity},
and the bound
\begin{equation}\label{proof:dir:moment-bound}
 0\le\mathcal P(N)\le\mathcal D(N)\ll_\phi\log N.
\end{equation}
\end{lemma}

\begin{proof}
The selector $(1-\chi(-1))/2$ in the full character sum gives
\[
 \frac2{\varphi(N)}\sum_{\chi\in X_N^-}\chi(n)\overline{\chi(m)}
 =\frac1{\varphi(N)}\sum_{\chi\ \mathrm{mod}\ N}
       (1-\chi(-1))\chi(n)\overline{\chi(m)}.
\]
Full character orthogonality gives, with the coprimality restriction
coming from the zero extension,
\[
 \frac2{\varphi(N)}\sum_{\chi\in X_N^-}\chi(n)\overline{\chi(m)}
 =\mathbf1_{(nm,N)=1}
   \left(\mathbf1_{n\equiv m\ \mathrm{mod}\ N}
              -\mathbf1_{n\equiv-m\ \mathrm{mod}\ N}\right),
\]
as in~\eqref{gb:orthogonality}.
If the weight is nonzero, $N/3\le n,m\le2N/3$, so
$|n-m|\le N/3$ and $2N/3\le n+m\le4N/3$.
The only permitted multiples of $N$ are therefore $0$ for $n-m$
and $N$ for $n+m$. Thus
\[
 \mathbf1_{(nm,N)=1}\mathbf1_{n+m=N}
 =\mathbf1_{(n,N)=1}\mathbf1_{n=m}
  -\frac2{\varphi(N)}\sum_{\chi\in X_N^-}
                 \chi(n)\overline{\chi(m)}
\]
on this support. Multiply by
$\Lambda(n)\Lambda(m)\phi(n/N)\phi(m/N)/N$ and sum.
The diagonal is $\mathcal D(N)$, while the character term factors
as $P_\chi P_{\bar\chi}=|P_\chi|^2$, since $\phi$ is real.
This proves the finite projection with its exact normalization.
Since both $\mathcal G(N)$ and $\mathcal P(N)$ are nonnegative,
$0\le\mathcal P(N)\le\mathcal D(N)$. Chebyshev's estimate
$\sum_{n\le x}\Lambda(n)\ll x$ gives
\[
 \mathcal D(N)\le\frac{\log N}{N}\sum_{n\le N}\Lambda(n)\ll\log N.
\]
This proves the lemma.
\end{proof}

The arithmetic problem is to estimate the difference between these
two moments at the smaller scale $\mathfrak S(N)$.

\subsection{Holomorphy of the character kernel}
\label{proof:v3:sec-geometry}\label{proof:v7:correlation-residual}

To insert the Mellin weight into both character polynomials, we use
the two-variable kernel
\[
 \mathscr P_N(s,z)=\frac2{\varphi(N)}\sum_{\chi\in X_N^-}
                    g_\chi(s)g_{\bar\chi}(z),
\]
as in~\eqref{eq:goldbach-moment},
where $g_\chi$ is the logarithmic derivative
\[
 g_\chi(s)=\frac{L'}L(1/2+s,\chi),
\]
as in~\eqref{qp:Gram}.
The variables $s,z$ are independent; conjugation here acts only on
the character. Thus $\mathscr P_N(s,z)$ is holomorphic in both
variables wherever the logarithmic derivatives are holomorphic.
The Hermitian kernel used later is $\mathscr P_N(s,\bar z)$:
it is holomorphic in the first variable and antiholomorphic in
the second, as is usual for kernels on spaces of holomorphic functions. We keep these two
forms distinct when applying Mellin inversion.

Write $\chi^*$ for the primitive character inducing $\chi$ and
$d_\chi\mid N$ for its conductor. In this section,
$\mathrm{GRH}(N)$ means that every nontrivial zero of every
primitive Dirichlet $L$-function whose conductor divides $N$
has real part $1/2$. On $\Re s>1$,
$L(s,\chi)=\sum_{n\ge1}\chi(n)n^{-s}$, and its logarithmic derivative is
\[
 -\frac{L'}L(1/2+s,\chi)
 =\sum_{n\ge1}\frac{\Lambda(n)\chi(n)}{n^{1/2+s}}
 \qquad(\Re s>1/2).
\]

\begin{proposition}[The character kernel]\label{proof:v3:kernel-holomorphy}
For $\Re s,\Re z>1/2$, the kernel has the absolutely convergent
expansion
\begin{equation}\label{proof:dir:kernel-expansion}
 \mathscr P_N(s,z)=
 \sum_{\substack{n,m\ge1\\(nm,N)=1}}
 \frac{\Lambda(n)\Lambda(m)}{n^{1/2+s}m^{1/2+z}}
 \left(\mathbf1_{n\equiv m\ \mathrm{mod}\ N}
              -\mathbf1_{n\equiv-m\ \mathrm{mod}\ N}\right).
\end{equation}
The kernel $\mathscr P_N$ is meromorphic on $\mathbb C^2$, with
possible poles only where one of its logarithmic derivatives has a
pole. Under $\mathrm{GRH}(N)$ it is jointly holomorphic for
$\Re s,\Re z>0$ and satisfies
\begin{equation}\label{proof:v3:kernel-symmetries}
 \mathscr P_N(s,z)=\mathscr P_N(z,s),\qquad
 \overline{\mathscr P_N(s,z)}=\mathscr P_N(\bar s,\bar z).
\end{equation}
The Hermitian kernel $\mathscr P_N(s,\bar z)$ is positive
semidefinite, of rank at most $\varphi(N)/2$.
\end{proposition}

\begin{proof}
Expand both logarithmic derivatives in their absolutely convergent
Dirichlet series and apply the odd-character orthogonality formula
proved above.
The two signs from the logarithmic derivatives cancel; the remaining
minus sign comes from the odd-character projection. This proves
\eqref{proof:dir:kernel-expansion}.

Every odd character is nonprincipal, so its $L$-function is entire
and its logarithmic derivative is meromorphic, with poles at the zeros.
Under GRH the nontrivial zeros lie on $\Re s=1/2$; the primitive
trivial zeros are negative odd integers, and the additional Euler-factor
zeros of imprimitive characters lie on $\Re s=0$. Hence each
$g_\chi(s)$ is holomorphic for $\Re s>0$.
The character sum is finite, so products of these one-variable
functions are jointly holomorphic on the product domain.
No limiting interchange is needed to pass from the individual
logarithmic derivatives to the kernel.

Conjugation gives $g_{\bar\chi}(\bar z)=\overline{g_\chi(z)}$.
Consequently the Hermitian form of the kernel is
\[
 \mathscr P_N(s,\bar z)=\frac2{\varphi(N)}\sum_{\chi\in X_N^-}
                    g_\chi(s)\overline{g_\chi(z)},
\]
as in~\eqref{qp:Gram}.
This Gram representation proves positivity and the rank bound,
since $|X_N^-|=\varphi(N)/2$. Reindexing by
$\chi\mapsto\bar\chi$ gives \eqref{proof:v3:kernel-symmetries}.
This proves the proposition.
\end{proof}

\subsection{Bounds for logarithmic derivatives under GRH}

The contour shifts require vertical bounds on fixed positive strips.
We derive them from the primitive zero expansion, retaining the Euler
factors that distinguish a character modulo $N$ from its inducing character.

For a primitive odd character $\chi^*$ of conductor $d_\chi$, write
\[
 \mathcal L(u,\chi^*)=(d_\chi/\pi)^{(u+1)/2}
 \Gamma((u+1)/2)L(u,\chi^*).
\]
For the following results, see~\cite{IK2004}.

\begin{lemma}\label{proof:v9:primitive-zero-data}
For a primitive odd character $\chi^*$, the nontrivial zeros $\rho$,
counted with multiplicity, satisfy
\begin{align}
 \frac{\mathcal L'}{\mathcal L}(u,\chi^*)
 -\frac{\mathcal L'}{\mathcal L}(u_0,\chi^*)
 &=\sum_\rho\left(\frac1{u-\rho}-\frac1{u_0-\rho}\right),
       \label{proof:gb:hadamard-difference}\\
 \#\{\rho:t\le\Im\rho<t+1\}
 &\ll\log\bigl(d_\chi(|t|+2)\bigr),\label{proof:gb:local-zero-count}
\end{align}
where the first identity holds away from zeros and its series
converges absolutely and locally uniformly in $u,u_0$.
\end{lemma}

Under $\mathrm{GRH}(N)$, the nontrivial zeros of $L(s,\chi^*)$
have the form $1/2+i\gamma_\chi$, where $\chi^*$ is the primitive
character inducing $\chi$. Set
\begin{align}
 Z_\chi(s)&=\sum_{\gamma_\chi}
   \left(\frac1{s-i\gamma_\chi}-\frac1{3/2-i\gamma_\chi}\right),
       &&\Re s>0,\label{proof:gb:zero-resolvent}\\
 e_\chi(u)&=\sum_{\substack{p\mid N\\p\nmid d_\chi}}
       \frac{\chi^*(p)\log p}{p^u-\chi^*(p)},
       &&\Re u>0.\label{proof:gb:Euler-correction}
\end{align}

\begin{lemma}\label{proof:lem:zero-expansion}
Assume $\mathrm{GRH}(N)$. The series \eqref{proof:gb:zero-resolvent}
converges absolutely and locally uniformly for $\Re s>0$, and
\begin{equation}\label{proof:gb:zero-identity}
 \begin{split}
 g_\chi(s)={}&Z_\chi(s)+\frac{L'}L(2,\chi)
 -\frac12\left\{\frac{\Gamma'}\Gamma(3/4+s/2)
                       -\frac{\Gamma'}\Gamma(3/2)\right\}\\
 &+e_\chi(1/2+s)-e_\chi(2).
 \end{split}
\end{equation}
For every compact interval $I\subset(0,\infty)$,
\begin{equation}\label{proof:gb:logderivative-bound}
 g_\chi(\sigma+it)\ll_I\log^2\bigl(N(|t|+2)\bigr)
 \qquad(\sigma\in I,\ t\in\mathbb R,\ \chi\in X_N^-).
\end{equation}
\end{lemma}

\begin{proof}
The paired summand in \eqref{proof:gb:zero-resolvent} is
$(3/2-s)/((s-i\gamma_\chi)(3/2-i\gamma_\chi))$.
The local zero count therefore proves locally uniform absolute
convergence. Apply \eqref{proof:gb:hadamard-difference} at
$u=1/2+s$ and $u_0=2$. The conductor terms cancel.
Subtracting the Gamma terms and using
\[
 L(u,\chi)=L(u,\chi^*)
       \prod_{\substack{p\mid N\\p\nmid d_\chi}}
                           (1-\chi^*(p)p^{-u})
\]
gives \eqref{proof:gb:zero-identity}.
On a compact subset of the right half-plane, the real parts stay
uniformly positive. Thus the finitely many low zeros cannot spoil
the local uniform convergence of the paired series.

For the bound, let $I=[\delta,M]$ with $\delta>0$ and put
$T=|t|+2$. For $|\gamma_\chi|\le2T$, group each fraction in
\eqref{proof:gb:zero-resolvent} by unit distance from its imaginary
part. Equation~\eqref{proof:gb:local-zero-count} gives
\[
 \ll_I\log\bigl(N(T+2)\bigr)
       \left(1+\sum_{1\le j\le3T+2}\frac1j\right)
 \ll_I\log^2\bigl(N(T+2)\bigr).
\]
The denominators nearest $t$ are bounded below by $\delta$.
For $|\gamma_\chi|>2T$, keep the paired fractions together:
their difference is $O_I(T/|\gamma_\chi|^2)$.
Dyadic summation and \eqref{proof:gb:local-zero-count} give a total
$O_I(\log(N(T+2)))$ from this range.
Here $d_\chi\le N$ makes the zero-count bound uniform in the
inducing conductor, so the implied constant depends only on $I$.

The remaining terms in \eqref{proof:gb:zero-identity} are smaller.
The series at $u=2$ is bounded absolutely, Stirling's formula bounds
the Gamma difference by $O_I(\log T)$, and
\[
 |e_\chi(u)|\le\sum_{p\mid N}\frac{\log p}{p^{1/2+\delta}-1}
       \le\frac{\log N}{2^{1/2+\delta}-1}
       \qquad(\Re u\ge1/2+\delta).
\]
The same estimate at $u=2$ gives $e_\chi(2)\ll\log N$.
This proves \eqref{proof:gb:logderivative-bound}.
This proves the lemma.
\end{proof}

\subsection{Mellin inversion and the signed families}
\label{proof:dir:character-mellin-summary}

We connect the unconditional finite projection to the contours used
in the model calculation and identify the three signed families.

For the companion weight $\Phi$ defined in \eqref{eq:fixed-goldbach-weight},
our Mellin transform and inversion conventions are
\[
 \widetilde\Phi(s)=\int_0^\infty\Phi(u)u^{s-1}\,\d u,
 \qquad
 \Phi(u)=\frac1{2\pi i}\int_{(\sigma)}
              \widetilde\Phi(s)u^{-s}\,\d s
 \quad(\sigma\in\mathbb R),
\]
as in~\eqref{gb:mellin-convention},
where $\int_{(\sigma)}$ is directed upwards along $\Re s=\sigma$.
Smoothness and compact support give, in particular,
\[
 \widetilde\Phi(\sigma+it)\ll_{\phi,I}(1+|t|)^{-6}
 \quad(\sigma\in I),
\]
for every fixed compact real interval $I$. The transform inserts
the weight into the logarithmic-derivative series, and its decay
controls the integrals when the contour is moved.

\begin{lemma}[Mellin projection and signed families]
\label{proof:v9:mellin-signed-families}
The following Mellin representation holds unconditionally for
$\sigma>1/2$ and, under $\mathrm{GRH}(N)$, also for
$0<\sigma\le1/9$:
\[
 \mathcal P(N)=\frac1{(2\pi i)^2}
 \int_{(\sigma)}\int_{(\sigma)}
 N^{s+z}\widetilde\Phi(s)\widetilde\Phi(z)\mathscr P_N(s,z)
 \,\d s\,\d z,
\]
as in~\eqref{gb:mellin-moment}.
The off-diagonal relations in \eqref{proof:dir:kernel-expansion} are
$m=n+kN$, $n=m+kN$, and $n+m=kN$, with signs $+,+,-$.
After insertion of the weight, only the sum family with $k=1$ remains.
\end{lemma}

\begin{proof}
For $\sigma>1/2$, the absolutely convergent logarithmic-derivative
series and the rapid decay of $\widetilde\Phi$ justify termwise
Mellin inversion:
\[
 P_\chi=-\frac1{2\pi i}\int_{(\sigma)}
       N^s\widetilde\Phi(s)\frac{L'}L(1/2+s,\chi)\,\d s,
\]
as in~\eqref{gb:mellin-P}.
The factor $N^s$ comes from $(n/N)^{-s}$. Apply the same formula
to $P_{\bar\chi}$, multiply, and average over $X_N^-$.
The two minus signs cancel, and absolute convergence permits the
product integrals and finite character sum to be interchanged.
This proves the displayed Mellin representation on the initial contours.
Under GRH, Proposition~\ref{proof:v3:kernel-holomorphy} excludes
poles between any positive contour and this initial contour.
Lemma~\ref{proof:lem:zero-expansion} gives logarithmic growth on
the intervening fixed strip, whereas $\widetilde\Phi$ decays faster
than every power. The horizontal sides of the contour rectangles
therefore tend to zero, and the vertical integrals converge absolutely.
Moving the contours successively proves the stated continuation.

Separating equal from unequal indices in the first congruence of
\eqref{proof:dir:kernel-expansion} gives the diagonal and the two
positive difference families. The second congruence gives the negative
sum family. On the support of the weight, $|n-m|\le N/3$ and
$2N/3\le n+m\le4N/3$. Thus both difference families and all sums
with $k\ge2$ have zero weighted contribution.
This proves the lemma.
\end{proof}

The signed combination of all three families must remain
in the analytic kernel. It supplies the cancellation used by the Euler model.

\section{Euler: the Beta model and the analytic residual}
\label{proof:sec:beta-mellin}\label{proof:v7:euler}

We construct a continuous model for the off-diagonal part of the
character kernel and subtract it, together with the exact diagonal,
to obtain an analytic residual. The model is specified by the
three arithmetic relations and their local densities. Its
construction does not assume an asymptotic formula for prime pairs.

Throughout this section, $N\ge4$ is even and $w=s+z$.
The character kernel is
\[
 \mathscr P_N(s,z)=\frac2{\varphi(N)}
 \sum_{\substack{\chi\ \mathrm{mod}\ N\\\chi(-1)=-1}}
 \frac{L'}L(1/2+s,\chi)\frac{L'}L(1/2+z,\bar\chi),
\]
as in~\eqref{eq:goldbach-moment},
where all characters, including imprimitive ones, are extended
by zero at nonunits. Its expansion
\eqref{proof:dir:kernel-expansion} has the diagonal $n=m$ and
the three remaining families
\[
 m=n+kN,\qquad n=m+kN,\qquad n+m=kN\qquad(k\ge1),
\]
with signs $+,+,-$, respectively. We keep the diagonal exactly and
assign to each remaining family its continuous integral multiplied
by the local density $\mathfrak S(kN)$. For any even $h$, this density is
\[
 \mathfrak S(h)=2C_2\prod_{\substack{p\mid h\\p>2}}\frac{p-1}{p-2},
 \qquad C_2=\prod_{p>2}\left(1-\frac1{(p-1)^2}\right),
\]
as in~\eqref{eq:singular-series}.
The factorizations below continue the resulting signed model to
the same domain as the arithmetic kernel.

\subsection{The signed Beta combination}
\label{proof:v7:signed-beta-model}

Fix $h=kN$ with $k\ge1$. For the relations $m=n+h$, $n=m+h$
and $n+m=h$, replace $(n,m)$ by $(x,x+h)$, $(x+h,x)$ and
$(x,h-x)$, respectively. Integrating the corresponding power weight
$n^{-1/2-s}m^{-1/2-z}$ over $x>0$ in the first two cases and
$0<x<h$ in the third gives the three integrals in
\eqref{v3:beta-scaled-model}. Scaling $x=ht$ produces the common
factor $h^{-s-z}$. The signs $+,+,-$ from the character projection
then give the Beta combination
\[
 K_-(s,z)=B(1/2-s,s+z)+B(1/2-z,s+z)
                         -B(1/2-s,1/2-z),
\]
as in~\eqref{eq:beta-kernel}.

Euler's Gamma function is
$\Gamma(s)=\int_0^\infty e^{-t}t^{s-1}\d t$ for $\Re s>0$,
with meromorphic continuation, and
$B(s,z)=\int_0^1t^{s-1}(1-t)^{z-1}\d t$ for $\Re s,\Re z>0$.
The change of variable $u=yt$ gives, for $y>0$,
\begin{equation}\label{v3i:Euler}
 \int_0^y u^{s-1}(y-u)^{z-1}\d u
 =B(s,z)y^{s+z-1},\quad
 B(s,z)=\frac{\Gamma(s)\Gamma(z)}{\Gamma(s+z)}
 \quad(\Re s,\Re z>0).
\end{equation}

\begin{lemma}[The continuous signed model]\label{proof:v9:signed-beta-model}
For $h>0$, $\Re s,\Re z<1/2$ and $\Re(s+z)>0$, put $w=s+z$.
The following integrals converge absolutely:
\begin{equation}\label{v3:beta-scaled-model}
 \begin{aligned}
 \int_0^\infty x^{-1/2-s}(x+h)^{-1/2-z}\,\mathrm{d}x
    &=h^{-w}B(1/2-s,w),\\
 \int_0^\infty x^{-1/2-z}(x+h)^{-1/2-s}\,\mathrm{d}x
    &=h^{-w}B(1/2-z,w),\\
 \int_0^h x^{-1/2-s}(h-x)^{-1/2-z}\,\mathrm{d}x
    &=h^{-w}B(1/2-s,1/2-z).
 \end{aligned}
\end{equation}
Their signed sum, with signs $+,+,-$, is $h^{-s-z}K_-(s,z)$.
For $h=kN$, the common local density for these three families is
$\mathfrak S(h)$, so their model is
$\mathfrak S(h)h^{-s-z}K_-(s,z)$.
\end{lemma}

\begin{proof}
For both sums and differences, the two entries exclude one
residue class modulo $p$ if $p\mid h$ and two otherwise.
Relative to independent prime densities, the correction is
\[
 \frac{1-\nu_p(h)/p}{(1-1/p)^2},\qquad
 \nu_p(h)=\begin{cases}1,&p\mid h,\\2,&p\nmid h.\end{cases}
\]
Since $h=kN$ is even, the product is
\[
 2C_2\prod_{\substack{p\mid h\\p>2}}\frac{p-1}{p-2}
 =\mathfrak S(h).
\]
In particular, $N\mid h$ makes the factors at primes dividing $N$
consistent with the coprimality restriction in the arithmetic kernel.
The substitutions $x=ht/(1-t)$ and $x=ht$ prove the three
identities in \eqref{v3:beta-scaled-model}. For example,
\[
 \int_0^\infty x^{-1/2-s}(x+h)^{-1/2-z}\,\mathrm{d}x
 =h^{-s-z}\int_0^1t^{-1/2-s}(1-t)^{s+z-1}\,\mathrm{d}t.
\]
The conditions $\Re s,\Re z<1/2$ control the finite endpoints,
and $\Re(s+z)>0$ controls infinity. In their common region,
the signed sum is therefore $h^{-s-z}K_-(s,z)$, with
$K_-$ as defined at the start of this subsection.
This proves the lemma.
\end{proof}

\subsection{Gamma factorization}
\label{proof:euler:gamma-section}

Reflection and duplication separate the signed Beta expression
into a factor in $w=s+z$ and one Gamma quotient in each variable.
Define the one-variable factor
\[
 \mathcal X_1(s)=\frac{\Gamma((3/2-s)/2)}{\Gamma((3/2+s)/2)}.
\]
The following lemma gives a bound and a Gamma factorization for $K_-(s,z)$.

\begin{lemma}[Gamma factorization]\label{proof:v9:gamma-factorization}
On $\Re s=\Re z=\sigma$, one has
\begin{equation}\label{proof:gb:beta-uniform-bound}
 |K_-(s,z)|\le
 2B(1/2-\sigma,2\sigma)+B(1/2-\sigma,1/2-\sigma)
 \qquad(0<\sigma\le1/9).
\end{equation}
For this Gamma quotient $\mathcal X_1$, the factorization is
\begin{equation}\label{proof:gb:gamma-beta}
 K_-(s,z)=2^{1-w}\Gamma(w)\cos(\pi w/2)
                      \mathcal X_1(s)\mathcal X_1(z),\qquad w=s+z.
\end{equation}
The latter equality is an identity of meromorphic functions.
\end{lemma}

\begin{proof}
Taking absolute values in the definition of $K_-(s,z)$ and applying
the triangle inequality proves \eqref{proof:gb:beta-uniform-bound}.
For the factorization, put $r=1/2-s$, $t=1/2-z$, so $r+t+w=1$.
Reflection gives
\[
 K_-(s,z)=\frac{\Gamma(r)\Gamma(t)\Gamma(w)}\pi
             (\sin\pi r+\sin\pi t-\sin\pi w),
\]
whereas duplication and reflection give
\[
 \mathcal X_1(s)=\frac{2^{1-r}}{\sqrt\pi}
                    \Gamma(r)\sin(\pi r/2).
\]
Use the corresponding expression for $z,t$ and the identity
\[
 \sin\pi r+\sin\pi t-\sin\pi w
 =4\sin(\pi r/2)\sin(\pi t/2)\cos(\pi w/2).
\]
This proves \eqref{proof:gb:gamma-beta} initially away from poles and
then as a meromorphic identity.
This proves the lemma.
\end{proof}

\subsection{The density series and its bounds}
\label{proof:v7:singular-series}\label{proof:v7:arithmetic-normalization}

We now collect the densities at all scales $h=kN$ and record the
exact diagonal. The corresponding Dirichlet series are
\[
 \begin{aligned}
 \mathscr D_N(w)&=\sum_{\substack{n\ge1\\(n,N)=1}}
                       \frac{\Lambda^2(n)}{n^{1+w}}
                       &&(\Re w>0),\\
 \mathscr S_N(w)&=\sum_{k\ge1}\frac{\mathfrak S(kN)}{k^w}
                       &&(\Re w>1),
 \end{aligned}
\]
as in~\eqref{eq:goldbach-arithmetic-kernels} and~\eqref{v3i:series}.
The former records the exact arithmetic diagonal; the latter collects
the local densities $\mathfrak S(kN)$. The density series is
initially defined for $\Re w>1$, whereas the three Beta integrals
have $0<\Re w<1$ in their common domain. To use the density factor
there, we first continue it by the Euler product, which we restate here:
\[
 \mathscr S_N(w)=\mathfrak S(N)\zeta(w)P_N(w),
\]
where $P_N$ is the Euler product
\[
 P_N(w)=\prod_{\substack{p>2\\p\nmid N}}
                  \left(1+\frac{p^{-w}}{p-2}\right),
\]
as in~\eqref{eq:goldbach-arithmetic-kernels}.
This Euler product is distinct from the character polynomial
$P_\chi$ of the preceding section.
Separating the zeta factor identifies the pole at $w=1$; controlling
the remaining product gives bounds uniform in the modulus.

\begin{lemma}\label{proof:lem:singular-N}
The diagonal series $\mathscr D_N(w)$ converges absolutely for
$\Re w>0$. The density series $\mathscr S_N(w)$ has the Euler
product displayed above and continues meromorphically to
$\Re w>0$, with a single simple pole at $w=1$ of residue
\begin{equation}\label{proof:gb:singular-series-residue}
 \operatorname*{Res}_{w=1}\mathscr S_N(w)=\frac N{\varphi(N)}.
\end{equation}
For fixed $0<\sigma_-<\sigma_+<\infty$,
\begin{equation}\label{proof:gb:singular-series-growth}
 |\mathscr S_N(w)|\ll_{\sigma_-,\sigma_+}
 \frac N{\varphi(N)}(1+|\Im w|)
 \quad(\sigma_-\le\Re w\le\sigma_+,\ |\Im w|\ge1).
\end{equation}
Moreover, $\mathfrak S(N)>1$ and
\begin{equation}\label{proof:sp:family-bounds}
 C_2\frac N{\varphi(N)}\le\mathfrak S(N)
       \le\frac N{\varphi(N)}\ll\log\log N.
\end{equation}
\end{lemma}

\begin{proof}
Absolute convergence of the diagonal follows from
$\Lambda^2(n)\le\log^2n$. The function
$\mathfrak S(kN)/\mathfrak S(N)$ is multiplicative in $k$.
At $p^j$, $j\ge1$, its value is $1$ if $p\mid N$ and
$(p-1)/(p-2)$ otherwise. Thus the local factor is
$(1-p^{-w})^{-1}$ in the first case and
$(1+p^{-w}/(p-2))/(1-p^{-w})$ in the second. This proves the
Euler product for $\Re w>1$.

For every $\delta>0$, the product $P_N(w)$ is holomorphic and
uniformly bounded in $N$ on $\Re w\ge\delta$, since
$\sum_{p>2}p^{-\delta}/(p-2)<\infty$.
Omitting the primes dividing $N$ only decreases this absolute
majorant. At $w=1$ all Euler factors are positive, so they do not
cancel the pole of the zeta factor.
Only the pole of $\zeta(w)$ at $1$ remains. The identity
\[
 \left(1-\frac1{(p-1)^2}\right)
 \left(1+\frac1{p(p-2)}\right)=1
\]
gives \eqref{proof:gb:singular-series-residue} and
\[
 \mathfrak S(N)=\frac N{\varphi(N)}
 \prod_{\substack{p>2\\p\nmid N}}
             \left(1-\frac1{(p-1)^2}\right).
\]
The remaining product lies between $C_2$ and $1$.
Also $C_2>\prod_{m=2}^\infty(1-m^{-2})=1/2$, so
$\mathfrak S(N)\ge2C_2>1$.

Euler summation yields
\[
 \zeta(w)=\frac w{w-1}
       -w\int_1^\infty\{x\}x^{-w-1}\,\mathrm{d}x
       \qquad(\Re w>0,\ w\ne1).
\]
Together with the product bound this proves
\eqref{proof:gb:singular-series-growth}.
Finally, split the prime divisors of $N$ at $y=\log N$.
Mertens' product estimate bounds the contribution of $p\le y$
by $O(\log y)$. There are at most $\log N/\log y$ prime divisors above $y$,
and each contributes $O(1/y)$ to the logarithm of the product. In particular,
\[
 \log\prod_{\substack{p\mid N\\p>y}}(1-1/p)^{-1}
 \ll\frac{\log N}{y\log y}=\frac1{\log y}.
\]
This proves the last bound in \eqref{proof:sp:family-bounds}; the bounded
initial range is absorbed into the constant. The use of Mertens'
estimate is the standard one in \cite{IK2004}.
This proves the lemma.
\end{proof}

\subsection{Continuation of the model}
\label{proof:v7:model-continuation}

We now combine the two factorizations. This removes the apparent
pole at $s+z=1$ and gives the domain for the residual operator.
With the continued density factor, define the off-diagonal model by
\[
 \mathscr K_N(s,z)=N^{-w}\mathscr S_N(w)K_-(s,z),\qquad w=s+z,
\]
as in~\eqref{eq:goldbach-arithmetic-kernels}.
The equivalent form to be proved is
\[
 \mathscr K_N(s,z)=\mathfrak S(N)\left(\frac\pi N\right)^w
       P_N(w)\zeta(1-w)\mathcal X_1(s)\mathcal X_1(z),
\]
as in~\eqref{v3:model-functional-form}.
The second expression makes holomorphy beyond the original Beta
integral domain visible factor by factor.

\begin{proposition}[Continuation of the Beta model]\label{proof:v3:model-geometry}
The second expression for $\mathscr K_N$ above gives its continuation
and is holomorphic on
\[
 \Omega=\{(s,z):0<\Re s<3/2,\ 0<\Re z<3/2\}.
\]
\end{proposition}

\begin{proof}
Insert \eqref{proof:gb:gamma-beta} and the density Euler product into
$\mathscr K_N(s,z)=N^{-w}\mathscr S_N(w)K_-(s,z)$.
The zeta functional equation
\[
 \zeta(1-w)=2(2\pi)^{-w}\Gamma(w)\cos(\pi w/2)\zeta(w)
\]
gives the asserted functional form of the model.
The density identity has already been continued to the initial
Beta domain. Thus the substitution uses established meromorphic
identities, rather than interchanging the density series with an
integral outside its region of absolute convergence.
On $\Omega$, $P_N(s+z)$ converges locally uniformly and neither
Gamma quotient has a pole: each numerator has argument with
positive real part, and the reciprocal Gamma function is entire.
The pole of $\zeta(1-s-z)$ would require $s+z=0$, outside
$\Omega$. Hence the model is holomorphic there.
This proves the proposition.
\end{proof}

\subsection{Holomorphy and symmetry of the residual}

The residual is the exact difference between the character kernel
and its diagonal and off-diagonal model:

\[
 \mathscr E_N(s,z)=\mathscr P_N(s,z)
                  -\mathscr D_N(s+z)-\mathscr K_N(s,z).
\]
Here $\mathscr P_N$ is the character kernel at the start of this
section. Its exact diagonal is
\[
 \mathscr D_N(w)=\sum_{(n,N)=1}\Lambda^2(n)n^{-1-w},
\]
as in~\eqref{eq:goldbach-arithmetic-kernels},
and $\mathscr K_N$ is the continued model.
We now assume $\mathrm{GRH}(N)$: every nontrivial zero of every
primitive Dirichlet $L$-function of conductor dividing $N$ lies on
$\Re s=1/2$. Proposition~\ref{proof:v3:kernel-holomorphy} then
gives holomorphy of $\mathscr P_N$ for $\Re s,\Re z>0$;
the other two terms have just been treated unconditionally.
Their common domain is
$\Omega=\{(s,z):0<\Re s<3/2,\ 0<\Re z<3/2\}$.
As with the character kernel, $\mathscr E_N(s,z)$ is the form
holomorphic in both variables, while $\mathscr E_N(s,\bar z)$
is the Hermitian form used in the kernel of the integral operator
$T_N$ defined in~\eqref{qp:operator}.

\begin{proposition}[The analytic residual]\label{proof:v7:residual-holomorphy}
Under $\mathrm{GRH}(N)$, $\mathscr E_N$ is holomorphic on
$\Omega$ and satisfies
\[
 \mathscr E_N(s,z)=\mathscr E_N(z,s),\qquad
 \overline{\mathscr E_N(s,z)}=\mathscr E_N(\bar s,\bar z).
\]
In particular, $\mathscr E_N(s,\bar z)$ is Hermitian.
\end{proposition}

\begin{proof}
Subtract the diagonal and model from $\mathscr P_N$.
Proposition~\ref{proof:v3:kernel-holomorphy} and
Proposition~\ref{proof:v3:model-geometry}, together with
Lemma~\ref{proof:lem:singular-N}, prove
holomorphy. The diagonal has real coefficients and depends only on
$s+z$; the Beta model is symmetric in $s,z$ and commutes with
simultaneous conjugation. Equation~\eqref{proof:v3:kernel-symmetries}
therefore gives both residual identities.
The subtraction is exact, so these analytic conclusions require
no estimate for the size or sign of the residual.
This proves the proposition.
\end{proof}

The next section realizes this residual as a compact selfadjoint
operator. To state the arithmetic identity that this construction
will control, recall the fixed nonzero weight
$\phi\in C_c^\infty(\mathbb R)$, with $0\le\phi\le1$,
support in $[1/3,2/3]$, and $\phi(u)=\phi(1-u)$. Put
\[
 \Phi(u)=\sqrt u\,\phi(u),\qquad
 \widetilde\Phi(s)=\int_0^\infty\Phi(u)u^{s-1}\,\d u,
 \qquad c_\phi=\int_0^1\phi^2(u)\,\d u,
\]
as in~\eqref{eq:fixed-goldbach-weight}, \eqref{gb:mellin-convention}, and~\eqref{eq:c0}.
For $0<\sigma\le1/9$, define
\[
 r_\sigma(N)=\frac1{(2\pi i)^2}
 \int_{(\sigma)}\int_{(\sigma)}
 N^{s+z}\widetilde\Phi(s)\widetilde\Phi(z)
 \mathscr E_N(s,z)\,\d s\,\d z,
\]
as in~\eqref{v3w:residual},
where both contours are directed upwards.
Section~\ref{proof:sec:v3-mellin-operator} will evaluate the
model integral as $-c_\phi\mathfrak S(N)$ and the diagonal
integral as $\mathcal D(N)$. Combined with the finite identity
$\mathcal G(N)=\mathcal D(N)-\mathcal P(N)$, these evaluations give
\[
 \mathcal G(N)=c_\phi\mathfrak S(N)-r_\sigma(N),
\]
as in~\eqref{v3w:exact},
where $\mathcal G(N)$ is the weighted Goldbach correlation
\[
 \mathcal G(N)=\frac1N
 \sum_{\substack{n+m=N\\(nm,N)=1}}
 \Lambda(n)\Lambda(m)\phi(n/N)\phi(m/N),
\]
as in~\eqref{eq:goldbach-correlation}.
Thus the signed discrepancy $r_\sigma(N)$ is independent of
$\sigma$. The remaining task is to express it as a quadratic
form and bound it by the norm of the residual operator.

\section{Hilbert--P\'olya: a compact selfadjoint realization}
\label{proof:sec:qp-foundation}

We realize the arithmetic residual as an integral operator on a
space of holomorphic functions. The main analytic point is that its
boundary integral produces the two traces of a single holomorphic
function on the strip. Once this is established, the Hermitian
symmetry and square integrability of the kernel give a compact
selfadjoint operator. Point evaluation will then recover the
residual from this operator.

Throughout, $N>e^{10}$ is even, and we assume $\mathrm{GRH}(N)$:
every nontrivial zero of every primitive Dirichlet $L$-function
whose conductor divides $N$ has real part $1/2$.
The construction concerns fixed $N$ unless uniformity is stated.
We use the strip and boundary conventions
\[
 \begin{gathered}
 \alpha_N=\frac1{\log N},\qquad \beta=\frac18,\qquad
 \sigma_*=\frac19,\qquad \Delta_N=\beta-\alpha_N,\\
 S_N=\{s:\alpha_N<\Re s<\beta\},\qquad
 \mathcal H_N=H^2(S_N),\\
 \Gamma_N=(\alpha_N+i\mathbb R)\,\sqcup\,(\beta+i\mathbb R),\qquad
 \langle h,g\rangle_{\mathcal H_N}
       =\int_{\Gamma_N}h(s)\overline{g(s)}\,|\mathrm{d}s|,
 \end{gathered}
\]
as in~\eqref{hd:strip}, \eqref{hd:space-and-boundary}, and~\eqref{hd:inner-product}.
Here $H^2(S_N)$ is the strip Hardy space, whose definition and
boundary traces are given below. The inner product is linear in
its first variable. Both lines carry positive arc length $\mathrm{d}t$,
with no contour orientation and no factor $1/(2\pi)$.
The choice of endpoints places $\sigma_*$ strictly inside $S_N$
and gives $N^{2\alpha_N}=e^2$; this last identity will be used to
bound the Mellin functional uniformly in $N$.

For reference, we collect the three arithmetic kernels entering the operator.
Let $X_N^- = \{\chi\ \mathrm{mod}\ N:\chi(-1)=-1\}$, including imprimitive
characters and extending each character by zero at nonunits. The kernels are
\[
 \begin{aligned}
 \mathscr P_N(s,z)&=\frac2{\varphi(N)}
              \sum_{\chi\in X_N^-}g_\chi(s)g_{\bar\chi}(z),\\
 \mathscr D_N(w)&=\sum_{\substack{n\ge1\\(n,N)=1}}
                     \frac{\Lambda^2(n)}{n^{1+w}},\\
 \mathscr K_N(s,z)&=\mathfrak S(N)\left(\frac\pi N\right)^{s+z}
 P_N(s+z)\zeta(1-s-z)\mathcal X_1(s)\mathcal X_1(z),
 \end{aligned}
\]
as in~\eqref{eq:goldbach-moment}, \eqref{eq:goldbach-arithmetic-kernels}, and~\eqref{v3:model-functional-form},
where $g_\chi$ is the logarithmic derivative
\[
 g_\chi(s)=\frac{L'}L(1/2+s,\chi),
\]
as in~\eqref{qp:Gram}.
The kernel $\mathscr K_N$ is the continued Beta model. Its factors are
\[
 \begin{gathered}
 \mathfrak S(N)=2C_2\prod_{\substack{p\mid N\\p>2}}\frac{p-1}{p-2},
 \qquad C_2=\prod_{p>2}\left(1-\frac1{(p-1)^2}\right),\\
 P_N(w)=\prod_{\substack{p>2\\p\nmid N}}
          \left(1+\frac{p^{-w}}{p-2}\right),\qquad
 \mathcal X_1(s)=\frac{\Gamma((3/2-s)/2)}{\Gamma((3/2+s)/2)}.
 \end{gathered}
\]
The singular-series and Euler-product definitions are as in
\eqref{eq:singular-series} and~\eqref{eq:goldbach-arithmetic-kernels},
respectively. Here $\Lambda$ is the von Mangoldt function and $\varphi(N)$ is
Euler's totient. Under $\mathrm{GRH}(N)$, the preceding section
shows that
\[
 \mathscr E_N(s,z)=\mathscr P_N(s,z)
                  -\mathscr D_N(s+z)-\mathscr K_N(s,z)
\]
is holomorphic on $0<\Re s,\Re z<3/2$. The integral operator to
be justified is
\[
 (T_Nh)(s)=\int_{\Gamma_N}
       \frac{\mathscr E_N(s,\bar z)}{(1+s)(1+\bar z)}
       h(z)\,|\mathrm{d}z|,
\]
as in~\eqref{qp:operator}.
The denominators provide decay on the infinite boundary lines;
conjugation in the second variable supplies the Hermitian symmetry.

\subsection{The Hermitian kernel on the arithmetic strip}
\label{proof:v9:hermitian-kernel-section}

We apply the holomorphy and symmetry of $\mathscr E_N$ on the
strip $S_N$ just defined. Its two holomorphic variables become the
first variable and the conjugate of the second in the integral
kernel: $\mathscr E_N(s,\bar z)$ is holomorphic in $s$ and
antiholomorphic in $z$. The Hermitian condition is
$\mathscr E_N(s,\bar z)=\overline{\mathscr E_N(z,\bar s)}$.

The additional term \emph{Fueter regular} refers here to a
quaternionic first-order differential equation. Write
$\omega=s+zj$, with $s,z\in\mathbb C$ and $j^2=-1$, $ij=-ji$.
In these coordinates the equation we use is
$(\partial_{\bar s}+j\partial_z)F=0$. For the present kernel,
holomorphy in $s$ and antiholomorphy in $z$ make the two
derivatives vanish separately.

\begin{proposition}[The Hermitian arithmetic kernel]
\label{proof:v9:hermitian-kernel}
The weighted kernel of $T_N$ is
\[
 \frac{\mathscr E_N(s,\bar z)}{(1+s)(1+\bar z)},
\]
as in~\eqref{qp:operator}.
On the arithmetic strip, both this kernel and $\mathscr E_N(s,\bar z)$
are Hermitian and holomorphic in $s$ and $\bar z$. The weighted kernel
extends to a neighborhood of the closed strip in each holomorphic
coordinate.
In the coordinates $\omega=s+zj$, the residual satisfies
\[
 (\partial_{\bar s}+j\partial_z)\mathscr E_N(s,\bar z)=0.
\]
The fixed contour has distances $\sigma_*-\alpha_N>1/90$ and
$\beta-\sigma_*=1/72$ from the two strip boundaries.
\end{proposition}

\begin{proof}
Proposition~\ref{proof:v7:residual-holomorphy} gives holomorphy
on $0<\Re s,\Re z<3/2$ and the identities
\[
 \mathscr E_N(s,z)=\mathscr E_N(z,s),\qquad
 \overline{\mathscr E_N(s,z)}=\mathscr E_N(\bar s,\bar z).
\]
They imply $\mathscr E_N(s,\bar z)
=\overline{\mathscr E_N(z,\bar s)}$, the Hermitian symmetry needed here.
The denominator $(1+s)(1+\bar z)$ has no zeros near the closed
strip and preserves these properties. Holomorphy in $s$ and
$\bar z$ makes both derivatives in the displayed equation zero,
proving the stated Fueter regularity.
Finally, $\alpha_N<1/10$, $\sigma_*=1/9$ and $\beta=1/8$
give the two distances.
This proves the proposition.
\end{proof}

Here the arithmetic coordinates are the two complex Mellin variables
$s$ and $z$, viewed as the four real coordinates
$(\Re s,\Im s,\Re z,\Im z)$ of $\omega=s+zj$. For fixed $N$,
the Fueter equation for $\mathscr E_N(s,\bar z)$ holds because its
holomorphy in $s$ and antiholomorphy in $z$ make
$\partial_{\bar s}\mathscr E_N(s,\bar z)$ and
$\partial_z\mathscr E_N(s,\bar z)$ vanish separately.
The spectral variable $q$ in $\mathfrak X_N(q)$ is introduced
independently of $(s,z)$. The spectral argument below uses the
selfadjoint operator $T_N$ on a complex Hilbert space and its real
eigenvalues; no positivity is asserted for the residual.

\subsection{Hardy boundary traces}

In the usual disk Hardy space $H^2$, one controls the $L^2$
means on concentric circles and obtains an $L^2$ value on the
boundary circle. A strip has two boundary components, so the
corresponding norm here is the sum of two squared $L^2$ norms.
There is also a compatibility condition: the two boundary
functions must come from the same holomorphic function in the
interior. Thus the boundary realization occupies a subspace of
$L^2(\mathbb R)\oplus L^2(\mathbb R)$.

Precisely, for $S=\{s:\alpha<\Re s<\beta\}$, define $H^2(S)$
to consist of the holomorphic functions $h$ on $S$ for which
\begin{equation}\label{hd:interior-norm}
 \|h\|_{H^2(S)}^2
 =\sup_{0<\varepsilon<(\beta-\alpha)/2}
 \int_{\mathbb R}\left(
 |h(\alpha+\varepsilon+it)|^2+
 |h(\beta-\varepsilon+it)|^2\right)\,\mathrm{d}t<\infty.
\end{equation}
The boundary values below are understood as $L^2$ limits:
\[
 h_\alpha=\lim_{x\downarrow\alpha}h(x+i\,\cdot),\qquad
 h_\beta=\lim_{x\uparrow\beta}h(x+i\,\cdot).
\]
Their existence is part of the next lemma, so no pointwise
boundary smoothness is required in the definition.

The Fourier--Laplace representation makes the compatibility of
the two boundaries explicit. Set
$w_{H,\alpha,\beta}(\tau)=e^{2\alpha\tau}+e^{2\beta\tau}$.
For $\psi\in L^2(\mathbb R,w_{H,\alpha,\beta}(\tau)\,\mathrm{d}\tau)$,
define
\begin{equation}\label{hd:PW-map}
 (\mathcal U_H\psi)(s)=\frac1{\sqrt{2\pi}}
       \int_{\mathbb R}\psi(\tau)e^{s\tau}\,\mathrm{d}\tau.
\end{equation}
On the vertical line $s=x+it$, this is the inverse Fourier
transform of $e^{x\tau}\psi(\tau)$. The same density $\psi$
therefore describes both boundaries, with factors $e^{\alpha\tau}$
and $e^{\beta\tau}$. Plancherel accounts for the two terms in
$w_{H,\alpha,\beta}$ and gives the boundary norm identity
\begin{equation}\label{hd:boundary-norm}
 \|h\|_{H^2(S)}^2
 =\|h_\alpha\|_{L^2(\mathbb R)}^2
   +\|h_\beta\|_{L^2(\mathbb R)}^2
 =\int_{\mathbb R}|\psi(\tau)|^2
                     w_{H,\alpha,\beta}(\tau)\,\mathrm{d}\tau.
\end{equation}
This fixes our normalization: each boundary contributes its
ordinary $L^2(\mathbb R,\mathrm{d}t)$ norm, without an averaging factor.
On the arithmetic strip, write
$w_{H,N}(\tau)=w_{H,\alpha_N,\beta}(\tau)$.
The general endpoints will also be used in Appendix~\ref{proof:app:param}.

\begin{lemma}[Fourier--Laplace representation]\label{proof:hd:PW}\label{proof:bm:PW}
For a strip $S$ with real endpoints $\alpha<\beta$, equip $H^2(S)$
with \eqref{hd:interior-norm}.
The map \eqref{hd:PW-map} is unitary from
$L^2(\mathbb R,w_{H,\alpha,\beta}(\tau)\,\mathrm{d}\tau)$ onto this Hardy space.
Every function has $L^2$ traces on both boundary lines, and
\eqref{hd:boundary-norm} holds. The trace map identifies the Hardy
space with a closed subspace of boundary $L^2$; on the arithmetic
strip its inner product is
\begin{samepage}
\[
 \langle h,g\rangle_{\mathcal H_N}
       =\int_{\Gamma_N}h(s)\overline{g(s)}\,|\mathrm{d}s|,
\]
as in~\eqref{hd:inner-product}.
\end{samepage}
\end{lemma}

\begin{proof}
For an interior $x$, the quotient
$e^{2x\tau}/w_{H,\alpha,\beta}(\tau)$ decays exponentially at
both ends. Cauchy--Schwarz therefore justifies the integral in
\eqref{hd:PW-map} and its derivatives locally uniformly.
On a compact substrip, the same exponential domination remains
integrable after multiplication by any fixed power of $\tau$.
This justifies differentiation under the integral. It defines a holomorphic function $h$, and Plancherel gives
\[
 \int_{\mathbb R}|h(x+it)|^2\,\mathrm{d}t
   =\int_{\mathbb R}|\psi(\tau)|^2e^{2x\tau}\,\mathrm{d}\tau.
\]
As $x$ approaches an endpoint $x_0$, dominated convergence gives
$e^{x\tau}\psi\to e^{x_0\tau}\psi$ in $L^2$; the squared
difference is bounded by $4|\psi|^2w_{H,\alpha,\beta}$.
Plancherel then gives the boundary traces.

Convexity yields
\[
 e^{2(\alpha+\varepsilon)\tau}+e^{2(\beta-\varepsilon)\tau}
 \le e^{2\alpha\tau}+e^{2\beta\tau}
 \qquad(0<\varepsilon<(\beta-\alpha)/2).
\]
Taking $\varepsilon\downarrow0$ proves that the paired supremum
equals the boundary norm, and hence that the map is isometric.

Conversely, for a Hardy function $h$, take the $L^2$ Fourier
transform $\widehat h_x$ of $h(x+i\,\cdot)$ with the convention
\[
 \widehat h_x(\tau)=(2\pi)^{-1/2}
       \int h(x+it)e^{-it\tau}\,\mathrm{d}t.
\]
The Cauchy--Riemann equation gives
$\partial_x\widehat h_x=\tau\widehat h_x$ in distributions.
Thus $e^{-x\tau}\widehat h_x(\tau)$ is independent of $x$
on compact frequency intervals. These representations agree on
overlaps and define a measurable function $\psi$.
Plancherel and Fatou, along a sequence of paired lines tending
to the boundaries, give
\[
 \int_{\mathbb R}|\psi(\tau)|^2w_{H,\alpha,\beta}(\tau)\,\mathrm{d}\tau
 \le\|h\|_{H^2(S)}^2.
\]
The forward construction therefore applies. Fourier inversion
identifies its image with $h$ almost everywhere, and holomorphy
gives equality everywhere. This proves surjectivity. The boundary
norm identity makes the trace map an isometry from a complete Hilbert
space, so its image in boundary $L^2$ is closed. Polarization gives
the stated boundary inner product.
This proves the lemma.
\end{proof}

\subsection{The compact selfadjoint operator}

For an integral operator on an $L^2$ space, a square-integrable
kernel gives a Hilbert--Schmidt operator: its squared
Hilbert--Schmidt norm is the integral of the squared absolute
value of the kernel. This is the integral-kernel form of
$\sum_{j,k}|a_{jk}|^2$ for a matrix. Here the boundary has two
components, so that integral has four contributions, according
to the choices of the first and second boundary lines.

For fixed $N$, recall that $\mathcal H_N=H^2(S_N)$ consists of
holomorphic functions on $S_N=\{s:\alpha_N<\Re s<\beta\}$ whose
paired vertical $L^2$ means remain bounded as the two lines approach
the boundary. By Lemma~\ref{proof:hd:PW}, we identify this Hardy
space isometrically with its closed trace subspace in $L^2(\Gamma_N)$.
To construct $T_N$ on $\mathcal H_N$, we will show that the boundary
integral operator $\widetilde T_N$ associated with the weighted
residual maps $L^2(\Gamma_N)$ into this trace subspace.
Hermitian symmetry will then imply that $\widetilde T_N$ vanishes
on $\mathcal H_N^\perp$. Its restriction $T_N$ to $\mathcal H_N$
therefore has the same Hilbert--Schmidt norm as $\widetilde T_N$,
which gives the equality in the following proposition.

\begin{proposition}[The compact selfadjoint realization]\label{proof:bd:HS}
Consider the operator on $\mathcal H_N$ defined by
\[
 (T_Nh)(s)=\int_{\Gamma_N}
       \frac{\mathscr E_N(s,\bar z)}{(1+s)(1+\bar z)}h(z)\,|\mathrm{d}z|.
\]
Its weighted residual kernel is Hermitian, holomorphic near the
closed strip in each holomorphic coordinate, and belongs to
$L^2(\Gamma_N\times\Gamma_N)$ on the boundary.
The operator $T_N$ is compact, selfadjoint, and Hilbert--Schmidt. Writing $\mu_j(N)$ for its nonzero eigenvalues,
counted with multiplicity, one has
\begin{equation}\label{hd:boundary-HS-norm}
 \sum_j\mu_j^2(N)=\|T_N\|_{\mathrm{HS}}^2
 =\int_{\Gamma_N}\int_{\Gamma_N}|\mathscr E_N(s,\bar z)|^2
      \frac{|\mathrm{d}s|}{|1+s|^2}\frac{|\mathrm{d}z|}{|1+\bar z|^2}.
\end{equation}
Its nonzero eigenvalues are real and satisfy $\|T_N\|=\sup_j|\mu_j(N)|$, with the maximum attained
unless $T_N=0$.
\end{proposition}

\begin{proof}
Holomorphy and symmetry follow from
Proposition~\ref{proof:v7:residual-holomorphy}, since the closed strip
lies inside its domain. For integrability, constants may depend
on $N$. For the arithmetic component, conjugation of the second
variable gives the Gram identity
\[
 \mathscr P_N(s,\bar z)=\frac2{\varphi(N)}
       \sum_{\chi\in X_N^-}g_\chi(s)\overline{g_\chi(z)},
 \qquad g_\chi(s)=\frac{L'}L(1/2+s,\chi),
\]
as in~\eqref{qp:Gram}.
Lemma~\ref{proof:lem:zero-expansion} bounds these logarithmic
derivatives on the fixed strip by
\[
 |g_\chi(x+iv)|\ll_N\log^2(2+|v|)
       \qquad(\alpha_N\le x\le\beta).
\]
The nonnegative coefficients of the diagonal series give
$|\mathscr D_N(s+\bar z)|\le\mathscr D_N(2\alpha_N)$ on the
closed strip. The continued model has the factorization
\[
 \mathscr K_N(s,\bar z)=\mathfrak S(N)
 \left(\frac\pi N\right)^{s+\bar z}P_N(s+\bar z)
 \zeta(1-s-\bar z)\mathcal X_1(s)\mathcal X_1(\bar z).
\]
The Euler product $P_N$ and the Gamma quotients $\mathcal X_1$
are bounded on the strip for fixed $N$, leaving the zeta factor
to control. For
$3/4\le\Re s\le1-2\alpha_N$, Euler summation at
$X=2+|\Im s|$ gives
$|\zeta(s)|\ll_N(1+|\Im s|)^{1/4}$: the finite sum and
integral term are $O_N(X^{1-\Re s})$, and the remainder is
$O(|s|X^{-\Re s}/\Re s)$.
Using $1+|v-t|\le(1+|v|)(1+|t|)$ and absorbing logarithms
into these powers, we obtain
\begin{equation}\label{proof:hd:kernel-majorant}
 \left|\frac{\mathscr E_N(x+iv,y-it)}{(1+x+iv)(1+y-it)}\right|^2
 \ll_N(1+|v|)^{-3/2}(1+|t|)^{-3/2}
 \quad(\alpha_N\le x,y\le\beta).
\end{equation}
This majorant is integrable on all four boundary products and
uniform across the closed strip.

The boundary kernel defines a selfadjoint Hilbert--Schmidt
operator $\widetilde T_N$ on $L^2(\Gamma_N)$.
For $h$ in that space, Cauchy--Schwarz and
\eqref{proof:hd:kernel-majorant} give uniform vertical $L^2$ bounds
for the integral
\[
 \int_{\Gamma_N}
       \frac{\mathscr E_N(s,\bar z)}{(1+s)(1+\bar z)}h(z)\,|\mathrm{d}z|.
\]
Local domination and Morera's theorem give holomorphy. Dominated convergence of the
squared kernel differences gives its $L^2$ boundary limits.
Lemma~\ref{proof:hd:PW} therefore places its range in the Hardy trace space.

Selfadjointness and the range inclusion imply that $\widetilde T_N$
annihilates the orthogonal complement of the trace space.
Indeed, pairing the image of an orthogonal vector with any
boundary $L^2$ vector gives zero: selfadjointness transfers the
operator to the latter vector, whose image lies in the trace space. Hence on
$L^2(\Gamma_N)=\mathcal H_N\oplus\mathcal H_N^\perp$ it is
$T_N\oplus0$. The Hilbert--Schmidt kernel identity and the
compact selfadjoint spectral theorem now give
\eqref{hd:boundary-HS-norm} for $T_N$ itself, as well as the
reality of the eigenvalues and the asserted operator norm formula.
This proves the proposition.
\end{proof}

\subsection{Point evaluation and the arithmetic residual}

An arbitrary $L^2$ function has no well-defined value at a
specified point. In this Hardy space, holomorphy makes interior
point evaluation a bounded functional. Its Riesz vector $k_z$
satisfies $h(z)=\langle h,k_z\rangle$ and is called the
\emph{reproducing kernel} at $z$. For the disk with normalized
boundary measure, it is $(1-s\bar z)^{-1}$; our strip gives
the cosecant formula below, with a constant determined by our
boundary measure. Applying these evaluation vectors in both
variables will recover the arithmetic kernel as a matrix element
of $T_N$.

\begin{proposition}[Point evaluation]\label{proof:bd:evaluation}
For $z\in S_N$, define
\begin{equation}\label{hd:reproducing-kernel}
 k_z(s)=\frac1{4\Delta_N}
       \csc\!\left(\frac{\pi(s+\bar z-2\alpha_N)}{2\Delta_N}\right).
\end{equation}
Then $k_z\in\mathcal H_N$ and
$h(s)=\langle h,k_s\rangle_{\mathcal H_N}$.
Writing $d_N(s)=\min(\Re s-\alpha_N,\beta-\Re s)$,
one has
\begin{equation}\label{proof:hd:evaluation-norm}
 \|k_s\|_{\mathcal H_N}^2
 =\frac1{4\Delta_N}
       \csc\!\left(\frac{\pi(\Re s-\alpha_N)}{\Delta_N}\right)
 \le\frac1{8d_N(s)}.
\end{equation}
Consequently,
\begin{equation}\label{proof:bd:pointwise}
 |\mathscr E_N(s,\bar z)|\le
 \frac{\|T_N\|\,|1+s|\,|1+z|}
                       {8\sqrt{d_N(s)d_N(z)}}.
\end{equation}
On the fixed interior contour this gives
\begin{equation}\label{proof:bd:pointwise-fixed}
 |\mathscr E_N(\sigma_*+iv,\sigma_*+it)|
 \ll\|T_N\|(1+|v|+|t|)^2.
\end{equation}
\end{proposition}

\begin{proof}
By Lemma~\ref{proof:hd:PW}, the evaluation vector at $z$ has coefficient
$(2\pi)^{-1/2}e^{\bar z\tau}/w_{H,N}(\tau)$ in the weighted
Fourier space. The conjugate of $z$ in this coefficient is forced by the
convention that the inner product is linear in its first variable. Its inverse transform is
\[
 k_z(s)=\frac1{2\pi}\int_{\mathbb R}
       \frac{e^{(s+\bar z)\tau}}
                  {e^{2\alpha_N\tau}+e^{2\beta\tau}}\,\mathrm{d}\tau.
\]
With $t=e^{2\Delta_N\tau}$ and
$\lambda=(s+\bar z-2\alpha_N)/(2\Delta_N)$, this becomes
\[
 \frac1{4\pi\Delta_N}\int_0^\infty
             \frac{t^{\lambda-1}}{1+t}\,\mathrm{d}t
 =\frac{B(\lambda,1-\lambda)}{4\pi\Delta_N},
 \qquad 0<\Re\lambda<1.
\]
Euler's reflection identity proves \eqref{hd:reproducing-kernel}.
Evaluation at $s=z$ and
$\sin(\pi t)\ge2\min(t,1-t)$, $0\le t\le1$, give
\eqref{proof:hd:evaluation-norm}.

Applying reproduction in both variables yields
\[
 \frac{\mathscr E_N(s,\bar z)}{(1+s)(1+\bar z)}
      =\langle T_Nk_z,k_s\rangle_{\mathcal H_N}.
\]
Cauchy--Schwarz proves \eqref{proof:bd:pointwise}. The distances of
$\sigma_*$ from both boundaries are at least $1/90$, giving
\eqref{proof:bd:pointwise-fixed} after conjugating the second argument.
This proves the proposition.
\end{proof}

\section{Riemann: spectral zeros and operator norm bounds}
\label{proof:sec:spectral-analysis}\label{proof:v7:riemann}

The residual operator converts the Goldbach error into a spectral
question: how large can its eigenvalues be? We encode their magnitudes
as zeros of a product. Its quadratic factors will make the relation
between zero location and operator norm explicit. We then establish
the norm bound available under GRH and examine the scale of the left
Hardy boundary. The resulting spectral estimates will be used in
the next section to prove the main theorems.

Throughout, $N>e^{10}$ is even and satisfies $\mathrm{GRH}(N)$.
We recall the Hardy space and operator
\[
 \begin{gathered}
 \mathcal H_N=H^2(S_N),\qquad
 S_N=\{s:\alpha_N<\Re s<\beta\},\qquad
 \alpha_N=\frac1{\log N},\quad\beta=\frac18,\\
 (T_Nh)(s)=\int_{\Gamma_N}
 \frac{\mathscr E_N(s,\bar z)}{(1+s)(1+\bar z)}h(z)\,|\mathrm{d}z|,
 \end{gathered}
\]
as in~\eqref{hd:strip}, \eqref{hd:space-and-boundary}, and~\eqref{qp:operator}.
Here
$\Gamma_N=(\alpha_N+i\mathbb R)\sqcup(\beta+i\mathbb R)$,
both lines carrying positive arc length. The elements of $\mathcal H_N$
are holomorphic functions whose paired interior $L^2$ means remain
bounded. Their traces on the two boundaries give the norm
\[
 \|h\|_{\mathcal H_N}^2
 =\int_{\mathbb R}|h(\alpha_N+it)|^2\,\mathrm{d}t
  +\int_{\mathbb R}|h(\beta+it)|^2\,\mathrm{d}t.
\]
The two traces must come from the same holomorphic function; they
cannot be chosen independently. The residual in the kernel is
\[
 \mathscr E_N(s,z)=\mathscr P_N(s,z)
                    -\mathscr D_N(s+z)-\mathscr K_N(s,z),
\]
where $\mathscr P_N$ is the odd-character correlation,
$\mathscr D_N$ is its exact diagonal, and $\mathscr K_N$ is the
continued Beta model. Their explicit formulas are recalled when we
estimate the three components below.
Proposition~\ref{proof:bd:HS} makes $T_N$
compact, selfadjoint, and Hilbert--Schmidt. Its nonzero eigenvalues
$\mu_j(N)$, counted with multiplicity, are real, with
\[
 \sum_j\mu_j^2(N)<\infty,\qquad
 \|T_N\|=\sup_j|\mu_j(N)|.
\]
The supremum is zero when $T_N=0$. The Goldbach singular series is
\[
 \mathfrak S(N)=2C_2\prod_{\substack{p\mid N\\p>2}}\frac{p-1}{p-2},
 \qquad C_2=\prod_{p>2}\left(1-\frac1{(p-1)^2}\right),
\]
as in~\eqref{eq:singular-series}.
We use the scale
\[
 \delta_N=\sqrt{\frac{\mathfrak S(N)}N}(\log N)^{3/2},
 \qquad \delta_N^2=\frac{\mathfrak S(N)(\log N)^3}{N},
\]
as in~\eqref{v3i:spectral-scale}.
We encode the size of these eigenvalues relative to $\delta_N$;
the Dirichlet zeros enter through GRH and are a separate set of data.

\subsection{Convergence of the spectral product}
\label{proof:v7:spectral-construction}

Define the spectral function by
\[
 \mathfrak X_N(q)=\prod_j
   \left(1-\frac{4\mu_j^2(N)}{\delta_N^2}q(1-q)\right),
 \qquad q\in\mathbb H,
\]
as in~\eqref{qp:spectral-function}.
The empty product is equal to one. Here $\mathbb H$ denotes the
quaternions. For a unit purely imaginary quaternion $I$, write
$\mathbb C_I=\{x+yI:x,y\in\mathbb R\}$. All factors at a fixed
$q\in\mathbb C_I$ lie in this commutative complex plane.

The term \emph{intrinsic entire slice regular} has a concrete meaning
here: the same real-coefficient power series defines an entire
function on every complex slice $\mathbb C_I$. It preserves each
slice and respects conjugation. This regularity concerns the spectral
variable $q$; the Fueter regularity of $\mathscr E_N(s,\bar z)$
concerns the separate arithmetic coordinates.

The ordinary Fredholm determinant extends the finite-dimensional formula
\[
 \det(\operatorname{Id}-tA)=\prod_j(1-ta_j)
\]
to a trace-class operator $A$, with eigenvalues $a_j$.
For positive compact $A$,
trace class means
$\sum_j a_j<\infty$, which ensures convergence of that product.
Here $T_N$ is known to be Hilbert--Schmidt, so its squared eigenvalues
are summable. Consequently we use $A=T_N^2$, whose eigenvalues are
$a_j=\mu_j(N)^2$, and obtain an ordinary Fredholm determinant.

\begin{proposition}[Convergence and determinant realization]
\label{proof:v9:spectral-convergence}
The product
\[
 \mathfrak X_N(q)=
 \prod_j\left(1-\frac{4\mu_j^2(N)}{\delta_N^2}q(1-q)\right),
 \qquad q\in\mathbb H,
\]
converges locally uniformly on
$\mathbb H$, independently of the enumeration of the eigenvalues, and
defines an intrinsic entire slice-regular function $\mathfrak X_N$.
The operator $T_N^2$ is positive and trace class, with
\[
 \operatorname{tr}(T_N^2)=\sum_j\mu_j^2(N)
                       =\|T_N\|_{\mathrm{HS}}^2.
\]
On the usual complex slice,
\[
 \mathfrak X_N(q)=\det\left(\operatorname{Id}_{\mathcal H_N}
       -\frac4{\delta_N^2}q(1-q)T_N^2\right).
\]
\end{proposition}

\begin{proof}
By Proposition~\ref{proof:bd:HS}, $\sum_j\mu_j^2(N)<\infty$.
On a compact subset of $\mathbb H$, each factor differs from $1$ by
at most a fixed multiple of $\mu_j^2(N)$. Thus the product converges
locally uniformly and is independent of the enumeration of the spectrum.
Within each complex slice its factors commute and are holomorphic.
Their real coefficients give the same real Taylor series on every slice,
so the limit is intrinsic and slice regular.

The spectral theorem gives
\[
 \operatorname{tr}(T_N^2)=\sum_j\mu_j^2(N)
                       =\|T_N\|_{\mathrm{HS}}^2<\infty.
\]
Hence $T_N^2$ is positive and trace class, and its Fredholm determinant
is exactly the product $\mathfrak X_N(q)$ on the usual complex slice. The operator and its determinant are therefore formed
on the complex Hilbert space $\mathcal H_N$; the real coefficients
then extend the scalar function to the other slices.
This proves the proposition.
\end{proof}

\subsection{Symmetries and zeros}
\label{proof:sec:spectral-symmetry}\label{proof:v7:no-trivial-zeros}
\label{proof:v7:zero-classification}

The quadratic factors determine both the symmetries and all zeros.
On a single complex slice the nonreal roots occur in conjugate
pairs. Allowing the imaginary unit $I$ to vary turns such a pair
into a sphere with fixed real part and fixed imaginary length.
The exceptional roots are real pairs, and only these affect the
norm criterion.

\begin{proposition}[Symmetries and normalization]
\label{proof:v9:spectral-symmetries}
The normalization and symmetries are
\begin{gather}
 \mathfrak X_N(0)=\mathfrak X_N(1)=1,
       \label{proof:sym:normalization}\\
 \mathfrak X_N(1-q)=\mathfrak X_N(q),
       \label{proof:sym:functional-equation}\\
 \mathfrak X_N(\bar q)=\overline{\mathfrak X_N(q)}.
       \label{proof:sym:conjugation}
\end{gather}
Combining reflection and conjugation gives
\begin{equation}\label{sym:hyperplane-reflection}
 \mathfrak X_N(1-\bar q)=\overline{\mathfrak X_N(q)}
 \qquad(q\in\mathbb H).
\end{equation}
The map $q\mapsto1-q$ is central inversion about $1/2$, whereas
$q\mapsto1-\bar q$ reflects across $\Re q=1/2$.
For the centered variable $\xi\in\mathbb H$,
\begin{equation}\label{sym:centered-product}
 \mathfrak X_N(1/2+\xi)
 =\prod_j\left(1-\frac{\mu_j^2(N)}{\delta_N^2}
                  +\frac{4\mu_j^2(N)}{\delta_N^2}\xi^2\right).
\end{equation}
This is even in $\xi$. On the critical hyperplane, for real $t$
and any unit purely imaginary quaternion $I$,
\begin{equation}\label{sym:critical-values}
 \mathfrak X_N(1/2+It)
 =\prod_j\left(1-\frac{\mu_j^2(N)}{\delta_N^2}(1+4t^2)\right)
 \in\mathbb R.
\end{equation}
The value is independent of $I$. Moreover,
$\mathfrak X_N(x)\ge1$ for real $x\notin(0,1)$.
\end{proposition}

\begin{proof}
The invariance of $q(1-q)$ under $q\mapsto1-q$ proves
\eqref{proof:sym:functional-equation}. Real coefficients give
\eqref{proof:sym:conjugation}, and combining the two gives
\eqref{sym:hyperplane-reflection}. Evaluation at $0$ and $1$ proves
\eqref{proof:sym:normalization}. Finally,
$(1/2+\xi)(1/2-\xi)=1/4-\xi^2$ gives the centered product
\eqref{sym:centered-product}; taking $\xi=It$ gives
\eqref{sym:critical-values}. All identities pass from finite products
to the limit by local uniform convergence.
For real $x\notin(0,1)$, each factor is at least one, giving the
asserted lower bound.
This proves the proposition.
\end{proof}

\begin{proposition}[Complete zero classification]
\label{proof:v9:spectral-zeros}
For each nonzero eigenvalue $\mu_j(N)$, put
$d_j=\delta_N^2/\mu_j^2(N)$. The zeros of $\mathfrak X_N$ are
exactly the zeros contributed by these factors, namely
\[
 \begin{cases}
 \displaystyle \Re q=\frac12,\quad
       |\Im q|=\frac12\sqrt{d_j-1},&d_j>1,\\[2mm]
 q=\frac12\quad\text{\normalfont (a double zero of the factor)},&d_j=1,\\[1mm]
 \displaystyle q=\frac12\pm\frac12\sqrt{1-d_j},&0<d_j<1.
 \end{cases}
\]
Thus every zero off the critical hyperplane is real, there are only
finitely many exceptional real pairs, and every real zero lies in
$(0,1)$. Each nonreal zero sphere is invariant under conjugation by
unit quaternions. On each complex slice, reflection and conjugation
preserve multiplicities, and the order of a zero at $1/2$ is even.
\end{proposition}

\begin{proof}
For the ratio $d_j=\delta_N^2/\mu_j^2(N)$ just defined, a
factor vanishes when $q(1-q)=d_j/4$. Writing $q=x+yI$ gives
\[
 q(1-q)=x(1-x)+y^2+(1-2x)yI.
\]
If $y\ne0$, this forces $x=1/2$ and $y^2=(d_j-1)/4$.
The equation is independent of $I$, which gives the stated two-sphere.
For real $q$, the ordinary quadratic formula gives the exceptional
pair when $d_j<1$ and the double central root when $d_j=1$.
For $x\notin(0,1)$, every factor is at least $1$, so there are no
other real zeros.

There are no zeros arising solely from the infinite product. Indeed,
at a point where all factors $1-u_j$ are nonzero, choose $J$ such that
$|u_j|\le1/2$ for $j\ge J$. Then
\[
 \prod_{j\ge J}|1-u_j|
 \ge\prod_{j\ge J}(1-|u_j|)>0,
\]
because $\sum_j|u_j|<\infty$ and
$\log(1-|u_j|)\ge-2|u_j|$. The finite preceding product is nonzero.
Compactness also allows only finitely many eigenvalues with
$|\mu_j|>\delta_N$, so there are only finitely many exceptional pairs.
On each complex slice, conjugation and reflection preserve multiplicity;
at $1/2$ the order is even by the centered product.
Conjugation by a unit quaternion preserves the real part and the
length of the imaginary part, so it preserves each nonreal zero sphere.
This proves the proposition.
\end{proof}

\subsection{From a zero-free region to a norm bound}

The zero classification turns the exclusion condition into an exact
bound for the largest eigenvalue modulus. The excluded half-space is
open, so equality is allowed in the norm bound. Here
$\mathrm R(N)$ denotes the assertion that every zero of
$\mathfrak X_N$ has real part $1/2$.

\begin{theorem}\label{proof:qp:dictionary}
Under the standing assumptions,
\[
 \mathrm R(N)\quad\Longleftrightarrow\quad\|T_N\|\le\delta_N.
\]
For $0<\eta<1/2$,
\[
 \mathfrak X_N(q)\ne0\quad(\Re q>1-\eta)
 \quad\Longleftrightarrow\quad
 \|T_N\|\le\frac{\delta_N}{2\sqrt{\eta(1-\eta)}}.
\]
In particular, if $\|T_N\|\le M_N$ with $M_N>\delta_N$, the
width is
\[
 \eta(M_N)=\frac{1-\sqrt{1-\delta_N^2/M_N^2}}2.
\]
The corresponding half-space is zero-free, with $\eta(M_N)\asymp\delta_N^2/M_N^2$
when $\eta(M_N)\le1/4$.
\end{theorem}

\begin{proof}
A factor has zeros off the critical hyperplane exactly when
$|\mu_j|>\delta_N$. Taking the supremum of the eigenvalue magnitudes
proves the first equivalence, including the zero operator.
The preceding product argument shows that these are all possible zeros.
It is therefore enough to test the spectral factors individually; no
separate condition on the limiting product is needed.

Only the larger root of an exceptional pair can enter
$\Re q>1-\eta$. Its exclusion is equivalent to
\[
 \frac12+\frac12\sqrt{1-\frac{\delta_N^2}{\mu_j^2}}\le1-\eta,
 \quad\text{or}\quad
 |\mu_j|\le\frac{\delta_N}{2\sqrt{\eta(1-\eta)}}.
\]
Squaring is reversible because $1-2\eta>0$. Nonexceptional eigenvalues
already satisfy the latter bound, which proves the second equivalence.
Finally, suppose $\|T_N\|\le M_N$ with $M_N>\delta_N$.
Choosing $\eta\in(0,1/2)$ to satisfy
$2\sqrt{\eta(1-\eta)}=\delta_N/M_N$ gives the stated width
$\eta(M_N)$, and the second equivalence ensures that the
corresponding half-space is zero-free. Rationalizing yields
\[
 \eta(M_N)=\frac{\delta_N^2/M_N^2}
 {2\left(1+\sqrt{1-\delta_N^2/M_N^2}\right)}.
\]
The denominator lies between $2$ and $4$, so
$\eta(M_N)\asymp\delta_N^2/M_N^2$.
This proves the theorem.
\end{proof}

We use this conversion to compare the proved GRH estimate with the
stronger bounds required by S and R.

\subsection{GRH proves a baseline zero-free region}
\label{proof:sec:grh-bound}

We estimate the odd-character correlation, the diagonal, and the
Beta model separately. An average of squared logarithmic derivatives
controls the first component; positivity controls the second; and
the continued model formula controls the third. These bounds give
a zero-free region without requiring cancellation between the
components of the residual.

\begin{theorem}[The baseline bound under GRH]\label{proof:grh:main-theorem}
For every sufficiently large even $N$ satisfying $\mathrm{GRH}(N)$,
one has
\begin{equation}\label{proof:grh:operator-bound}
 \|T_N\|\le\|T_N\|_{\mathrm{HS}}\ll(\log N)^2.
\end{equation}
Consequently the spectral function has the zero-free region
\begin{equation}\label{proof:grh:zero-free-width}
 \mathfrak X_N(q)\ne0\quad(\Re q>1-\eta_N^{\mathrm{GRH}}),
 \qquad
 \eta_N^{\mathrm{GRH}}\asymp\frac{\mathfrak S(N)}{N\log N}.
\end{equation}
The constants and the required lower bound on $N$ are absolute.
\end{theorem}

For this proof, retain $N>e^{10}$ and take the auxiliary reference line
$\sigma_{\mathrm{ref}}=2/\log N<1/5$. This line may lie to the right
of the Hardy boundary $\beta=1/8$; the auxiliary Poisson comparison
below is therefore established for $0<x\le1/5$. The Hardy strip
keeps its boundaries $\alpha_N$ and $\beta$ and its fixed interior
line $\sigma_*=1/9$.
The required estimate must be uniform as the left boundary approaches
the critical line; compactness for each fixed $N$ is not enough.

\subsubsection{An averaged bound for logarithmic derivatives}

The correlation component requires a mean-square bound that remains
uniform as the left Hardy boundary approaches the critical line.
We first prove such a bound on a reference line by smoothing a short
Dirichlet polynomial, then transfer it to the required lines using
the positive zero sum and Poisson kernels. Here $X_N^-$ consists
of all odd characters modulo $N$, including imprimitive ones, and
\[
 g_\chi(s)=\frac{L'}L(1/2+s,\chi).
\]
For each $\chi$, write $\chi^*$ for its primitive inducing character
and $d_\chi\mid N$ for the conductor. Under GRH its nontrivial
zeros are $1/2+i\gamma_\chi$, counted with multiplicity.

\begin{lemma}[Averaged logarithmic derivatives]
\label{proof:grh:averaged-energy-lemma}
For $N>e^{10}$ and $0<\sigma\le\beta$,
\begin{equation}\label{grh:averaged-energy}
 \frac2{\varphi(N)}\sum_{\chi\in X_N^-}
 \int_{\mathbb R}\frac{|g_\chi(\sigma+it)|^2}{1+t^2}\,\mathrm{d}t
 \ll(\log N)^2+\sigma^{-2}.
\end{equation}
The implied constant is absolute.
\end{lemma}

\begin{proof}
\emph{The reference line.}
Fix $\chi\in X_N^-$, and let $\chi^*$ denote the primitive
Dirichlet character inducing $\chi$, of conductor $d_\chi\mid N$.
Write $1/2+i\gamma_\chi$ for the nontrivial zeros of
$L(s,\chi^*)$, and define
\[
 Q_\chi(x,t)=\sum_{\gamma_\chi}
             \frac{x}{x^2+(t-\gamma_\chi)^2},
 \qquad x>0,\quad t\in\mathbb R.
\]
The sum runs over all such ordinates $\gamma_\chi$, with each zero
counted according to its multiplicity.
We first establish the comparison
\begin{equation}\label{proof:grh:Poisson-comparison}
 |Q_\chi(x,t)-\Re g_\chi(x+it)|
       \ll\log\bigl(N(|t|+2)\bigr)
 \qquad(0<x\le1/5)
\end{equation}
and the reference-line estimate
\begin{equation}\label{proof:v7l:reference-energy-bound}
 \frac2{\varphi(N)}\sum_{\chi\in X_N^-}
 \int_{\mathbb R}\frac{|g_\chi(\sigma_{\mathrm{ref}}+it)|^2}{1+t^2}\,\mathrm{d}t
 \ll(\log N)^2.
\end{equation}
Compare the logarithmic-derivative formulas from the Hadamard
products at $u=1/2+x+it$ and $u_0=2+it$.
For the completed primitive function
$\mathcal L(u,\chi^*)=(d_\chi/\pi)^{(u+1)/2}
\Gamma((u+1)/2)L(u,\chi^*)$, the identity in
Lemma~\ref{proof:v9:primitive-zero-data} reads
\[
 \frac{\mathcal L'}{\mathcal L}(u,\chi^*)
 -\frac{\mathcal L'}{\mathcal L}(u_0,\chi^*)
 =\sum_\rho\left(\frac1{u-\rho}-\frac1{u_0-\rho}\right).
\]
Taking real parts gives $Q_\chi(x,t)-Q_\chi(3/2,t)$ on the
right. The conductor terms cancel, leaving the logarithmic
derivatives of $L$ and a Gamma difference on the left.
The reference sum satisfies
$Q_\chi(3/2,t)\ll\log\bigl(d_\chi(|t|+2)\bigr)$:
group its terms according to
$m\le|t-\gamma_\chi|<m+1$ and use
the unit-interval bound in Lemma~\ref{proof:v9:primitive-zero-data},
\[
 \#\{\rho:v\le\Im\rho<v+1\}
 \ll\log\bigl(d_\chi(|v|+2)\bigr),
\]
with the summable factor $(1+m^2)^{-1}$.
The logarithmic derivative at $2+it$
is uniformly bounded by its Dirichlet series. The Gamma
difference is uniformly bounded for $0<x\le1/5$, and the imprimitive
Euler corrections contribute $O(\log N)$, as in
Lemma~\ref{proof:lem:zero-expansion}. These observations prove
\eqref{proof:grh:Poisson-comparison} without any zero-spacing assumption.

For the reference-line estimate, set $X=N/2$ and $x=\sqrt X$
and use the smoothed Dirichlet polynomial
\[
 P_{\chi,X}(u)=
 \sum_{n\le X}\Lambda(n)\chi(n)n^{-u}w_X(n).
\]
Its cutoff weight is
\[
 w_X(n)=\begin{cases}
 1,&n\le x,\\
 \log(X/n)/\log x,&x<n\le X,\\
 0,&n>X.
 \end{cases}
\]
For $u=1/2+\sigma+it$, the Perron kernel
$(X^v-x^v)/(v^2\log x)$ has a simple pole of residue $1$
at zero. In particular, there is no derivative term at that
point. The residue of $-L'/L(u+v,\chi)$ at a zero
$v=\rho-u$ of multiplicity $m$ is $-m$. Shifting the contour gives
\begin{equation}\label{proof:grh:smoothed-truncation}
 -\frac{L'}L(u,\chi)-P_{\chi,X}(u)
 =\frac1{\log x}\sum_\rho
          \frac{X^{\rho-u}-x^{\rho-u}}{(\rho-u)^2}.
\end{equation}
The sum includes, with multiplicity, the primitive nontrivial and
trivial zeros and the zeros of the additional Euler factors.
For primes $p\mid N$ with $p\nmid d_\chi$, these zeros have
real part zero and ordinates
$(\arg\chi^*(p)+2\pi k)/\log p$, $k\in\mathbb Z$.
Since $\chi$ is odd, there is no pole at $u=1$. For the contour
shift, choose left vertical lines on which $\Re(u+v)$ is a negative
even integer. The local zero count permits horizontal ordinates
of size $T$ at distance $\gg1/\log(N(T+2))$ from the zero ordinates,
including those of the Euler factors. On each fixed horizontal
segment, the logarithmic derivative is then
$O(\log^2(N(T+2)))$, so its integral against the Perron kernel
tends to zero. First let the height tend to infinity, then move
the left line to minus infinity. The functional equation bounds the
logarithmic derivative on these lines; $x>1$ supplies exponential
decay on the left. The quadratic denominator and the local zero
count give absolute convergence of the nontrivial and Euler
zero sums. The trivial-zero sum has additional geometric decay.

The nontrivial-zero contribution is bounded by
$\epsilon_\sigma Q_\chi(\sigma,t)$, where
\[
 \epsilon_\sigma=
       \frac{x^{-\sigma}+X^{-\sigma}}{\sigma\log x}.
\]
The remaining zeros contribute
$O((\log N/\log x)x^{-1/2-\sigma})$, uniformly in $t$.
Here the Euler zeros are the zeros of the additional factors
$1-\chi^*(p)p^{-u}$, where $p$ ranges over the prime divisors
of $N$ that do not divide $d_\chi$. Since $|\chi^*(p)|=1$,
the zeros of each such factor are
\[
 \rho_{p,k}=i\left(t_p+\frac{2\pi k}{\log p}\right),
 \qquad t_p=\frac{\arg\chi^*(p)}{\log p},\quad k\in\mathbb Z.
\]
Thus, for each fixed prime $p$, their ordinates have spacing
$2\pi/\log p$, which gives
\[
 \sum_{k\in\mathbb Z}
 \frac1{(1/2+\sigma)^2+(t-t_p-2\pi k/\log p)^2}
 \ll 1+\log p.
\]
Since $\sum_{p\mid N}(1+\log p)\ll\log N$, summing over these
primes gives the stated error bound. The primitive trivial zeros
give a smaller bound.
At the reference line, put
\[
 a_N=\sigma_{\mathrm{ref}}\log x
     =1-\frac{\log2}{\log N}>\frac9{10}.
\]
Thus
\[
 \epsilon_{\sigma_{\mathrm{ref}}}
 =\frac{e^{-a_N}+e^{-2a_N}}{a_N}
 \le\frac{10}{9}\left(e^{-9/10}+e^{-9/5}\right)<\frac23.
\]
The comparison \eqref{proof:grh:Poisson-comparison} therefore allows the logarithmic
derivative on the right of \eqref{proof:grh:smoothed-truncation} to be
absorbed into the left side. Explicitly, the comparison
$Q_\chi(\sigma_{\mathrm{ref}},t)\le |g_\chi(\sigma_{\mathrm{ref}}+it)|+
C\log(N(|t|+2))$ gives
\[
 (1-\epsilon_{\sigma_{\mathrm{ref}}})|g_\chi(\sigma_{\mathrm{ref}}+it)|
 \le |P_{\chi,X}(1/2+\sigma_{\mathrm{ref}}+it)|
       +C\log(N(|t|+2)).
\]
The bound $1-\epsilon_{\sigma_{\mathrm{ref}}}>1/3$ is uniform in $N$,
the character, and the height. Hence
\[
 |g_\chi(\sigma_{\mathrm{ref}}+it)|
 \ll |P_{\chi,X}(1/2+\sigma_{\mathrm{ref}}+it)|
                      +\log\bigl(N(|t|+2)\bigr).
\]
Odd-character orthogonality leaves exactly the diagonal for this
polynomial: $n,m\le N/2$, and $w_X(N/2)=0$ removes the possible
endpoint term.
Thus averaging removes every oscillatory cross term before the
integration in height. The resulting estimate does not require a bound
for an individual character polynomial.
Hence
\[
 \frac2{\varphi(N)}\sum_{\chi\in X_N^-}
 |P_{\chi,X}(1/2+\sigma+it)|^2
 =\sum_{(n,N)=1}\frac{\Lambda^2(n)w_X^2(n)}{n^{1+2\sigma}}
 \ll\sigma^{-2}.
\]
For the last step, Chebyshev's bound gives
$B(u):=\sum_{n\le u}\Lambda^2(n)\ll u\log(2u)$, so partial summation yields
\[
 \sum_{n\ge1}\frac{\Lambda^2(n)}{n^{1+2\sigma}}
 =(1+2\sigma)\int_1^\infty\frac{B(u)}{u^{2+2\sigma}}\,\mathrm{d}u
 \ll\int_1^\infty\frac{\log(2u)}{u^{1+2\sigma}}\,\mathrm{d}u
 \ll\sigma^{-2}.
\]
Average the square of the preceding pointwise bound and integrate
against $(1+t^2)^{-1}$. The diagonal estimate then shows that the
left side of \eqref{proof:v7l:reference-energy-bound} is $O((\log N)^2)$.

\emph{Transfer to the Hardy boundaries.}
For each character, write
\begin{equation}\label{proof:v7l:weighted-energy-definition}
 \mathscr E_\chi(\sigma)=
 \int_{\mathbb R}\frac{|g_\chi(\sigma+it)|^2}{1+t^2}\,\mathrm{d}t.
\end{equation}
For $0<\sigma\le\beta$ and any $0<\sigma_{\mathrm{ref}}\le1/5$,
we claim that
\begin{equation}\label{proof:v7l:energy-transfer-bound}
 \mathscr E_\chi(\sigma)\ll
 \begin{cases}
 (\sigma_{\mathrm{ref}}/\sigma)^2
    \bigl(\mathscr E_\chi(\sigma_{\mathrm{ref}})+(\log N)^2\bigr),
       &\sigma\le\sigma_{\mathrm{ref}},\\
 \mathscr E_\chi(\sigma_{\mathrm{ref}}),&\sigma_{\mathrm{ref}}\le\sigma.
 \end{cases}
\end{equation}
For any $0<\sigma\le\beta$, we first relate the
weighted norm of the full logarithmic derivative to that of its
real part:
\begin{equation}\label{proof:grh:weighted-realpart}
 \int_{\mathbb R}\frac{|g_\chi(\sigma+it)|^2}{1+t^2}\,\mathrm{d}t
 \le2\int_{\mathbb R}
       \frac{|\Re g_\chi(\sigma+it)|^2}{1+t^2}\,\mathrm{d}t
       +\pi|g_\chi(\sigma+1)|^2.
\end{equation}
To apply Parseval with the present weight, use the disk function
\[
 G(\zeta)=g_\chi\!\left(\sigma+\frac{1+\zeta}{1-\zeta}\right).
\]
Here the disk Hardy space $H^2$ uses one circular boundary:
bounded radial square means are equivalent to square summability
of the Taylor coefficients. The Cayley map converts its normalized
boundary measure into the weight used in the displayed inequality.
The GRH logarithmic-derivative bound and the absolutely
convergent Dirichlet series farther to the right give
\[
 |G(\zeta)|\ll_{N,\sigma}\log^2(2+|1-\zeta|^{-1}).
\]
This bound gives uniformly bounded radial square means, so $G\in H^2$.
With normalized circle measure, if $G(\zeta)=\sum a_n\zeta^n$, then
\[
 \|G\|_2^2=2\|\Re G\|_2^2-\Re(a_0^2).
\]
The Cayley map sends that measure to $\mathrm{d}t/(\pi(1+t^2))$.
This change of variables and $-\Re(a_0^2)\le|a_0|^2$ give
\eqref{proof:grh:weighted-realpart}; its last term is uniformly
$O(1)$ by the Dirichlet series at $3/2+\sigma$.

For $0<\sigma\le\min(\sigma_{\mathrm{ref}},\beta)$, positivity gives
$Q_\chi(\sigma,t)\le(\sigma_{\mathrm{ref}}/\sigma)Q_\chi(\sigma_{\mathrm{ref}},t)$.
Writing $K=\sigma_{\mathrm{ref}}/\sigma\ge1$,
\eqref{proof:grh:Poisson-comparison} implies
\[
 |\Re g_\chi(\sigma+it)|
 \le K|g_\chi(\sigma_{\mathrm{ref}}+it)|
       +C(K+1)\log\bigl(N(|t|+2)\bigr).
\]
Square this inequality and integrate against $(1+t^2)^{-1}$.
Since $\int_{\mathbb R}\log^2(|t|+2)(1+t^2)^{-1}\,\mathrm{d}t<\infty$,
the integrated logarithmic error is $O((\log N)^2)$.
Applying \eqref{proof:grh:weighted-realpart} therefore gives
\[
 \mathscr E_\chi(\sigma)\ll
 K^2\bigl(\mathscr E_\chi(\sigma_{\mathrm{ref}})+(\log N)^2\bigr).
\]
The bounded center term is absorbed since $K\ge1$ and $N\ge4$.
This proves the first case of \eqref{proof:v7l:energy-transfer-bound}.
For $\sigma\ge\sigma_{\mathrm{ref}}$, the disk $H^2$ estimate bounds
the squared modulus by the Poisson integral of its boundary squared
modulus. This is the harmonic majorant used here. Under the same
Cayley change of variable it becomes
\[
 |g_\chi(\sigma+it)|^2
 \le(P_\ell*|g_\chi(\sigma_{\mathrm{ref}}+i\,\cdot)|^2)(t),\qquad
 P_\ell(t)=\frac{\ell}{\pi(\ell^2+t^2)},\quad \ell=\sigma-\sigma_{\mathrm{ref}}.
\]
Here $*$ denotes convolution in $t$; for $\ell=0$ the formula is
understood as the identity. The disk $H^2$ boundary representation
has no additional point mass at the point corresponding to infinity.
Multiply the resulting majorant by $\pi P_1(t)$ and integrate.
By nonnegativity, Tonelli's theorem permits interchanging the two integrals.
Then $P_1*P_\ell=P_{1+\ell}$ and $P_{1+\ell}\le(1+\ell)P_1$ give
$\mathscr E_\chi(\sigma)\le(1+\ell)\mathscr E_\chi(\sigma_{\mathrm{ref}})$.
The factor is bounded since $0\le \ell\le\beta$. This proves the
second case of \eqref{proof:v7l:energy-transfer-bound}.

\emph{The character average.}
Set $\sigma_{\mathrm{ref}}=2/\log N$.
Equation~\eqref{proof:v7l:reference-energy-bound} bounds the weighted
square integrals of the logarithmic derivatives on this line,
averaged over all odd characters modulo $N$:
\[
 \frac2{\varphi(N)}\sum_{\chi\in X_N^-}
       \mathscr E_\chi(\sigma_{\mathrm{ref}})
 =\frac2{\varphi(N)}\sum_{\chi\in X_N^-}
       \int_{\mathbb R}
       \frac{|g_\chi(\sigma_{\mathrm{ref}}+it)|^2}{1+t^2}\,\mathrm{d}t
 \ll(\log N)^2.
\]
For $0<\sigma\le\min(\sigma_{\mathrm{ref}},\beta)$, average the first case of
\eqref{proof:v7l:energy-transfer-bound} over the odd characters.
With $K=\sigma_{\mathrm{ref}}/\sigma$, this gives
$O(K^2(\log N)^2)=O(\sigma^{-2})$ on the left of the reference line.
Here $K^2(\log N)^2=4\sigma^{-2}$ and
$2|X_N^-|/\varphi(N)=1$, so the additional $(\log N)^2$
term has the same bound after averaging.
For $\sigma_{\mathrm{ref}}\le\sigma\le\beta$, average the second case of
\eqref{proof:v7l:energy-transfer-bound} to obtain $O((\log N)^2)$.
These two ranges prove \eqref{grh:averaged-energy} with the
asserted uniform constant.
This proves the lemma.
\end{proof}

\subsubsection{The three components of the residual operator}

We apply the averaged estimate to the correlation operator and
bound the diagonal and Beta-model operators on the same Hardy space.
Positivity gives trace bounds for the first two components; the
continued Beta model is controlled through its boundary kernel.

The three components act by boundary integration against the
respective kernels
\[
 \frac{\mathscr P_N(s,\bar z)}{(1+s)(1+\bar z)},\qquad
 \frac{\mathscr D_N(s+\bar z)}{(1+s)(1+\bar z)},\qquad
 \frac{\mathscr K_N(s,\bar z)}{(1+s)(1+\bar z)}.
\]
We denote these operators by $\mathcal P_N$, $\mathcal D_N$, and
$\mathcal K_N$, so that
$T_N=\mathcal P_N-\mathcal D_N-\mathcal K_N$.
\begin{proposition}\label{proof:grh:component-bounds}
The operators $\mathcal P_N$ and $\mathcal D_N$ are positive
and trace class, and
\begin{align*}
 \|\mathcal P_N\|_{\mathrm{HS}}
       &\le\operatorname{tr}\mathcal P_N\ll (\log N)^2+\alpha_N^{-2},\\
 \|\mathcal D_N\|_{\mathrm{HS}}
       &\le\operatorname{tr}\mathcal D_N\ll\alpha_N^{-2},\\
 \|\mathcal K_N\|_{\mathrm{HS}}&\ll\alpha_N^{-2}.
\end{align*}
The implied constants are independent of $N$ and $\alpha_N$.
\end{proposition}

\begin{proof}
For the correlation component, the Gram representation is
\[
 \mathscr P_N(s,\bar z)=\frac2{\varphi(N)}
       \sum_{\chi\in X_N^-}g_\chi(s)\overline{g_\chi(z)},
\]
as in~\eqref{qp:Gram}.
It makes $\mathcal P_N$ positive and finite rank.
The boundary norm is the sum over $\Re s=\alpha_N$ and
$\Re s=\beta$, with $|1+s|^{-2}\asymp(1+|\Im s|^2)^{-1}$.
Thus the averaged bound for weighted square integrals in
\eqref{grh:averaged-energy} applies to each boundary and gives
\[
 \|\mathcal P_N\|_{\mathrm{HS}}
 \le\operatorname{tr}\mathcal P_N
 =\frac2{\varphi(N)}\sum_{\chi\in X_N^-}
       \left\|\frac{g_\chi(s)}{1+s}\right\|_{\mathcal H_N}^2
 \ll (\log N)^2+\alpha_N^{-2}.
\]
For the diagonal component, recall from
\eqref{eq:goldbach-arithmetic-kernels} that
$\mathscr D_N(w)=\sum_{(n,N)=1}\Lambda^2(n)n^{-1-w}$,
and set $d_n(s)=n^{-s}/(1+s)$.
Each $d_n$ belongs to $\mathcal H_N$, and
\[
 \|d_n\|^2=\pi\left(\frac{n^{-2\alpha_N}}{1+\alpha_N}
                  +\frac{n^{-2\beta}}{1+\beta}\right).
\]
Because the inner product is linear in its first variable,
$d_n\otimes d_n$ denotes the operator
$h\mapsto\langle h,d_n\rangle d_n$.
The positive rank-one series
$\sum_{(n,N)=1}\Lambda^2(n)n^{-1}d_n\otimes d_n$
therefore has total trace
\[
 \operatorname{tr}\mathcal D_N
 =\pi\left(\frac{\mathscr D_N(2\alpha_N)}{1+\alpha_N}
                     +\frac{\mathscr D_N(2\beta)}{1+\beta}\right)
 \ll\alpha_N^{-2}.
\]
The bound $\sum_{n\le u}\Lambda^2(n)\ll u\log(2u)$ and partial
summation give $\mathscr D_N(2\sigma)\ll\sigma^{-2}$ for
$0<\sigma\le\beta$. Hence the rank-one series converges in trace
norm and defines $\mathcal D_N$ with the required diagonal kernel.
For this positive series, convergence in trace norm means that the
sum of the traces of the omitted terms tends to zero; it implies
convergence in operator norm as well. The operators
$\mathcal P_N$ and $\mathcal D_N$ act on
$\mathcal H_N$, as does $T_N$ by Proposition~\ref{proof:bd:HS}.
The identity $\mathscr E_N=\mathscr P_N-\mathscr D_N-\mathscr K_N$ shows that
$\mathcal K_N=\mathcal P_N-\mathcal D_N-T_N$ has the boundary
kernel corresponding to $\mathscr K_N$.

For the model component, the arithmetic and analytic factors are
separated as follows:
\begin{equation}\label{proof:grh:model-functional-form}
 \mathscr K_N(s,\bar z)=\mathfrak S(N)
 \left(\frac\pi N\right)^{s+\bar z}P_N(s+\bar z)
 \zeta(1-s-\bar z)\mathcal X_1(s)\mathcal X_1(\bar z),
\end{equation}
as in~\eqref{v3:model-functional-form}.
Here
\[
 P_N(w)=\prod_{\substack{p>2\\p\nmid N}}
       \left(1+\frac{p^{-w}}{p-2}\right),\qquad
 \mathcal X_1(s)=\frac{\Gamma((3/2-s)/2)}{\Gamma((3/2+s)/2)}.
\]
We estimate these factors uniformly as the left boundary approaches
zero. Let $0<\sigma\le2\beta$ and put
$P_*(\sigma)=\prod_{p>2}(1+p^{-\sigma}/(p-2))$.
The convergent product comparison
\[
 \left(1+\frac{p^{-\sigma}}{p-2}\right)(1-p^{-1-\sigma})
 \le1+\frac2{p(p-2)}
\]
gives $P_*(\sigma)\ll\zeta(1+\sigma)\ll\sigma^{-1}$.
Moreover,
\[
 \mathfrak S(N)P_N(\sigma)
 =2C_2P_*(\sigma)
   \prod_{\substack{p\mid N\\p>2}}
        \frac{p-1}{p-2+p^{-\sigma}}.
\]
Each finite factor is at most $p^\sigma$, and
$\prod_{p\mid N}p\le N$. Hence, for $\Re w=\sigma$ in this range,
$|N^{-w}\mathfrak S(N)P_N(w)|\ll\sigma^{-1}$.
To bound the zeta factor, apply Euler summation at
$R=2+|\tau|$. The truncated sum and the pole term are each
$O(R^\sigma/\sigma)$, and the remainder is $O(R^\sigma)$.
Consequently
$|\zeta(1-\sigma-i\tau)|\ll\sigma^{-1}(2+|\tau|)^{2\beta}$
for $0<\sigma\le2\beta$. Stirling's formula gives
$\mathcal X_1(s)\ll(1+|\Im s|)^{-\Re s}$ uniformly for
$0\le\Re s\le\beta$, in particular a uniform bound there.
The functional form
\eqref{proof:grh:model-functional-form} therefore bounds the squared
model kernel on each boundary product by
\[
 C\alpha_N^{-4}
 \frac{(2+|t-y|)^{4\beta}}{(1+t^2)(1+y^2)}.
\]
The displayed majorant is integrable because $4\beta<1$; explicitly,
\[
 \int_{\mathbb R}
 \frac{\mathrm{d}t}{(1+t^2)(1+(t-u)^2)}=\frac{2\pi}{4+u^2}.
\]
Thus the boundary-kernel operator has Hilbert--Schmidt norm
$O(\alpha_N^{-2})$.
The preceding identity identifies the restriction of this
boundary-kernel operator to the Hardy trace space with
$\mathcal K_N$. Restricting an operator to a closed subspace
cannot increase its Hilbert--Schmidt norm, which proves the
model bound.

This proves the proposition.
\end{proof}

\subsubsection{Proof of the baseline spectral zero-free theorem}

The triangle inequality gives the operator bound; the zero--norm
conversion then gives the stated width.

\begin{proof}[Proof of Theorem~\ref{proof:grh:main-theorem}]
Applying the triangle inequality to the three bounds in
Proposition~\ref{proof:grh:component-bounds} gives
\[
 \begin{aligned}
 \|T_N\|_{\mathrm{HS}}
 &\le\|\mathcal P_N\|_{\mathrm{HS}}
      +\|\mathcal D_N\|_{\mathrm{HS}}
      +\|\mathcal K_N\|_{\mathrm{HS}}\\
 &\ll (\log N)^2+\alpha_N^{-2}\ll (\log N)^2.
 \end{aligned}
\]
Combining this with $\|T_N\|\le\|T_N\|_{\mathrm{HS}}$ proves
\eqref{proof:grh:operator-bound}.

Choose $\eta_N^{\mathrm{GRH}}$ to be a sufficiently small
absolute multiple of
$\mathfrak S(N)/(N\log N)$. Since $\mathfrak S(N)\ll\log\log N$, this width
tends to zero and lies in $(0,1/2)$ for sufficiently large $N$.
Using $\delta_N^2=\mathfrak S(N)(\log N)^3/N$, the corresponding
norm bound is
\[
 \frac{\delta_N}
      {2\sqrt{\eta_N^{\mathrm{GRH}}(1-\eta_N^{\mathrm{GRH}})}}
 \ge\frac{\delta_N}{2\sqrt{\eta_N^{\mathrm{GRH}}}}.
\]
The right-hand side is a fixed multiple of $(\log N)^2$.
Choosing the width's constant sufficiently small makes this multiple
large enough to dominate the operator bound.
The choice is made once, using the absolute constant in the operator
estimate, and hence remains uniform as the modulus varies.
Theorem~\ref{proof:qp:dictionary} then gives the stated zero-free half-space,
including its quaternionic extension. The chosen width has the
order asserted in \eqref{proof:grh:zero-free-width}.
This proves the theorem.
\end{proof}

\subsection{The scale of the left Hardy boundary}

The normalization used above also balances the available analytic
estimates. Keep $\beta=1/8$ fixed, let $0<\alpha<1/10$, and write
$\widehat T_{N,\alpha}$ for the operator on this strip and
$\mathcal Q_\alpha$ for the squared norm of its Mellin functional.
The proof of Proposition~\ref{proof:grh:component-bounds} applies
with $\alpha$ in place of $\alpha_N$. The Fourier calculation in
Proposition~\ref{proof:bm:norm}, performed with the same endpoints,
gives
\[
 \mathcal Q_\alpha=\frac1{2\pi}\int_0^\infty
 \frac{(Nu)^{2\alpha}|\Phi(u)-u\Phi'(u)|^2}
      {1+(Nu)^{-2(\beta-\alpha)}}\,\frac{\mathrm{d}u}{u}.
\]
On the fixed support, $u^{2\alpha}$ is bounded above and below
uniformly for $0<\alpha<1/10$, and the denominator lies between
one and two for sufficiently large $N$. Consequently,
\[
 \mathcal Q_\alpha\asymp_\phi N^{2\alpha},\qquad
 \|\widehat T_{N,\alpha}\|\le M_\alpha
 :=C\bigl((\log N)^2+\alpha^{-2}\bigr).
\]
The constants are uniform in $\alpha$ in this range, for
sufficiently large $N$. For left boundary $\alpha$, the
Cauchy--Schwarz argument in \eqref{v3f:master-error-formula} gives
an error bound $\mathcal Q_\alpha\|\widehat T_{N,\alpha}\|$.
Setting $x=\alpha\log N$, we have
\[
 \mathcal Q_\alpha M_\alpha
 \asymp_\phi(\log N)^2e^{2x}(1+x^{-2}).
\]
The last factor tends to infinity both as $x\downarrow0$ and
as $x\to\infty$, and has a positive finite minimum in between.
Thus $\alpha\asymp1/\log N$ gives the best order obtainable
from these separate bounds; it does not assert optimality of
the actual residual norm. The chosen $\alpha_N=1/\log N$
also gives $N^{2\alpha_N}=e^2$.
Appendix~\ref{proof:app:param} develops the parameter dependence.

The next section proves the exact quadratic-form identity and the
uniform functional bound, then combines them with the spectral
estimates to bound the Goldbach error and prove
Theorems~\ref{thm:goldbach-S} and \ref{thm:goldbach-R}.

\section{Hardy: the arithmetic error as an exact quadratic form}
\label{proof:sec:v3-mellin-operator}\label{proof:v7:hardy-error}\label{proof:detail:hardy}

We express the signed error in the weighted Goldbach correlation as
an exact quadratic form of the residual operator. The construction
has two parts: Mellin inversion identifies the arithmetic error,
and a bounded functional on the Hardy space selects the vector at
which the operator gives that error. We then combine the Hardy
and spectral estimates to bound the relative Goldbach error in
terms of a zero-free width and prove the two main theorems.

Throughout, $N\ge4$ is even. The weight $\phi$ is a fixed nonzero
function in $C_c^\infty(\mathbb R)$ supported in $[1/3,2/3]$, with
$0\le\phi\le1$ and $\phi(u)=\phi(1-u)$. Put
\[
 \Phi(u)=\sqrt u\,\phi(u),\qquad
 c_\phi=\int_0^1\phi(u)^2\,\mathrm{d}u.
\]
The arithmetic quantities to be compared are
\[
 \begin{aligned}
 \mathcal G(N)&=\frac1N
   \sum_{\substack{n+m=N\\(nm,N)=1}}
       \Lambda(n)\Lambda(m)\phi(n/N)\phi(m/N),\\
 \mathfrak S(N)&=2C_2\prod_{\substack{p\mid N\\p>2}}
                           \frac{p-1}{p-2},\qquad
 C_2=\prod_{p>2}\left(1-\frac1{(p-1)^2}\right),
 \end{aligned}
\]
as in~\eqref{eq:goldbach-correlation} and~\eqref{eq:singular-series}.
The sum is over ordered pairs of positive integers. Whenever
$\mathrm{GRH}(N)$ is assumed, it means that every nontrivial zero
of every primitive Dirichlet $L$-function whose conductor divides
$N$ has real part $1/2$. The Hardy construction will additionally
require $N>e^{10}$.
Our target is $c_\phi\mathfrak S(N)-\mathcal G(N)$; retaining its
sign gives an exact identity before any norm estimate is applied.

\subsection{Mellin inversion and the arithmetic residual}
\label{proof:detail:mellin}

We evaluate the character, diagonal, and model integrals on a common
Mellin contour. Their combination will represent the same arithmetic
discrepancy on every permitted contour. The residual is formed by
subtracting the exact diagonal and the continuous model from the
character kernel:
\[
 \mathscr E_N(s,z)=\mathscr P_N(s,z)
                    -\mathscr D_N(s+z)-\mathscr K_N(s,z).
\]
We evaluate these three terms in order, keeping the same Mellin
normalization in each.

\subsubsection{The character integral}
\label{proof:v7:character-mellin}

The odd-character family is
$X_N^-=\{\chi\ \mathrm{mod}\ N:\chi(-1)=-1\}$, with all characters extended
by zero at nonunits. The character moment is
\[
 \mathcal P(N)=\frac2{\varphi(N)}
                 \sum_{\chi\in X_N^-}|P_\chi|^2.
\]
Here each character contributes the polynomial
\[
 P_\chi=\sum_{n\ge1}
       \frac{\Lambda(n)\chi(n)}{\sqrt n}\Phi(n/N).
\]
Writing $g_\chi(s)=L'/L(1/2+s,\chi)$, the analytic kernel is
\[
 \mathscr P_N(s,z)=\frac2{\varphi(N)}\sum_{\chi\in X_N^-}
                          g_\chi(s)g_{\bar\chi}(z).
\]
Here $s$ and $z$ are independent holomorphic variables. We use the
Mellin convention
\[
 \widetilde\Phi(s)=\int_0^\infty\Phi(u)u^{s-1}\,\mathrm{d}u,
 \qquad
 \Phi(u)=\frac1{2\pi i}\int_{(\sigma)}\widetilde\Phi(s)u^{-s}\,\mathrm{d}s.
\]
Here $\int_{(\sigma)}$ runs upward along $\Re s=\sigma$; the same
tilde notation and Mellin convention apply to the other weights below.
The decay needed for the contour shifts is
\begin{equation}\label{gb:mellin-decay}
 \widetilde{\Phi}(\sigma+it)\ll_{\phi,I}(1+|t|)^{-6}
 \quad(\sigma\in I),
\end{equation}
for every fixed compact real interval $I$. We prove this estimate
together with the character inversion below.

\begin{lemma}\label{proof:gb:mellin-projection}
The Mellin inversion formula above and the decay estimate
\eqref{gb:mellin-decay} hold, with decay uniform on compact real
intervals for the contour. Under $\mathrm{GRH}(N)$, for every
$0<\sigma\le1/9$ one has the absolutely convergent identities
\begin{align}
 P_\chi&=-\frac1{2\pi i}\int_{(\sigma)}
               N^s\widetilde\Phi(s)g_\chi(s)\,\mathrm{d}s,
               \label{proof:gb:character-inversion}\\
 \mathcal P(N)&=\frac1{(2\pi i)^2}
   \int_{(\sigma)}\int_{(\sigma)}
   N^{s+z}\widetilde\Phi(s)\widetilde\Phi(z)
                 \mathscr P_N(s,z)\,\mathrm{d}s\,\mathrm{d}z.
                 \label{proof:gb:moment-inversion}
\end{align}
\end{lemma}

\begin{proof}
Write the Mellin transform as the Fourier integral of
$e^{\sigma x}\Phi(e^x)$. This function has fixed compact support, and
six integrations by parts give the asserted decay uniformly for
$\sigma$ in a compact real interval. Fourier inversion then gives
the absolutely convergent Mellin inversion formula.

On $\Re s=1$, absolute convergence follows from
$\sum_n\Lambda(n)n^{-3/2}<\infty$. Mellin decay therefore
permits termwise inversion and proves \eqref{proof:gb:character-inversion} on that line.
For an odd character, the inducing primitive $L$-function is entire.
Under GRH its nontrivial zeros lie on $\Re u=1/2$; its trivial zeros
and the zeros of the removed Euler factors lie in $\Re u\le0$.
Thus $L'/L(1/2+s,\chi)$ has no poles between $\sigma$ and $1$.

On the horizontal sides of the contour rectangle, the bound in
Lemma~\ref{proof:lem:zero-expansion} gives
\[
 \left|N^s\widetilde\Phi(s)\frac{L'}L(1/2+s,\chi)\right|
 \ll_{\phi,\sigma}N(1+T)^{-6}\log^2\bigl(N(T+2)\bigr),
 \qquad s=x\pm iT,\quad\sigma\le x\le1.
\]
For fixed $N$ and $\sigma$, these contributions vanish as $T\to\infty$.
The contour is shifted for each fixed modulus and positive left line.
Uniformity in the moving arithmetic boundary is supplied later by the
Hardy estimates, rather than by this shift.
Cauchy's theorem proves the shifted identity. Multiplying the formulas
for $\chi$ and $\bar\chi$, then averaging over $X_N^-$, proves
\eqref{proof:gb:moment-inversion}: $\overline{P_\chi}=P_{\bar\chi}$,
the two minus signs cancel, and absolute convergence permits Fubini.
This proves the lemma.
\end{proof}

\subsubsection{The diagonal and Beta integrals}
\label{proof:sec:goldbach-main}\label{proof:v7:main-term-evaluation}

The diagonal and continuous-model kernels are
\[
 \mathscr D_N(w)=\sum_{\substack{n\ge1\\(n,N)=1}}
                    \frac{\Lambda(n)^2}{n^{1+w}},\qquad
 \mathscr K_N(s,z)=N^{-s-z}\mathscr S_N(s+z)K_-(s,z).
\]
The corresponding finite diagonal is
\[
 \mathcal D(N)=N^{-1}\sum_{(n,N)=1}\Lambda(n)^2\phi(n/N)^2.
\]
The density series and signed Beta factor are
\[
 \begin{aligned}
 \mathscr S_N(w)&=\sum_{k\ge1}\mathfrak S(kN)k^{-w}
                                      \quad(\Re w>1),\\
 K_-(s,z)&=B(1/2-s,s+z)+B(1/2-z,s+z)
                           -B(1/2-s,1/2-z).
 \end{aligned}
\]
Here $B(s,z)=\Gamma(s)\Gamma(z)/\Gamma(s+z)$ is Euler's Beta
function. Lemma~\ref{proof:lem:singular-N} continues $\mathscr S_N$
to $\Re w>0$, with only a simple pole at $w=1$. In the range of
the contours below, $\mathscr K_N$ means the continued model
constructed in Section~\ref{proof:v7:euler}.

To encode the two positive difference families and the negative
sum family, define the following convolutions:
\begin{equation}\label{gb:convolutions}
 \begin{aligned}
 C_0(y)&=\int_0^\infty\phi(u)\phi(u+y)\d u,\qquad
 J_0(y)=\int_0^y\phi(u)\phi(y-u)\d u,\\
 H_0(y)&=2C_0(y)-J_0(y),\qquad J_0(1)=c_\phi.
 \end{aligned}
\end{equation}
The value $J_0(1)=c_\phi$ follows from $\phi(u)=\phi(1-u)$.
Write $\widetilde H_0$ for the Mellin transform of $H_0$.
The Beta--Mellin relation below reduces the model integral to one
Mellin variable, and a contour shift then gives its value.

\begin{lemma}\label{proof:gb:main-term-inversion}
For $0<\sigma\le1/9$, the diagonal and model integrals converge
absolutely and satisfy
\begin{align}
 &\frac1{(2\pi i)^2}\int_{(\sigma)}\int_{(\sigma)}
 N^{s+z}\widetilde\Phi(s)\widetilde\Phi(z)
                         \mathscr D_N(s+z)\,\mathrm{d}s\,\mathrm{d}z=\mathcal D(N),
                         \label{proof:gb:diagonal-inversion}\\
 I_N(\Phi)&:=\frac1{(2\pi i)^2}\int_{(\sigma)}\int_{(\sigma)}
 \widetilde\Phi(s)\widetilde\Phi(z)
                 \mathscr S_N(s+z)K_-(s,z)\,\mathrm{d}s\,\mathrm{d}z
                 =-c_\phi\mathfrak S(N).
                 \label{proof:gb:beta-inversion}
\end{align}
In the model integral, the factor $N^{s+z}$ has cancelled the
$N^{-s-z}$ in $\mathscr K_N$.
\end{lemma}

\begin{proof}
Since $|\mathscr D_N(s+z)|\le\mathscr D_N(2\sigma)$ on the product
contour, Mellin decay permits termwise integration of its Dirichlet
series. Inversion in both variables gives
\[
 \sum_{(n,N)=1}\frac{\Lambda(n)^2}{n}\Phi(n/N)^2
 =\frac1N\sum_{(n,N)=1}\Lambda(n)^2\phi(n/N)^2
 =\mathcal D(N),
\]
which proves \eqref{proof:gb:diagonal-inversion}.

For the model, Euler's integral gives
\begin{equation}\label{proof:gb:inner-beta}
 \begin{aligned}
 \int_0^r u^{-1/2-s}(r-u)^{w-1}\d u
 &=r^{w-1/2-s}B(1/2-s,w),\\
 \int_0^\infty u^{-1/2-s}(u+r)^{w-1}\d u
 &=r^{w-1/2-s}B(1/2-s,1/2+s-w).
 \end{aligned}
\end{equation}
Both identities are absolutely convergent when
$\Re s=\sigma$ and $0<\Re w<1/2+\sigma$.
Insert Mellin inversion for $\phi(u)=u^{-1/2}\Phi(u)$ into the
convolution transforms. The substitutions $r=u+y$ for $C_0$ and
$y=u+r$ for $J_0$, followed by \eqref{proof:gb:inner-beta}, give
\[
 \begin{aligned}
 \widetilde C_0(w)
 &=\frac1{2\pi i}\int_{(\sigma)}
   \widetilde\Phi(s)\widetilde\Phi(w-s)B(1/2-s,w)\d s,\\
 \widetilde J_0(w)
 &=\frac1{2\pi i}\int_{(\sigma)}
   \widetilde\Phi(s)\widetilde\Phi(w-s)
   B(1/2-s,1/2+s-w)\d s.
 \end{aligned}
\]
Here $r$ remains in the fixed support of $\phi$, while the real-part
versions of \eqref{proof:gb:inner-beta} and the integrability of
$\widetilde\Phi$ justify the interchanges. On $\Re w=2\sigma$,
replacing $s$ by $w-s$ gives the second Beta term for $C_0$.
Adding these two expressions and subtracting the expression for $J_0$
gives, for $0<\eta\le1/9$ and $\Re w=2\eta$,
\begin{equation}\label{gb:beta-mellin}
 \widetilde H_0(w)=\frac1{2\pi i}\int_{(\eta)}
 \widetilde{\Phi}(s)\widetilde{\Phi}(w-s)K_-(s,w-s)\d s.
\end{equation}

The support of $\phi$ makes $C_0$ vanish for $y>1/3$ and $J_0$
vanish outside $[2/3,4/3]$. Their derivatives through order six are
bounded, so $\widetilde H_0$ is holomorphic for $\Re w>0$.
For every compact interval $I\subset(0,\infty)$, the derivatives
through order six of $e^{\sigma x}H_0(e^x)$ are integrable uniformly
for $\sigma\in I$: they decay exponentially as $x\to-\infty$ and
vanish for large positive $x$. All boundary terms vanish, so six
integrations by parts give
\begin{equation}\label{proof:gb:H-decay}
 \widetilde H_0(\sigma+it)\ll_{\phi,I}(1+|t|)^{-6}
 \qquad(\sigma\in I,\ t\in\mathbb R).
\end{equation}
Put $a_0=\int_0^\infty\phi(u)\,\mathrm{d}u$. Integration over the positive
quadrant assigns half the product integral to $C_0$ and all of it
to $J_0$, giving
\begin{equation}\label{gb:pole-cancellation}
 \int_0^\infty C_0(y)\d y=\frac{a_0^2}{2},
 \quad
 \int_0^\infty J_0(y)\d y=a_0^2,
 \quad\widetilde H_0(1)=0.
\end{equation}
The support restrictions separately give the integer values
\begin{equation}\label{gb:integer-values}
 C_0(k)=0\ (k\geq1),\quad
 J_0(k)=0\ (k\geq2),\quad H_0(1)=-c_\phi.
\end{equation}

On the original product contour, Euler's integral bounds the three
Beta factors uniformly, while
$\mathscr S_N(2\sigma+it)\ll_{N,\sigma}1+|t|$.
The model integrand is therefore dominated by a constant times
\[
 (1+|v|)^{-6}(1+|y|)^{-6}(1+|v+y|),
\]
an integrable function on $\mathbb R^2$. The change of variables
$(v,y)\mapsto(v,v+y)$ has Jacobian $1$, so Fubini and
\eqref{gb:beta-mellin} give the single-contour expression
\[
 I_N(\Phi)=\frac1{2\pi i}\int_{(2\sigma)}
                  \mathscr S_N(w)\widetilde H_0(w)\,\mathrm{d}w,
 \qquad w=s+z.
\]
This absolute majorant permits the change of variables before either
integration is evaluated. It also shows that the orientation and the
normalization are unchanged when the two heights are replaced by their sum.

The possible pole at $w=1$ is removable in the product
$\mathscr S_N(w)\widetilde H_0(w)$, because $\mathscr S_N$ has a
simple pole there and $\widetilde H_0(1)=0$.
Applying \eqref{proof:gb:singular-series-growth} and
\eqref{proof:gb:H-decay} on $I=[2\sigma,3/2]\subset(0,\infty)$ gives
horizontal integrals
$O_{N,\phi,\sigma}(T^{-5})$. We may therefore move the $w$-contour
to $\Re w=3/2$. On that line the density series
$\mathscr S_N(w)=\sum_{k\ge1}\mathfrak S(kN)k^{-w}$
converges absolutely. Termwise Mellin inversion
therefore completes the evaluation:
\begin{equation}\label{gb:single-contour}
 \begin{aligned}
 I_N(\Phi)&=\frac1{2\pi i}\int_{(2\sigma)}
       \mathscr S_N(w)\widetilde H_0(w)\,\d w\\
       &=\sum_{k\ge1}\mathfrak S(kN)H_0(k)
        =-c_\phi\mathfrak S(N).
 \end{aligned}
\end{equation}
Only $k=1$ survives by \eqref{gb:integer-values}, proving
\eqref{proof:gb:beta-inversion}. The pole cancellation concerns
$\widetilde H_0(1)$; the surviving sampled value is $H_0(1)$.
This proves the lemma.
\end{proof}

\subsubsection{The arithmetic discrepancy}

We can now combine the three evaluations without any further contour
argument. Write the residual integral as
\begin{equation}\label{proof:gb:residual-integral}
 r_\sigma(N)=\frac1{(2\pi i)^2}
   \int_{(\sigma)}\int_{(\sigma)}
   N^{s+z}\widetilde\Phi(s)\widetilde\Phi(z)
                              \mathscr E_N(s,z)\,\mathrm{d}s\,\mathrm{d}z.
\end{equation}
The signs of the three kernels are essential: the model is subtracted
in $\mathscr E_N$, but its own weighted integral is negative.

\begin{proposition}\label{proof:v3w:exact-moving}
Under $\mathrm{GRH}(N)$, the residual integral $r_\sigma(N)$ is
absolutely convergent for $0<\sigma\le1/9$ and satisfies
\begin{equation}\label{proof:gb:exact-error}
 \mathcal G(N)=c_\phi\mathfrak S(N)-r_\sigma(N).
\end{equation}
In particular, $r_\sigma(N)$ is independent of $\sigma$; we write it as
$r(N)$.
\end{proposition}

\begin{proof}
Lemmas~\ref{proof:gb:mellin-projection} and \ref{proof:gb:main-term-inversion}
evaluate the three terms of the residual as
\[
 r_\sigma(N)=\mathcal P(N)-\mathcal D(N)+c_\phi\mathfrak S(N).
\]
The identity $\mathcal G(N)=\mathcal D(N)-\mathcal P(N)$ from
Lemma~\ref{proof:v9:finite-projection} gives
\eqref{proof:gb:exact-error}. Its arithmetic side has no contour parameter.
The model enters with the opposite sign to its own Mellin evaluation
because it is subtracted in the residual kernel.
The convergence constants may depend on $N$ and $\sigma$;
the uniform estimate follows from the Hardy norm below.
This proves the proposition.
\end{proof}

\subsection{The residual as a quadratic form}
\label{proof:detail:hardy-quadratic}

For even $N>e^{10}$ satisfying $\mathrm{GRH}(N)$, we construct a
vector $v_N$ of uniformly bounded norm such that
$r(N)=\langle T_Nv_N,v_N\rangle$. This transfers the spectral
estimates to the arithmetic discrepancy without a factor growing
with $N$. We use the strip and boundary
\[
 \begin{gathered}
 \alpha_N=\frac1{\log N},\qquad \beta=\frac18,\qquad
 \sigma_*=\frac19,\qquad \Delta_N=\beta-\alpha_N,\\
 S_N=\{s:\alpha_N<\Re s<\beta\},\qquad
 \Gamma_N=(\alpha_N+i\mathbb R)\sqcup(\beta+i\mathbb R).
 \end{gathered}
\]
The Hardy space $\mathcal H_N=H^2(S_N)$ consists of holomorphic
functions with bounded paired vertical $L^2$ means as the two lines
approach the boundary. As proved in Lemma~\ref{proof:hd:PW}, its
boundary traces $h_{\alpha_N},h_\beta$ have norm
\[
 \|h\|_{\mathcal H_N}^2
 =\int_{\mathbb R}|h_{\alpha_N}(t)|^2\,\mathrm{d}t
  +\int_{\mathbb R}|h_\beta(t)|^2\,\mathrm{d}t.
\]
Thus both boundary lines contribute to the norm, and the two traces
must come from the same holomorphic function in the strip. The
inner product is
$\langle h,g\rangle=\int_{\Gamma_N}h(s)\overline{g(s)}\,|\mathrm{d}s|$,
linear in its first variable. Both boundaries carry positive arc
length, with no factor $1/(2\pi)$. On this space the compact
selfadjoint operator constructed in Proposition~\ref{proof:bd:HS} is
\[
 (T_Nh)(s)=\int_{\Gamma_N}
     \frac{\mathscr E_N(s,\bar z)}{(1+s)(1+\bar z)}h(z)\,|\mathrm{d}z|.
\]
The next argument restores the weight to this normalized kernel.

\subsubsection{The functional norm}

Define the linear functional
\[
 \ell_N(h)=\frac1{2\pi i}\int_{(\sigma)}
       N^s\widetilde\Phi(s)(1+s)h(s)\,\mathrm{d}s.
\]
Here $\alpha_N<\sigma<\beta$. The bound proved below gives a
unique Riesz vector $v_N\in\mathcal H_N$ characterized by
$\ell_N(h)=\langle h,v_N\rangle$. Write
$\mathcal Q_N=\|\ell_N\|^2=\|v_N\|_{\mathcal H_N}^2$.
For the calculation of this norm, introduce the companion weight
\[
 V(u)=\Phi(u)-u\Phi'(u).
\]
We use the Fourier--Laplace representation
of Lemma~\ref{proof:hd:PW}:
\[
 \begin{aligned}
 (\mathcal U_H\psi)(s)
   &=\frac1{\sqrt{2\pi}}
       \int_{\mathbb R}\psi(\tau)e^{s\tau}\,\mathrm{d}\tau,\\
 \|\mathcal U_H\psi\|_{\mathcal H_N}^2
   &=\int_{\mathbb R}|\psi(\tau)|^2w_{H,N}(\tau)\,\mathrm{d}\tau.
 \end{aligned}
\]
Here $\psi$ ranges over $L^2(\mathbb R,w_{H,N}(\tau)\,\mathrm{d}\tau)$,
with weight
\[
 w_{H,N}(\tau)=e^{2\alpha_N\tau}+e^{2\beta\tau}.
\]
Its two terms correspond to the two boundary lines in the Hardy norm.
This representation identifies
both the norm and the contour independence of $\ell_N$; the factor
$1+s$ in the functional is encoded by $V$.

\begin{proposition}[The functional norm]\label{proof:bm:norm}
The functional $\ell_N$ is continuous and independent of the
interior contour. Its squared norm and uniform bound are
\begin{equation}\label{proof:hardy:functional-norm}
 \mathcal Q_N=\frac1{2\pi}\int_0^\infty
       \frac{(Nu)^{2\alpha_N}|V(u)|^2}
            {1+(Nu)^{-2\Delta_N}}\,\frac{\mathrm{d}u}{u}
 \le\frac{e^2}{2\pi}\int_0^\infty|V(u)|^2\,\frac{\mathrm{d}u}{u}.
\end{equation}
Moreover,
\begin{equation}\label{bm:norm-asymptotic}
 \mathcal Q_N
 =\frac{e^2}{2\pi}\int_0^\infty|V(u)|^2\,\frac{\mathrm{d}u}{u}
       +O_{\phi}\left(\frac1{\log N}\right).
\end{equation}
\end{proposition}

\begin{proof}
Integration by parts gives
$\widetilde V(s)=(1+s)\widetilde\Phi(s)$; the compact support of
$\Phi$ removes the boundary term.
Apply the Fourier--Laplace representation
\eqref{hd:PW-map}--\eqref{hd:boundary-norm}, first to a density
$\psi$ of compact support. Fubini and Mellin inversion give
\[
 \begin{aligned}
 \ell_N(\mathcal U_H\psi)
 &=\frac1{\sqrt{2\pi}}\int_{\mathbb R}\psi(\tau)
   \left\{\frac1{2\pi i}\int_{(\sigma)}
        \widetilde V(s)(Ne^\tau)^s\,\mathrm{d}s\right\}\,\mathrm{d}\tau\\
 &=\frac1{\sqrt{2\pi}}\int_{\mathbb R}
        \psi(\tau)V(e^{-\tau}/N)\,\mathrm{d}\tau.
 \end{aligned}
\]
The first interchange is absolute because $\widetilde V$ is integrable
on the contour and $\psi$ has compact support. In $L^2(w_{H,N})$,
the last integral has squared norm
\[
 \frac1{2\pi}\int_{\mathbb R}
        \frac{|V(e^{-\tau}/N)|^2}{w_{H,N}(\tau)}\,\mathrm{d}\tau.
\]
Weighted Cauchy--Schwarz gives the upper bound, with equality for a
density proportional to
$\overline{V(e^{-\tau}/N)}/w_{H,N}(\tau)$. This density has compact
support for each fixed $N$, so the norm is exact.
The weight is strictly positive, so this equality case is admissible
in the weighted space. Its norm is finite because its support is compact.

Compactly supported densities are dense in the weighted space.
The original contour integral is also continuous there:
$e^{2\sigma\tau}\le w_{H,N}(\tau)$ bounds restriction to the
interior line in $L^2$, and $(1+s)\widetilde\Phi(s)$ is square
integrable. The two expressions therefore agree for every Hardy
vector. The density expression proves contour independence.

The substitution $\tau=-\log(Nu)$ gives the $u$-integral in
\eqref{proof:hardy:functional-norm}. On the fixed support of $V$,
\[
 (Nu)^{2\alpha_N}
 =e^2\{1+O_\phi((\log N)^{-1})\},\qquad
 (Nu)^{-2\Delta_N}=O_\phi(N^{-1/4}),
\]
which proves \eqref{bm:norm-asymptotic}. Finally,
$(Nu)^{2\alpha_N}=e^2u^{2/\log N}\le e^2$ and the denominator in
\eqref{proof:hardy:functional-norm} is at least $1$ on this support.
This proves the uniform bound in~\eqref{proof:hardy:functional-norm},
hence $\mathcal Q_N\ll_\phi1$.
This proves the proposition.
\end{proof}

\subsubsection{Recovering the quadratic form}

Proposition~\ref{proof:bd:evaluation} gives the evaluation vectors
\[
 k_z(s)=\frac1{4\Delta_N}
       \csc\!\left(\frac{\pi(s+\bar z-2\alpha_N)}{2\Delta_N}\right),
 \qquad h(s)=\langle h,k_s\rangle.
\]
It also identifies the residual with the matrix elements of $T_N$:
\[
 \langle T_Nk_z,k_s\rangle
       =\frac{\mathscr E_N(s,\bar z)}{(1+s)(1+\bar z)}.
\]
Thus applying the weight functional to the evaluation vectors will
restore exactly the two Mellin factors. The conjugations are fixed
by the convention for the inner product.

\begin{theorem}[The exact residual identity]\label{proof:bm:quadratic}
For every even $N>e^{10}$ satisfying $\mathrm{GRH}(N)$, the residual
satisfies
\begin{equation}\label{proof:hardy:quadratic-bound}
 r(N)=\langle T_Nv_N,v_N\rangle,
 \qquad |r(N)|\le\mathcal Q_N\|T_N\|\ll_\phi\|T_N\|.
\end{equation}
Consequently,
$\mathcal G(N)=c_\phi\mathfrak S(N)-\langle T_Nv_N,v_N\rangle$, and
\begin{equation}\label{proof:hardy:goldbach-error}
 |\mathcal G(N)-c_\phi\mathfrak S(N)|\ll_\phi\|T_N\|.
\end{equation}
\end{theorem}

\begin{proof}
Put $\Psi_N(s)=N^s\widetilde\Phi(s)(1+s)$ and fix
$\sigma=\sigma_*=1/9$. The reproducing property and uniqueness of
the Riesz vector give
\[
 v_N=\frac1{2\pi}\int_{\mathbb R}
       \overline{\Psi_N(\sigma+iy)}k_{\sigma+iy}\,\mathrm{d}y.
\]
This is a Bochner integral, an integral of Hilbert-space vectors
defined by convergence in the Hilbert norm. In the present setting
it is the norm limit of the integrals over bounded intervals;
integrability of the vector norms ensures that this limit exists.
Indeed, \eqref{proof:hd:evaluation-norm} gives
$\|k_s\|^2\le[8\min(\Re s-\alpha_N,\beta-\Re s)]^{-1}$.
At $\Re s=\sigma_*$ both distances are at least $1/90$, so
$\sup_y\|k_{\sigma_*+iy}\|^2\le45/4$, while $\Psi_N$ is integrable
by Mellin decay. The coefficient is conjugated because the inner product
is conjugate-linear in its second entry.
When expanding the quadratic form, one may first perform both integrations
over bounded intervals.
The evaluation bound then permits passage to the full lines by dominated
convergence, preserving the order of the two inner-product entries.

The kernel identity established in the proof of
Proposition~\ref{proof:bd:evaluation} yields
\[
 \langle T_Nv_N,v_N\rangle
 =\frac1{(2\pi)^2}\int_{\mathbb R^2}
 \Psi_N(\sigma+ix)\overline{\Psi_N(\sigma+iy)}
 \frac{\mathscr E_N(\sigma+ix,\sigma-iy)}
      {(1+\sigma+ix)(1+\sigma-iy)}\,\mathrm{d}x\,\mathrm{d}y.
\]
The same evaluation bound gives absolute convergence. Since $\Phi$
is real, $\overline{\Psi_N(\sigma+iy)}=\Psi_N(\sigma-iy)$.
Replace $y$ by $-y$; reversing the limits cancels the sign from $\mathrm{d}y$.
Both Mellin variables now run upward, and their factors $1+s$ and
$1+z$ cancel the kernel denominators. The remaining integral is
exactly the residual integral in~\eqref{proof:gb:residual-integral}.
This proves the identity in~\eqref{proof:hardy:quadratic-bound}.

Cauchy--Schwarz and Proposition~\ref{proof:bm:norm} give
$|r(N)|\le\|T_N\|\,\|v_N\|^2=\mathcal Q_N\|T_N\|$.
Combining this with \eqref{proof:gb:exact-error} gives
\eqref{proof:hardy:goldbach-error}, the estimate used for
Theorems~\ref{thm:goldbach-S} and \ref{thm:goldbach-R}.
This proves the theorem.
\end{proof}

\subsection{Proofs of the main theorems}
\label{proof:main-theorem-proofs}
\label{proof:v7:siegel-hypothesis}\label{proof:v7:riemann-hypothesis}

We now combine the spectral bounds of Theorem~\ref{proof:qp:dictionary}
with the Hardy estimate of Theorem~\ref{proof:bm:quadratic}.
The following lemma expresses the relative Goldbach error in terms
of a zero-free width and gives a criterion for the asymptotic formula.

\begin{lemma}[An arithmetic estimate from a zero-free width]
\label{proof:v9:width-arithmetic-estimate}
Assume $\mathrm{GRH}(N)$ for an even $N>e^{10}$.
If $0<\eta_N\le1/4$ and $\mathfrak X_N(q)\ne0$ for
$\Re q>1-\eta_N$, then
\begin{equation}\label{proof:width:relative-error}
 \frac{|\mathcal G(N)-c_\phi\mathfrak S(N)|}{c_\phi\mathfrak S(N)}
 \ll_\phi\sqrt{\frac{(\log N)^3}{N\eta_N\mathfrak S(N)}}.
\end{equation}
As $N\to\infty$ through even integers satisfying these hypotheses,
the relative Goldbach asymptotic follows from the condition
\begin{equation}\label{proof:width:asymptotic-condition}
 \frac{N\eta_N\mathfrak S(N)}{(\log N)^3}\longrightarrow\infty.
\end{equation}
\end{lemma}

\begin{proof}
Insert the bound
$\|T_N\|\le\delta_N/(2\sqrt{\eta_N(1-\eta_N)})$
from Theorem~\ref{proof:qp:dictionary} into the Hardy estimate in
Theorem~\ref{proof:bm:quadratic} and divide by
$c_\phi\mathfrak S(N)$.
Using $\delta_N^2=\mathfrak S(N)(\log N)^3/N$ and
$1-\eta_N\ge3/4$ gives \eqref{proof:width:relative-error}.
Under \eqref{proof:width:asymptotic-condition}, the right-hand side
of \eqref{proof:width:relative-error} tends to zero.
This proves the lemma.
\end{proof}

We now prove Theorems~\ref{thm:goldbach-S} and \ref{thm:goldbach-R}
under Hypotheses S and R, respectively.

\begin{proof}[Proof of Theorem~\ref{thm:goldbach-S}]
Hypothesis S excludes zeros in $\Re q>1-\eta_N$, with the
prescribed width
\[
 \eta_N=\frac{h(N)^2(\log N)^3}{N\mathfrak S(N)}\le\frac14,
 \qquad h(N)\longrightarrow\infty,
\]
as in~\eqref{v3i:S-width}.
Here $h$ is fixed in advance. Substitution in
Theorem~\ref{proof:qp:dictionary} gives
\[
 \|T_N\|\le\frac{\mathfrak S(N)}{2h(N)\sqrt{1-\eta_N}}
          \le\frac{\mathfrak S(N)}{\sqrt3\,h(N)}.
\]
The estimate in Theorem~\ref{proof:bm:quadratic} now gives
\[
 \mathcal G(N)=c_\phi\mathfrak S(N)
                   +O_\phi\!\left(\frac{\mathfrak S(N)}{h(N)}\right),
\]
as in~\eqref{v3i:weak-result}.
Its implied constant depends only on $\phi$; the required lower
bound on $N$ may depend on the admissibility and growth of $h$.
Since $h(N)\to\infty$, the error is $o(\mathfrak S(N))$;
if $h(N)/\mathfrak S(N)\to\infty$, it is also $o(1)$.
This proves the theorem.
\end{proof}

\begin{proof}[Proof of Theorem~\ref{thm:goldbach-R}]
Hypothesis R puts every spectral zero on $\Re q=1/2$.
By Theorem~\ref{proof:qp:dictionary}, this is equivalent to
$\|T_N\|\le\delta_N$. The estimate in
Theorem~\ref{proof:bm:quadratic} yields
\[
 \mathcal G(N)=c_\phi\mathfrak S(N)
 +O_\phi\!\left(\sqrt{\frac{\mathfrak S(N)}N}(\log N)^{3/2}\right),
\]
as in~\eqref{v3i:strong-result}.
The zero classification also shows that R implies S for every
prescribed admissible width: all exceptional real pairs have been
excluded. The functional symmetries alone provide no such norm bound.
This proves the theorem.
\end{proof}

The next section places these results alongside the twin-prime
problem and discusses the possible scope of classical methods.

\section{Goldbach, twin primes, and classical methods}
\label{story:philosophy}\label{proof:v7:philosophy}

Goldbach's conjecture and the twin-prime conjecture both concern
the simultaneous primality of two linear expressions, but their
arithmetic questions are different. The twin-prime conjecture asks
for infinitely many prime pairs $(p,p+2)$, with the difference fixed.
Goldbach's conjecture asks for a prime pair $(p,N-p)$ for every even
$N\ge4$; its asymptotic version requires uniform control as the
target $N$ varies.

The twin-prime problem fits the general orbit-and-polynomial
formulation of the affine sieve of Bourgain, Gamburd, and
Sarnak~\cite{BGS2010}. For example, let
\[
 U=\begin{pmatrix}1&1\\0&1\end{pmatrix},\qquad
 v=\binom01,\qquad F(x,y)=x(x+2y).
\]
Then $\langle U\rangle v=\{(n,1)^{\mathsf T}:n\in\mathbb Z\}$ and
$F(U^nv)=n(n+2)$. The desired prime pairs are therefore the orbit
points with $n>1$ at which both linear factors are prime.
This is a unipotent orbit; the expansion theorems for semisimple
groups do not apply to it directly.

For Goldbach, the corresponding polynomial is $F_N(n)=n(N-n)$,
and the relevant positive range is finite for each fixed $N$.
The polynomial, the range, and the local congruence data all vary
with $N$. An orbit description is still possible, but obtaining
representations for every sufficiently large even target requires
estimates uniform in this varying family. This distinction matters
when comparing the two conjectures through a common sieve framework.

We do not know whether sufficiently ingenious refinements of sieve
methods or the circle method can yield the Goldbach asymptotic
formula. The present difficulties do not warrant excluding this
possibility.

\begin{samepage}
The central advance of this paper is the exact identification of the
signed Goldbach error with a quadratic form of the residual operator.
For every even $N>e^{10}$ satisfying $\mathrm{GRH}(N)$,
Theorem~\ref{proof:bm:quadratic} gives
\[
 c_\phi\mathfrak S(N)-\mathcal G(N)
 =\langle T_Nv_N,v_N\rangle.
\]
\end{samepage}
The operator $T_N$ is compact and selfadjoint, and the vector $v_N$
determined by the weight satisfies $\|v_N\|^2\ll_\phi1$.
The identity thus retains the full signed arithmetic discrepancy and
transfers operator norm bounds to error bounds without a factor
growing with $N$. Hypotheses S and R enter only after this identity
has been established: they control $\|T_N\|$ and hence the size of
the inner product. They yield, respectively, the relative asymptotic
formula of Theorem~\ref{thm:goldbach-S} and the sharper remainder of
Theorem~\ref{thm:goldbach-R}.

Nor can we rule out that sieve methods might produce a positive
lower bound sufficient to establish Goldbach representations,
even without the full asymptotic formula. In pursuing such lower
bounds, we have found some remarkable phenomena, whose analysis
will be presented elsewhere.

\medskip
\noindent
{\bf Acknowledgements}. The authors are supported by the National Key
Research and Development Program of China (No. 2021YFA1000700), the
National Natural Science Foundation of China (Nos. 12631001 \& 123B2001),
and the Key R\&D Program of Shandong Province, China (No. 2026ZLYS01).
\medskip

\clearpage
\appendix

\section{Dependence on the strip parameter}\label{proof:app:param}

The main text fixes $\alpha_N=1/\log N$. To discuss that choice,
we introduce the family $\alpha_{N,c}=c/\log N$, so the main
normalization is $c=1$. We track the dependence of the Hardy
realization and spectral conditions on $c$, then distinguish
fixed-$c$ asymptotics from the limit $c\downarrow0$ at fixed $N$.

\subsection{The Hardy realization with a variable boundary}
\label{proof:app:param:realization}

We construct the Hardy realization as the left boundary varies
and identify its unchanged arithmetic pairing. The arithmetic
weights and residual kernel remain those of the main text, while
the Hardy space, operator and representing vector acquire $c$.

Fix $c>0$ and an even $N\ge4$ with $N>e^{10c}$. Put
\begin{equation}\label{proof:app:param:strip}
 \alpha_{N,c}=\frac{c}{\log N},\qquad
 \Delta_{N,c}=\beta-\alpha_{N,c}.
\end{equation}
Use \eqref{hd:strip}, \eqref{hd:space-and-boundary}, \eqref{qp:operator} and
\eqref{bm:functional} with this left boundary to obtain
$S_{N,c}$, $\mathcal H_{N,c}$, $\Gamma_{N,c}$, $T_{N,c}$ and
$\ell_{N,c}$; write $v_{N,c}$ for the Riesz vector and
$\mathcal Q_{N,c}=\|v_{N,c}\|^2$.

\begin{proposition}[Hardy realization with a variable boundary]
\label{proof:v9:variable-boundary-realization}
Assume $\mathrm{GRH}(N)$. Then $\alpha_{N,c}<1/10$,
$\Delta_{N,c}>1/40$, and $T_{N,c}$ is a compact selfadjoint
Hilbert--Schmidt operator on $\mathcal H_{N,c}$.
The common contour $\sigma_*=1/9$ remains at distance greater
than $1/90$ from the left boundary and exactly $1/72$ from the right.
The point-evaluation estimate \eqref{proof:bd:pointwise}
holds with the distance $d_{N,c}$ to the new boundary.
With $V$ as in Subsubsection~\ref{story:hardy-quadratic},
the functional norm satisfies
\begin{equation}\label{proof:app:param:functional-norm}
 \begin{aligned}
 \mathcal Q_{N,c}
 &=\frac{e^{2c}}{2\pi}\int_0^\infty |V(u)|^2\,\frac{\mathrm{d}u}{u}
       +O_{\phi,c}\!\left(\frac1{\log N}\right),\\
 \mathcal Q_{N,c}
 &\le\frac{e^{2c}}{2\pi}\int_0^\infty |V(u)|^2\,\frac{\mathrm{d}u}{u}.
 \end{aligned}
\end{equation}
The first line holds as $N\to\infty$ with $c$ fixed; the second
holds throughout the permitted range. Moreover,
\begin{equation}\label{proof:app:param:quadratic-form}
 r(N)=\langle T_{N,c}v_{N,c},v_{N,c}\rangle,\qquad
 |r(N)|\le\mathcal Q_{N,c}\|T_{N,c}\|.
\end{equation}
The left side is independent of $c$: varying the realization
changes the operator and vector, but not their arithmetic pairing.
At $c=1$ all these objects reduce to those of the main text.
\end{proposition}

\begin{proof}
The inequalities for the strip and contour follow from
$\alpha_{N,c}=c/\log N<1/10$, $\beta=1/8$ and $\sigma_*=1/9$.
The proof of Proposition~\ref{proof:bd:HS} applies on every such closed
strip, and the Fourier--Laplace and point-evaluation arguments
apply to arbitrary strip endpoints. This proves the operator
properties and the stated pointwise estimate.

In \eqref{bm:norm-integral}, $N^{2\alpha_{N,c}}=e^{2c}$.
On the fixed support of $V$ in $(0,1)$, the denominator is at
least one and $u^{2\alpha_{N,c}}\le1$, giving the upper bound
in \eqref{proof:app:param:functional-norm}. For fixed $c$,
$u^{2\alpha_{N,c}}=1+O_{\phi,c}(1/\log N)$ uniformly on this
support, while the remaining denominator tends uniformly to one
with a power saving in $N$. This gives its asymptotic formula.
The proof of Theorem~\ref{proof:bm:quadratic}, on the same fixed
interior contour, gives \eqref{proof:app:param:quadratic-form}.
Its arithmetic side depends only on the residual and the weight,
so it is independent of $c$. The specialization at $c=1$ follows
from the definitions.
This proves the proposition.
\end{proof}

\subsection{Spectral conditions and their quantifiers}
\label{proof:app:param:families}

We establish the zero--norm correspondence for each realization
and state its asymptotic consequences with $c$ fixed before $N$
grows. This order of quantifiers is essential for the estimates.

Let $\mathfrak X_{N,c}$ be the product
\eqref{qp:spectral-function} formed from $T_{N,c}$, retaining the
scale \eqref{v3i:spectral-scale}, which is independent of $c$.
Use $\mathrm R(N,c)$ and $\mathrm S(N,c;h_c)$ for the corresponding
conditions of Section~\ref{story:section1}.
For each fixed $c>0$, preassign a positive $h_c$ on the even
integers, tending to infinity and satisfying
\begin{equation}\label{proof:app:param:S-width}
 \eta_{N,c}=\frac{h_c^2(N)(\log N)^3}{N\mathfrak S(N)}\le\frac14
\end{equation}
eventually.

\begin{proposition}[Spectral conditions for a fixed boundary parameter]
\label{proof:v9:parameter-spectral-conditions}
Assume $\mathrm{GRH}(N)$ and $N>e^{10c}$. Then
\begin{equation}\label{proof:app:param:dictionary}
 \begin{gathered}
 \mathrm R(N,c)\quad\Longleftrightarrow\quad\|T_{N,c}\|\le\delta_N,\\
 \mathfrak X_{N,c}(q)\ne0\quad(\Re q>1-\eta)
 \quad\Longleftrightarrow\quad
 \|T_{N,c}\|\le\frac{\delta_N}{2\sqrt{\eta(1-\eta)}} ,
 \qquad 0<\eta<\tfrac12.
 \end{gathered}
\end{equation}
The half-space is open, so zeros on its boundary are allowed.
For every admissible $\eta_{N,c}$,
\begin{equation}\label{proof:app:param:S-norm}
 \mathrm S(N,c;h_c)\quad\Longleftrightarrow\quad
 \|T_{N,c}\|\le
 \frac{\mathfrak S(N)}{2h_c(N)\sqrt{1-\eta_{N,c}}}.
\end{equation}
For fixed $c>0$ and sufficiently large even $N$ satisfying GRH,
\[
 \|T_{N,c}\|\ll_c(\log N)^2,\qquad
 \eta_{N,c}^{\mathrm{GRH}}\asymp_c\frac{\mathfrak S(N)}{N\log N},
\]
and
\[
 \mathcal G(N)=c_\phi\mathfrak S(N)+
 \begin{cases}
 O_{\phi,c}\!\left(\mathfrak S(N)/h_c(N)\right),&\mathrm S(N,c;h_c),\\
 O_{\phi,c}\!\left(\delta_N\right),&\mathrm R(N,c).
 \end{cases}
\]
GRH is assumed in both cases. In the S estimate the implied
constant is independent of $h_c$, although the required lower
bound on $N$ depends on the growth and admissibility of $h_c$.
\end{proposition}

\begin{proof}
The proof of Theorem~\ref{proof:qp:dictionary} uses only compactness,
selfadjointness and the Hilbert--Schmidt property, supplied here
by Proposition~\ref{proof:v9:variable-boundary-realization}.
It therefore proves \eqref{proof:app:param:dictionary} for
$\mathfrak X_{N,c}$. Substituting \eqref{proof:app:param:S-width}
and the scale \eqref{v3i:spectral-scale} gives
\eqref{proof:app:param:S-norm}.

In the component estimates of
Proposition~\ref{proof:grh:component-bounds}, replacing
$1/\log N$ by $c/\log N$ changes the constants to ones depending
on the fixed $c$ and gives $\|T_{N,c}\|\ll_c(\log N)^2$.
The zero--norm correspondence then supplies the stated GRH width.
Combining \eqref{proof:app:param:quadratic-form} with the upper
bound in \eqref{proof:app:param:functional-norm} gives the two
arithmetic estimates. For S, use $1-\eta_{N,c}\ge3/4$;
the functional bound is independent of $h_c$. For R, use
$\|T_{N,c}\|\le\delta_N$.
This proves the proposition.
\end{proof}

The eventual hypotheses assert these conditions for all sufficiently
large even $N$ satisfying $\mathrm{GRH}(N)$, with lower bounds
$N_0(c)$ and $N_0(c,h_c)$ on $N$. Neither these lower bounds nor
the admissibility and divergence of $h_c$ need be uniform in $c$.
The main hypotheses are their specializations at $c=1$.
These conclusions concern sufficiently large $N$; covering every
even $N\ge4$ requires a separate check of the finite initial range.

\subsection{Why the boundary has order \texorpdfstring{$1/\log N$}{1/log N}}
\label{proof:app:param:balance}\label{proof:bm:boundary-balance}

The comparison in Subsection~6.5 balances the
functional norm against the available operator bound. Here we
justify its uniformity and record the dependence on $c$.

Use $\widehat T_{N,\alpha}$, $\mathcal Q_\alpha$ and $M_\alpha$
from that comparison, with $0<\alpha<1/10$.

\begin{lemma}[Balance of the functional and operator bounds]
\label{proof:v9:boundary-balance}
For sufficiently large even $N$ satisfying $\mathrm{GRH}(N)$,
uniformly for $0<\alpha<1/10$,
\[
 \mathcal Q_\alpha\asymp_\phi N^{2\alpha},\qquad
 \|\widehat T_{N,\alpha}\|\le
 M_\alpha=C((\log N)^2+\alpha^{-2}),
\]
where $C$ is absolute. Setting $\alpha=c/\log N$ gives
\[
 \mathcal Q_\alpha M_\alpha
 \asymp_\phi(\log N)^2e^{2c}(1+c^{-2}).
\]
The factor $e^{2c}(1+c^{-2})$ tends to infinity as $c\downarrow0$
and as $c\to\infty$.
The ratio of the
sufficient width for error $O(\mathfrak S(N)/h_c(N))$ to the
width supplied by $M_\alpha$ is, when both widths are admissible,
\[
 \asymp_\phi
 \frac{h_c^2(N)(\log N)^4}{\mathfrak S^2(N)}
       e^{4c}(1+c^{-2})^2.
\]
\end{lemma}

\begin{proof}
The denominator of \eqref{bm:norm-integral} lies between one and
two on the fixed support for large $N$, and $u^{2\alpha}$ is
bounded above and below there, uniformly in $\alpha$.
This proves $\mathcal Q_\alpha\asymp_\phi N^{2\alpha}$.
The component estimates with left boundary $\alpha$ give
$M_\alpha=C((\log N)^2+\alpha^{-2})$ with an absolute $C$.
Substitution of $\alpha=c/\log N$ gives the stated product.
As $c\downarrow0$, its factor $e^{2c}(1+c^{-2})$ is at least
$c^{-2}$; as $c\to\infty$, it is at least $e^{2c}$.
This proves divergence at both ends.

For widths at most $1/4$, Theorem~\ref{proof:qp:dictionary}
makes a width comparable to $\delta_N^2/M^2$ equivalent to
a norm bound of size $M$. The stated arithmetic error is ensured
by a norm bound of size $\mathfrak S(N)/(h_c(N)\mathcal Q_\alpha)$
by the Cauchy--Schwarz argument in
\eqref{v3f:master-error-formula} with left boundary $\alpha$,
retaining the factor $\mathcal Q_\alpha$.
Dividing this sufficient width by the width supplied by $M_\alpha$
therefore gives
\[
 \asymp_\phi
 \frac{h_c^2(N)\mathcal Q_\alpha^2M_\alpha^2}
      {\mathfrak S^2(N)},
\]
which is the asserted ratio.
This proves the lemma.
\end{proof}

Lemma~\ref{proof:v9:boundary-balance} shows that a fixed positive
$c$ gives the best order
available from these bounds; they do not select a universal
optimal constant.
Changing $c$ therefore changes the constants, not the logarithmic
gap in the available estimates.

\subsection{The obstruction at a fixed modulus}
\label{proof:app:param:fixed-modulus}\label{proof:sec:small-c-obstruction}

We examine the different limit in which $N$ is fixed and
$c\downarrow0$. A pole of a logarithmic derivative supplies
a lower bound for the residual operator norm and obstructs
a zero-free width uniform in $c$.

\begin{proposition}[The obstruction at a fixed modulus]
\label{proof:v9:fixed-modulus-obstruction}
Assume $\mathrm{GRH}(N)$ for a fixed even $N\ge4$.
Choose a nontrivial zero $\rho=1/2+i\gamma$ of an odd character
modulo $N$, and let $m_\chi(\rho)$ be its multiplicity for
$L(s,\chi)$, with value zero when it is not a zero. Put
\[
 C_{N,\gamma}=\frac2{\varphi(N)}
     \sum_{\chi\in X_N^-}m_\chi^2(\rho)>0.
\]
Then
\[
 \liminf_{c\downarrow0}c\|T_{N,c}\|
 \ge\frac{2C_{N,\gamma}\log N}{1+\gamma^2}>0.
\]
In particular, there are $z_N^{\mathrm{obs}}>0$ and
$0<c_N^{\mathrm{obs}}\le\log N/24$ such that
\begin{equation}\label{proof:app:param:obstruction}
 \|T_{N,c}\|\ge z_N^{\mathrm{obs}}/c
 \qquad(0<c<c_N^{\mathrm{obs}}).
\end{equation}
The condition $\mathrm R(N,c)$ fails for all sufficiently small
$c$. In the same range, every zero-free width $0<\eta\le1/4$
satisfies
\[
 \eta\le\frac{\delta_N^2c^2}{3(z_N^{\mathrm{obs}})^2}.
\]
\end{proposition}

\begin{proof}
At $s=\sigma+i\gamma$, Laurent expansion gives
\begin{equation}\label{proof:qp:small-c-correlation}
 \mathscr P_N(s,\bar s)=\frac{C_{N,\gamma}}{\sigma^2}
                         +O_{N,\gamma}(\sigma^{-1}).
\end{equation}
The difference between $\mathscr D_N(w)$ and
$(\zeta'/\zeta)'(1+w)$ consists of higher prime powers and
corrections at primes dividing $N$, and is analytic near zero. Hence
\begin{equation}\label{proof:qp:small-c-diagonal}
 \mathscr D_N(2\sigma)=\frac1{4\sigma^2}+O_N(1).
\end{equation}
For the model, factor $P_N(w)=\zeta(1+w)R_N(w)$, where
\[
 R_N(w)=\prod_{p\mid N}(1-p^{-1-w})
 \prod_{\substack{p>2\\p\nmid N}}
 \left(1+\frac{2p^{-w}-p^{-2w}}{p(p-2)}\right).
\]
This product is analytic near zero, and its Euler factors give
$\mathfrak S(N)R_N(0)=1$. Consequently
$\mathfrak S(N)P_N(2\sigma)=1/(2\sigma)+O_N(1)$.
Using $|\mathcal X_1(i\gamma)|=1$ and
$\zeta(1-2\sigma)=-1/(2\sigma)+O(1)$ in
\eqref{v3:model-functional-form}, we obtain
\begin{equation}\label{proof:qp:small-c-model}
 \mathscr K_N(s,\bar s)=-\frac1{4\sigma^2}
                         +O_{N,\gamma}(\sigma^{-1}).
\end{equation}
The quadratic singularities of the diagonal and model cancel, so
\[
 \mathscr E_N(s,\bar s)=\frac{C_{N,\gamma}}{\sigma^2}
                         +O_{N,\gamma}(\sigma^{-1}).
\]

Take $\sigma=2c/\log N=2\alpha_{N,c}$. For $c<\log N/24$,
the distance to the strip boundary is $\sigma/2$.
The general-strip version of \eqref{proof:bd:pointwise} gives
\[
 \|T_{N,c}\|\ge
 \frac{4\sigma|\mathscr E_N(s,\bar s)|}{|1+s|^2}.
\]
Since $c/\sigma=\log N/2$, the residual expansion proves the
stated lower limit and hence \eqref{proof:app:param:obstruction}.
Since $\delta_N$ is fixed, $\mathrm R(N,c)$ fails for sufficiently
small $c$. Likewise, a zero-free width $0<\eta\le1/4$ must satisfy
\[
 \frac{z_N^{\mathrm{obs}}}{c}
 \le\frac{\delta_N}{2\sqrt{\eta(1-\eta)}},
 \qquad
 \eta\le\frac{\delta_N^2c^2}{3(z_N^{\mathrm{obs}})^2}.
\]
Here the last inequality uses $1-\eta\ge3/4$.
This proves the proposition.
\end{proof}

No positive width is uniform in this fixed-modulus limit.
This does not decide the eventual validity of R or S with $c$
fixed before $N$ tends to infinity.

\section{A numerical experiment at \texorpdfstring{$N=30030$}{N=30030}}
\label{proof:app:numerical}

We recover the number of unordered Goldbach representations of $30030$
from character integrals, exact finite corrections, and an error interval
containing a unique integer, and compare the calculation with the classical
Hardy--Littlewood integral prediction. This route uses finite prime and
prime-power data for corrections and error checks, without requiring an
enumeration of all Goldbach pairs. The count is unconditional; its Hardy
interpretation and the finite spectral picture assume $\mathrm{GRH}(N)$.

Throughout, $N=30030=2\cdot3\cdot5\cdot7\cdot11\cdot13$ and
$\varphi(N)=5760$. The counting moments use all $5760$ characters;
the Hardy residual uses the $2880$ odd characters.

\subsection{A wide weight and exact finite corrections}
\label{proof:app:num:weight}

A wider weight than $\phi$ is needed to count all representations.
Reciprocal logarithms remove the von Mangoldt weights, with finite
corrections for proper prime powers and the cutoff near the endpoints.
Set $A=16$ and $L=1000$. We use the cutoff
\[
 \eta(x)=P\!\left(\frac{x-A}{L-A}\right)
          P\!\left(\frac{N-x-A}{L-A}\right),
\]
where $P(t)=0$ for $t\le0$, $P(t)=1$ for $t\ge1$, and
\[
 P(t)=\sum_{j=7}^{13}\binom{13}{j}t^j(1-t)^{13-j}
 \qquad(0<t<1).
\]
With zero extension off the support, put
\[
 W(u)=\frac{\sqrt{1-u}\,\eta(Nu)}
                {\sqrt{\log(Nu)\log(N(1-u))}},\qquad f(u)=\frac{W(u)}{\sqrt u},
\]
and define
\[
 \mathcal J=\frac1N\sum_{\substack{n+m=N\\(nm,N)=1}}
        \Lambda(n)\Lambda(m)f(n/N)f(m/N).
\]
Let $r_G(N)$ count unordered prime pairs. Write
$\mathcal E=\sum_{p\le q,\ p+q=N,\ p<L}(1-\eta(p)^2)$, with $p,q$ prime,
and let $\mathcal C_{\rm pp}$ be the ordered sum of
$\eta(p^k)^2/(k\ell)$ over $p^k+q^\ell=N$, $(pq,N)=1$, $k,\ell\ge1$,
$\max(k,\ell)\ge2$.

\begin{proposition}[The wide-weight counting identity]
\label{proof:v9:wide-weight-count}
The cutoff $\eta$ is $C^6$, symmetric about $N/2$, and equal to $1$
on $[L,N-L]$. On its support,
\[
 f(u)f(1-u)=\frac{\eta(Nu)^2}{\log(Nu)\log(N(1-u))}.
\]
At $N=30030$, the exact counting identity is
\begin{equation}\label{proof:app:num:count}
 r_G(N)=\frac N2\mathcal J-\frac{\mathcal C_{\rm pp}}2+\mathcal E.
\end{equation}
\end{proposition}

\begin{proof}
Since $P'(t)=12012t^6(1-t)^6$, its first six derivatives match the
constant extensions. The cutoff properties and the formula for $f$ follow.
A prime-power pair contributes $\eta(p^k)^2/(k\ell)$ to $N\mathcal J$;
subtracting $\mathcal C_{\rm pp}$ leaves the weighted ordered prime pairs.
Every prime pair is coprime to $N$: otherwise a common prime divisor
would force $N=2p$. Also, $N/2=15015$ is not a prime power.
Division by two therefore counts unordered pairs, and $\mathcal E$
restores the missing edge weights.
This proves the proposition.
\end{proof}

Exact rational evaluation of the $75$ edge pairs and $26$ unordered
proper-prime-power pairs gives
\[
 \mathcal E=44.35776816334822525\ldots,\qquad
 \mathcal C_{\rm pp}/2=10.16758713116445818\ldots.
\]

\subsection{Character integrals and their numerical evaluation}
\label{proof:app:num:method}

Using the Mellin convention \eqref{gb:mellin-convention}, we
evaluate the character integrals
\[
 I_\chi=\frac1{2\pi}\int_{\mathbb R}N^{1+it}\widetilde W(1+it)
       \frac{L'}L(3/2+it,\chi)\,\mathrm{d}t.
\]
We also use the diagonal moment
\[
 \mathcal D=\sum_{\substack{1\le n<N\\(n,N)=1}}a_n^2,
\]
where the coefficients are
\[
 a_n=\frac{\Lambda(n)}{\sqrt n}W(n/N).
\]

\begin{lemma}[Character moments on the counting line]
\label{proof:v9:wide-character-moments}
The integrals converge absolutely and satisfy
\begin{equation}\label{proof:app:num:moments}
 \begin{gathered}
 I_\chi=-\sum_{1\le n<N}a_n\chi(n),\\
 \mathcal J=\frac1{5760}\sum_\chi\chi(-1)|I_\chi|^2,\quad
 \mathcal D=\frac1{5760}\sum_\chi|I_\chi|^2.
 \end{gathered}
\end{equation}
All characters, including the principal character, occur; no GRH or
spectral hypothesis is required.
\end{lemma}

\begin{proof}
The line $\Re(1/2+s)=3/2$ gives absolute convergence of $L'/L$.
The $C^6$ Mellin decay permits termwise inversion, yielding the finite sum.
Orthogonality selects $n+m=N$ in $\mathcal J$ and $n=m$ in $\mathcal D$,
since $0<n,m<N$ on the support. For $n+m=N$,
$a_na_m=N^{-1}\Lambda(n)\Lambda(m)f(n/N)f(m/N)$, proving the normalization.
This proves the lemma.
\end{proof}

We integrate with mesh $h=0.1$ and height $T=1600$, using full weights
at the two outer endpoints. A Fourier transform of the weight in the
logarithmic variable, sampled on $2^{20}$ points over a period $2\pi/h$,
supplies the Mellin values. The Hurwitz formula for $L$ and $L'$ uses
$24$ Euler--Maclaurin corrections and a height-dependent cutoff chosen
from the remainder bound. Phase recurrences are reset every $32$ heights.
Character transforms on cyclic groups of orders $2,4,6,10,12$ give all
$5760$ values; conjugation supplies negative heights. The integrated
moments are
\[
 \widehat{\mathcal J}\approx0.05799599193959489,\qquad
 \widehat{\mathcal D}\approx0.29331844084848136.
\]
The finite identity in \eqref{proof:app:num:moments} will be used to
check their error, without presupposing a representation count.

\subsection{The main integral, residual, and integer enclosure}
\label{proof:app:num:interval}

The classical Hardy--Littlewood prediction and the wide-weight main
integral are, respectively,
\begin{align*}
 \mathfrak S(N)\int_2^{N/2}\frac{\mathrm{d}x}{\log x\log(N-x)}
     &\approx908.1405628685,\\
 \mathcal A:=\mathfrak S(N)\int_A^{N/2}
             \frac{\eta(x)^2\,\mathrm{d}x}{\log x\log(N-x)}
     &\approx851.831147718741.
\end{align*}
These use $65$-digit quadrature split at the cutoff transition;
$C_2$ is evaluated by an accelerated zeta product, removing primes up
to $97$ and retaining logarithmic terms through order $60$.
The wide-weight main term with finite corrections is
$\mathcal A-\mathcal C_{\rm pp}/2+\mathcal E\approx886.021328750925$.
Put $c_f=\int_0^1f(u)f(1-u)\,\mathrm{d}u$ and
$\mathcal B=\mathfrak S(N)c_f-\mathcal J$. Since
$\mathcal A=N\mathfrak S(N)c_f/2$, the identity \eqref{proof:app:num:count}
also reads
\[
 r_G(N)=\mathcal A-\frac N2\mathcal B
                   -\frac{\mathcal C_{\rm pp}}2+\mathcal E.
\]
The integrated moments give the residual contribution
$-N\widehat{\mathcal B}/2\approx18.978671254276$, and hence
\begin{equation}\label{proof:app:num:integral-centre}
 \begin{aligned}
 \widehat r_G
 &\approx851.831147718741+18.978671254276\\
 &\hspace{8mm}-10.167587131164+44.357768163348\\
 &\approx905.0000000052.
 \end{aligned}
\end{equation}
This displays the correction to the main integral. To turn the numerical
value into an exact count, we bound the error in the character moment.

\begin{lemma}[From moment error to a counting interval]
\label{proof:v9:moment-error-interval}
Form $\widehat{\mathcal J}$ and $\widehat{\mathcal D}$ from approximations
$\widehat I_\chi$ by \eqref{proof:app:num:moments}. If
$\varepsilon_I^2\ge5760^{-1}\sum_\chi|I_\chi-\widehat I_\chi|^2$ with
$\varepsilon_I\ge0$, then
\[
 |\mathcal J-\widehat{\mathcal J}|
 \le2\sqrt{\widehat{\mathcal D}}\,\varepsilon_I+\varepsilon_I^2.
\]
With exact finite corrections, set
\[
 \widehat r_G=\frac N2\widehat{\mathcal J}
                  -\frac{\mathcal C_{\rm pp}}2+\mathcal E,
 \qquad
 \varepsilon_r=\frac N2
       (2\sqrt{\widehat{\mathcal D}}\,\varepsilon_I+\varepsilon_I^2).
\]
Then $r_G(N)\in[\widehat r_G-\varepsilon_r,
\widehat r_G+\varepsilon_r]\cap\mathbb Z$.
\end{lemma}

\begin{proof}
Writing $e_\chi=I_\chi-\widehat I_\chi$ and using $|\chi(-1)|=1$,
the difference of the moments is bounded by the average of
$2|\widehat I_\chi||e_\chi|+|e_\chi|^2$.
Cauchy--Schwarz gives the stated bound; \eqref{proof:app:num:count}
and integrality give the enclosure. The term $\mathcal A$ cancels,
so its quadrature error does not enter this interval.
This proves the lemma.
\end{proof}

For the error check, the $3280$ nonzero coefficients $a_n$ are enclosed
using rational bounds for logarithms and square roots. Discrete logarithms
to the roots $(2,2,3,2,2)$ modulo $(3,5,7,11,13)$ place them in the
character-transform array. Rational enclosures of the roots of unity
and each transform stage, compared with the stored integral approximations,
give
\[
 \left(\frac1{5760}\sum_\chi|I_\chi-\widehat I_\chi|^2\right)^{1/2}
 <\varepsilon_I=1.68935\times10^{-5}.
\]
This finite-transform check bounds the total integral error, including
quadrature, tails, $L$-evaluation and roundoff. It yields
$\varepsilon_r\approx0.274758682716$; outward rounding gives
\begin{equation}\label{proof:app:num:enclosure}
 [904.725241,\ 905.274759]\cap\mathbb Z=\{905\}.
\end{equation}
Thus the integral value $905.0000000052$ determines the exact integer
$905$, compared with the Hardy--Littlewood prediction $908.1405628685$.
As a final check, direct rational summation in \eqref{proof:app:num:count}
and a separate prime enumeration both give $905$ unordered pairs
($1810$ ordered pairs). This enumeration is a check on the count already
certified from the integrals and finite corrections.

\subsection{The Hardy interpretation and the spectral picture}
\label{proof:app:num:spectrum}

We now take $c=1$ in Appendix~\ref{proof:app:param}, as in the main text.
Here $\alpha_{N,1}=1/\log N\approx0.09699366$, so the Hardy line
$\sigma_*=1/9$ lies between $\alpha_{N,1}$ and $\beta=1/8$.
Under GRH the same operator $T_{N,1}$ realizes the wide-weight residual;
changing the weight changes its testing vector.

\begin{proposition}[The Hardy residual for the wide weight]
\label{proof:v9:wide-hardy-residual}
Assume $\mathrm{GRH}(N)$ at $N=30030$. The wide-weight residual
has the representation
\[
 \mathcal B=\langle T_{N,1}v_W,v_W\rangle,
\]
where $v_W$ is the Riesz vector of the bounded functional on
$\mathcal H_{N,1}$ defined by
\[
 \ell_W(h)=\frac1{2\pi i}\int_{(1/9)}
       N^s\widetilde W(s)(1+s)h(s)\,\mathrm{d}s.
\]
\end{proposition}

\begin{proof}
The real weight $f$ is $C^6$ with compact support in $(0,1)$.
Define $C_f$, $J_f$, and $H_f$ by \eqref{gb:convolutions},
with $f$ in place of $\phi$. At integer arguments, they satisfy
$C_f(k)=0$ for $k\ge1$, $J_f(k)=0$ for $k\ge2$, and $J_f(1)=c_f$;
integration also gives $\widetilde H_f(1)=0$.
These identities and $C^6$ Mellin decay justify model inversion and
boundedness of $\ell_W$. Applying the odd-character argument of
Section~\ref{proof:sec:v3-mellin-operator} gives the quadratic identity
with the unchanged residual kernel.
This proves the proposition.
\end{proof}

The principal-character integral remains on its absolute-convergence
line; this identification does not invoke Hypothesis R.
For the finite spectral model, use $2562$ nodes on the two boundary
lines at height $40$ with mesh $1/16$. The positive trapezoidal weights
$w_i$ have half weight at the endpoints and no $1/(2\pi)$ factor.
The finite spectral function is
\[
 X_{30030}(q):=\det\!\left(I-\frac{4q(1-q)}{\delta_N^2}C^2\right),
\]
where $C=(C_{ij})$ is the matrix with entries
\[
 C_{ij}=\sqrt{w_i}\,
 \frac{\mathscr E_N(s_i,\overline{s_j})}{(1+s_i)(1+\overline{s_j})}
 \sqrt{w_j}.
\]

\begin{proposition}[Zeros of the finite spectral model]
\label{proof:v9:finite-spectral-zeros}
Let $\lambda_1,\ldots,\lambda_{2562}$ be the real eigenvalues of $C$.
Then
\[
 X_{30030}(q)=\prod_{j=1}^{2562}
     \left(1-\frac{4\lambda_j^2}{\delta_N^2}q(1-q)\right),
 \qquad X_{30030}(1-q)=X_{30030}(q).
\]
If $\|C\|\le\delta_N$, every complex zero has real part $1/2$.
\end{proposition}

\begin{proof}
Unitary diagonalization gives the product and symmetry. For
$\lambda_j\ne0$, the zeros satisfy $q(1-q)=\delta_N^2/(4\lambda_j^2)$,
whose roots have real part $1/2$ when $|\lambda_j|\le\delta_N$.
A zero eigenvalue contributes the constant factor one.
This proves the proposition.
\end{proof}

Boundary $L'/L$ values use the Hurwitz formula and character FFT with
$64$ direct terms and $24$ Euler--Maclaurin corrections. Extracting zeta
factors accelerates the prime corrections in $\mathscr D_N$ and $P_N$.
At prime cutoff $2\times10^5$, the extreme eigenvalues are approximately
$-0.353421$ and $0.357462$, with
$\|C\|/\delta_N\approx0.826864<1$ and
$\delta_N=0.432310239587548\ldots$. Doubling the cutoff from $10^5$
changes either endpoint by less than $3\times10^{-8}$.
The computed finite-model zeros therefore lie on the critical line.

\begin{center}
 \begin{minipage}{0.93\textwidth}
 \centering
 \includegraphics[width=0.7742\linewidth]{X30030-heatmap.pdf}
 \par\smallskip
 {\small\raggedright
 \textbf{Figure B.1.} The finite spectral product
 $\log_{10}|X_{30030}(q)|$; the circles mark its displayed zeros.\par}
 \medskip
 \includegraphics[width=0.7742\linewidth]{X30030-zeros.pdf}
 \par\smallskip
 {\small\raggedright
 \textbf{Figure B.2.} The displayed zeros of $X_{30030}(q)$.
 Both figures use the full computed eigenvalue list.\par}
 \end{minipage}
\end{center}
Cutoff stability and pointwise high-precision checks do not bound the
full operator's truncation or discretization error. The figures illustrate
the finite spectral model; they do not certify Hypothesis R for $T_{30030,1}$.

\begin{thebibliography}{99}

\bibitem{BGS2010}
J.~Bourgain, A.~Gamburd, and P.~Sarnak,
\emph{Affine linear sieve, expanders, and sum-product},
Invent. Math. \textbf{179} (2010), no.~3, 559--644.

\bibitem{Chen6673}
J.~R.~Chen,
\emph{On the representation of a large even integer as the sum
of a prime and the product of at most two primes},
Kexue Tongbao \textbf{17} (1966), 385--386 (in Chinese);
Sci. Sinica \textbf{16} (1973), 157--176.

\bibitem{GoldbachEuler1742}
P.-H.~Fuss (ed.),
\emph{Correspondance math\'ematique et physique de quelques
c\'el\`ebres g\'eom\`etres du XVIIIe si\`ecle},
vol.~1, St.~P\'etersbourg, 1843, pp.~125--129
(Goldbach--Euler letters of 7 and 30 June 1742).

\bibitem{Hardy1922}
G.~H.~Hardy,
\emph{Goldbach's theorem},
Mat. Tidsskr. B (1922), 1--16.

\bibitem{HL23}
G.~H.~Hardy and J.~E.~Littlewood,
\emph{Some problems of `Partitio numerorum'; III:
On the expression of a number as a sum of primes},
Acta Math. \textbf{44} (1923), 1--70.

\bibitem{HL24V}
G.~H.~Hardy and J.~E.~Littlewood,
\emph{Some problems of `Partitio numerorum'; V:
A further contribution to the study of Goldbach's problem},
Proc. Lond. Math. Soc. (2) \textbf{22} (1924), 46--56.

\bibitem{IK2004}
H.~Iwaniec and E.~Kowalski,
\emph{Analytic Number Theory},
American Mathematical Society Colloquium Publications, vol.~53,
American Mathematical Society, Providence, RI, 2004.

\bibitem{LeblSCV}
J.~Lebl,
\emph{Tasty Bits of Several Complex Variables},
\href{https://math.libretexts.org/Bookshelves/Analysis/Tasty_Bits_of_Several_Complex_Variables_(Lebl)/01:_Holomorphic_Functions_in_Several_Variables/1.01:_Onto_Several_Variables}
{\S1.1, ``Onto Several Variables''}, online edition.

\bibitem{LiLiu19}
J.~Li and J.~Liu,
\emph{Theorem $(1+1.9)$ on the Goldbach Conjecture},
\href{https://arxiv.org/abs/2606.05224}{arXiv:2606.05224} (2026).

\bibitem{Linnik1963}
Yu.~V.~Linnik,
\emph{The Dispersion Method in Binary Additive Problems},
translated by S.~Schuur,
Translations of Mathematical Monographs, vol.~4,
American Mathematical Society, Providence, RI, 1963.

\bibitem{PanPan}
C.~D.~Pan and C.~B.~Pan,
\emph{Goldbach Conjecture},
Science Press, Beijing, 1981 (in Chinese); 1992 (English edition).

\bibitem{Ren48}
A.~R\'enyi,
\emph{\"Uber die Darstellung der geraden Zahlen als Summe einer
Primzahl und einer Fast-Primzahl},
Izv. Akad. Nauk SSSR Ser. Mat. \textbf{12} (1948), no.~1, 57--78.

\bibitem{Riemann1859}
B.~Riemann,
\emph{Ueber die Anzahl der Primzahlen unter einer gegebenen Gr\"osse},
Monatsberichte der Berliner Akademie (1859), 671--680.
English translation by D.~R.~Wilkins (1998),
\href{https://www.maths.tcd.ie/pub/HistMath/People/Riemann/Zeta/EZeta.pdf}
{\emph{On the number of prime numbers less than a given quantity}}.

\bibitem{Vin37}
I.~M.~Vinogradov,
\emph{Representation of an odd number as a sum of three primes},
Dokl. Akad. Nauk SSSR \textbf{15} (1937), 291--294.

\end{thebibliography}

\end{document}
